% La morale — Philosophie 1re Tech — Cours signé OnBuch+
% Programme officiel MINESEC. Compiler avec : tectonic N07-morale.tex
\newcommand{\DOCMATIERE}{Philosophie}
\newcommand{\DOCNIVEAU}{1re Tech}
\newcommand{\DOCMODULE}{Philosophie – classe de Première technique}
\newcommand{\DOCLECON}{N07}
\newcommand{\DOCTITRE}{La morale}
% OnBuch+ — préambule « cours signés » ENRICHI (compilé avec tectonic / XeLaTeX)
% Système visuel « Pop » d'OnBuch+ : fond crème (jamais blanc), bordures encre
% épaisses, ombres « dures » décalées, titres Archivo Black en capitales.
% Version MATHÉMATIQUES : graphes (pgfplots/tikz), boîtes pédagogiques
% (définition, propriété, méthode, attention, le savais-tu, exercice résolu,
% exercice, TP, bilan), page de garde et signature OnBuch+ ancrée en bas.
%
% Macros attendues dans main.tex AVANT \input{preamble} :
%   \DOCMATIERE  \DOCNIVEAU  \DOCMODULE  \DOCLECON (numéro)  \DOCTITRE
\documentclass[11pt]{article}

\usepackage{fontspec}
\usepackage[french]{babel}
\usepackage[a4paper,margin=2cm,top=3.4cm,bottom=2.8cm,headheight=1.4cm,footskip=1.3cm]{geometry}
\usepackage{amsmath,amssymb,amsfonts}
\usepackage{mathtools} % \xmapsto, \coloneqq... souvent utilisés par les modèles
\usepackage{cancel} % simplifications barrées (bilans de forces)
% \abs{z} (module, valeur absolue) et \norm{u} (norme), utilisés par les modèles
\providecommand{\abs}{}\let\abs\relax\DeclarePairedDelimiter{\abs}{\lvert}{\rvert}
\providecommand{\norm}{}\let\norm\relax\DeclarePairedDelimiter{\norm}{\lVert}{\rVert}
\usepackage[table]{xcolor} % [table] : fournit aussi \rowcolors (alternance de lignes)
\usepackage{pagecolor}
\usepackage{tikz}
\usetikzlibrary{arrows.meta,positioning,calc,shapes.geometric,shapes.misc,shapes.arrows,decorations.pathmorphing,decorations.pathreplacing,decorations.markings,patterns,shadows,mindmap,trees,fit,backgrounds,matrix,intersections,angles,quotes}
\usepackage{pgfplots}
\usepackage[version=4]{mhchem} % équations nucléaires \ce{^{235}_{92}U}
\usepackage[european,siunitx]{circuitikz} % schémas électriques
\pgfplotsset{compat=1.18}
\usepgfplotslibrary{fillbetween,groupplots} % aires entre courbes (calcul intégral)
\usepackage{enumitem}
\usepackage{fancyhdr}
% [most] charge toutes les bibliothèques tcolorbox (listings, minted, raster,
% documentation...) et gonfle inutilement la mémoire TeX sur un document long
% (~300 boîtes) : on ne charge que ce qui est réellement utilisé.
\usepackage[skins,breakable]{tcolorbox}
\usepackage{graphicx}
\usepackage{colortbl}
\usepackage{booktabs}
\usepackage{tabularx}
\usepackage{multirow}
\usepackage{diagbox} % case d'en-tête diagonale des tableaux à double entrée
% Colonnes extensibles centrée / gauche / droite (C, L, R), souvent utilisées par les modèles
\newcolumntype{C}{>{\centering\arraybackslash}X}
\newcolumntype{L}{>{\raggedright\arraybackslash}X}
\newcolumntype{R}{>{\raggedleft\arraybackslash}X}
\usepackage{array}
\usepackage{multicol}
\usepackage{siunitx}
% parse-numbers=false : accepte aussi \SI{2,5 \times 10^{-3}}{...}
\sisetup{locale=FR,output-decimal-marker={,},group-minimum-digits=4,parse-numbers=false}
\DeclareSIUnit\ohm{\ensuremath{\symup{\Omega}}} % U+2126 absent de la police
\DeclareSIUnit\division{div} % réglages d'oscilloscope (ms/div, V/div)
\DeclareSIUnit\tr{tr} % tours (vitesse de rotation, ex. \SI{1500}{\tr\per\minute})
\DeclareSIUnit\molar{M} % concentration molaire (ex. \SI{3}{\molar}, \SI{1}{\micro\molar})
\DeclareSIUnit\an{an} % année (le modèle confond parfois avec \year, un compteur TeX) — ex. \SI{5}{\kilo\meter\per\an}
\DeclareSIUnit\FCFA{FCFA} % franc CFA (ex. \SI{500}{\FCFA\per\kilo\gram})
\DeclareSIUnit\degreeBrix{\textdegree Brix} % teneur en sucre d'un sirop/jus (réfractométrie)
\DeclareSIUnit\month{mois} % le modèle utilise parfois \month, un compteur TeX, comme unité de durée
\DeclareSIUnit\microliter{\micro\liter} % le modèle écrit parfois \microliter en un seul token
\DeclareSIUnit\million{millions} % dénombrement (ex. \SI{15}{\million\per\mL} spermatozoïdes)
\usepackage{qrcode}
% \degree : commande fréquemment utilisée par les modèles pour un angle ou une
% température en dehors de \SI{}{} (non fournie par les paquets chargés).
\providecommand{\degree}{^{\circ}}
% \qed (fin de démonstration), utilisé par les modèles sans amsthm
\providecommand{\qed}{\hfill$\square$}
% \barre : double barre des valeurs interdites dans un tableau de variations (tkz-tab)
\providecommand{\barre}{\ensuremath{\|}}
% environnement proof (amsthm non chargé) utilisé par les modèles
\ifdefined\proof\else\newenvironment{proof}[1][Démonstration]{\par\noindent\textit{#1.}\ }{\qed\par}\fi
\usepackage{lastpage}
\usepackage{hyperref}
\hypersetup{hidelinks}

% ============================================================
% Palette « Pop » OnBuch+ — mêmes hex que la classe P (plus_ui.dart)
% ============================================================
\definecolor{poporange}{HTML}{F59321}
\definecolor{popdark}{HTML}{E07A0C}
\definecolor{popcream}{HTML}{FFF6E8}
\definecolor{popink}{HTML}{1C1714}
\definecolor{poppurple}{HTML}{7A5AE0}
\definecolor{popgreen}{HTML}{2E9E6A}
\definecolor{poppink}{HTML}{E0457B}
\definecolor{popblue}{HTML}{2F7DE1}
\definecolor{popgold}{HTML}{C98A12}
\definecolor{popmuted}{HTML}{8A7F76}
\definecolor{popline}{HTML}{EADFD2}
\definecolor{popgray}{HTML}{8A7F76}
\colorlet{popgrey}{popgray}
% Alias tolérants (le modèle nomme parfois les couleurs différemment)
\colorlet{popwhite}{white}
\colorlet{popblack}{popink}
\colorlet{popred}{poppink}
\colorlet{popyellow}{popgold}
% Teintes claires (fonds de tableaux/encadrés) — le modèle nomme parfois
% les couleurs avec un suffixe L (ex. popblueL) sans les avoir définies.
\colorlet{poporangeL}{poporange!20}
\colorlet{popdarkL}{popdark!20}
\colorlet{poppurpleL}{poppurple!20}
\colorlet{popgreenL}{popgreen!20}
\colorlet{poppinkL}{poppink!20}
\colorlet{popblueL}{popblue!20}
\colorlet{popgoldL}{popgold!20}
\colorlet{popredL}{popred!20}
\colorlet{popgrayL}{popgray!20}
% Teintes claires pour fonds de tableaux / zones
\colorlet{popblueL}{popblue!10!white}
\colorlet{poporangeL}{poporange!14!white}
\colorlet{popgreenL}{popgreen!12!white}
\colorlet{poppurpleL}{poppurple!12!white}
\colorlet{poppinkL}{poppink!12!white}
\colorlet{popgoldL}{popgold!14!white}

\pagecolor{popcream}

% ============================================================
% Typographie
% ============================================================
\defaultfontfeatures{Ligatures=TeX}
\setmainfont{PlusJakartaSans-Medium.ttf}[Path=../../../fonts/,BoldFont=PlusJakartaSans-Bold.ttf]
% Maths en sans-serif (Fira Math) avec chiffres et lettres droites de Plus
% Jakarta, pour que les formules se fondent dans le texte.
\usepackage[math-style=ISO,bold-style=ISO]{unicode-math}
\setmathfont{FiraMath-Regular.otf}[Path=../../../fonts/]
\setmathfont{PlusJakartaSans-Medium.ttf}[Path=../../../fonts/,range={up/num,up/{Latin,latin},\mathup}]
\usepackage{icomma} % virgule décimale sans espace en mode math
% Caractères mathématiques Unicode absents des polices de texte (Plus Jakarta,
% Archivo Black) : rendus par leur équivalent mathématique, en texte comme en
% formule — sinon ils s'affichent en carré vide (ex. « Arithmétique dans ℤ »).
\usepackage{newunicodechar}
\newunicodechar{ℝ}{\ensuremath{\mathbb{R}}}
\newunicodechar{ℤ}{\ensuremath{\mathbb{Z}}}
\newunicodechar{ℂ}{\ensuremath{\mathbb{C}}}
\newunicodechar{ℕ}{\ensuremath{\mathbb{N}}}
\newunicodechar{ℚ}{\ensuremath{\mathbb{Q}}}
\newunicodechar{𝔻}{\ensuremath{\mathbb{D}}}
\newunicodechar{α}{\ensuremath{\alpha}}
\newunicodechar{β}{\ensuremath{\beta}}
\newunicodechar{γ}{\ensuremath{\gamma}}
\newunicodechar{δ}{\ensuremath{\delta}}
\newunicodechar{ε}{\ensuremath{\varepsilon}}
\newunicodechar{θ}{\ensuremath{\theta}}
\newunicodechar{λ}{\ensuremath{\lambda}}
\newunicodechar{μ}{\ensuremath{\mu}}
\newunicodechar{σ}{\ensuremath{\sigma}}
\newunicodechar{τ}{\ensuremath{\tau}}
\newunicodechar{φ}{\ensuremath{\varphi}}
\newunicodechar{ψ}{\ensuremath{\psi}}
\newunicodechar{ω}{\ensuremath{\omega}}
\newunicodechar{Σ}{\ensuremath{\Sigma}}
% majuscules produites par \MakeUppercase dans les titres (α-aminés -> Α)
\newunicodechar{Α}{\ensuremath{\alpha}}
\newunicodechar{Β}{\ensuremath{\beta}}
\newunicodechar{Γ}{\ensuremath{\Gamma}}
\newunicodechar{Δ}{\ensuremath{\Delta}}
\newunicodechar{Ε}{\ensuremath{\varepsilon}}
\newunicodechar{Θ}{\ensuremath{\Theta}}
\newunicodechar{Λ}{\ensuremath{\Lambda}}
\newunicodechar{Μ}{\ensuremath{\mu}}
\newunicodechar{Τ}{\ensuremath{\tau}}
\newunicodechar{Φ}{\ensuremath{\Phi}}
\newunicodechar{Ψ}{\ensuremath{\Psi}}
\newunicodechar{Ω}{\ensuremath{\Omega}}
\newunicodechar{↦}{\ensuremath{\mapsto}}
\newunicodechar{∈}{\ensuremath{\in}}
\newunicodechar{∉}{\ensuremath{\notin}}
\newunicodechar{∧}{\ensuremath{\wedge}}
\newunicodechar{⇔}{\ensuremath{\Leftrightarrow}}
\newunicodechar{⟺}{\ensuremath{\Longleftrightarrow}}
\newunicodechar{⇒}{\ensuremath{\Rightarrow}}
\newunicodechar{⟹}{\ensuremath{\Longrightarrow}}
\newunicodechar{∘}{\ensuremath{\circ}}
\newunicodechar{∪}{\ensuremath{\cup}}
\newunicodechar{∩}{\ensuremath{\cap}}
\newunicodechar{≡}{\ensuremath{\equiv}}
\newunicodechar{⊂}{\ensuremath{\subset}}
\newunicodechar{⊥}{\ensuremath{\perp}}
\newunicodechar{∃}{\ensuremath{\exists}}
\newunicodechar{∀}{\ensuremath{\forall}}
\newunicodechar{∛}{\ensuremath{\sqrt[3]{\,}}}
\newunicodechar{∜}{\ensuremath{\sqrt[4]{\,}}}
\newunicodechar{⃗}{\ensuremath{{}^{\to}}}
\newunicodechar{ⁿ}{\ifmmode^{n}\else\textsuperscript{\ensuremath{n}}\fi}
\newunicodechar{ˣ}{\ifmmode^{x}\else\textsuperscript{\ensuremath{x}}\fi}
\newunicodechar{ᵇ}{\ifmmode^{b}\else\textsuperscript{\ensuremath{b}}\fi}
\newunicodechar{ʳ}{\ifmmode^{r}\else\textsuperscript{\ensuremath{r}}\fi}
\newunicodechar{ᵘ}{\ifmmode^{u}\else\textsuperscript{\ensuremath{u}}\fi}
\newunicodechar{ʸ}{\ifmmode^{y}\else\textsuperscript{\ensuremath{y}}\fi}
\newunicodechar{ᵃ}{\ifmmode^{a}\else\textsuperscript{\ensuremath{a}}\fi}
\newunicodechar{ᵉ}{\ifmmode^{e}\else\textsuperscript{\ensuremath{e}}\fi}
\newunicodechar{ᵏ}{\ifmmode^{k}\else\textsuperscript{\ensuremath{k}}\fi}
\newunicodechar{ᵖ}{\ifmmode^{p}\else\textsuperscript{\ensuremath{p}}\fi}
\newunicodechar{ᴺ}{\ifmmode^{N}\else\textsuperscript{\ensuremath{N}}\fi}
\newunicodechar{ᶜ}{\ifmmode^{c}\else\textsuperscript{\ensuremath{c}}\fi}
\newunicodechar{ʰ}{\ifmmode^{h}\else\textsuperscript{\ensuremath{h}}\fi}
\newunicodechar{ᵠ}{\ifmmode^{q}\else\textsuperscript{\ensuremath{q}}\fi}
\newunicodechar{ₙ}{\ifmmode_{n}\else\textsubscript{\ensuremath{n}}\fi}
\newunicodechar{ₐ}{\ifmmode_{a}\else\textsubscript{\ensuremath{a}}\fi}
\newunicodechar{ᵢ}{\ifmmode_{i}\else\textsubscript{\ensuremath{i}}\fi}
\newunicodechar{ᵣ}{\ifmmode_{r}\else\textsubscript{\ensuremath{r}}\fi}
\newunicodechar{ₖ}{\ifmmode_{k}\else\textsubscript{\ensuremath{k}}\fi}
\newunicodechar{ᵦ}{\ifmmode_{\beta}\else\textsubscript{\ensuremath{\beta}}\fi}
\newunicodechar{ₓ}{\ifmmode_{x}\else\textsubscript{\ensuremath{x}}\fi}
\newfontfamily\popdisplay{ArchivoBlack-Regular.ttf}[Path=../../../fonts/]
\newfontfamily\poplogo{SpaceGrotesk-Bold.ttf}[Path=../../../fonts/]
\newfontfamily\popsemi{PlusJakartaSans-SemiBold.ttf}[Path=../../../fonts/]
\newfontfamily\popxbold{PlusJakartaSans-ExtraBold.ttf}[Path=../../../fonts/]
\newfontfamily\popxboldls{PlusJakartaSans-ExtraBold.ttf}[Path=../../../fonts/,LetterSpace=3.0]

\setlist[enumerate]{leftmargin=1.7em,itemsep=5pt,topsep=4pt}
\setlist[itemize]{leftmargin=1.7em,itemsep=4pt,topsep=4pt,label={\color{poporange}\textbullet}}
\setlength{\parindent}{0pt}
\setlength{\parskip}{5pt}
\renewcommand{\arraystretch}{1.25}

% pgfplots : style Pop par défaut
\pgfplotsset{
  every axis/.append style={axis line style={popink,line width=1pt},tick style={popink},
    grid style={popline,line width=0.5pt},label style={font=\small},tick label style={font=\footnotesize}},
  popcurve/.style={poporange,line width=1.8pt,no markers,smooth},
}
% Même style disponible dans un \draw TikZ simple (hors pgfplots)
\tikzset{popcurve/.style={poporange,line width=1.8pt}}
\tikzset{popgrid/.style={popline,very thin}}
\colorlet{popcurve}{poporange} % le nom du style est parfois utilisé comme couleur

% ============================================================
% Pill / squiggle
% ============================================================
\newcommand{\poppill}[3]{%
  \tikz[baseline=-0.55ex]\node[draw=popink,line width=1.2pt,rounded corners=2.6pt,%
    fill=#2,inner xsep=6.5pt,inner ysep=3pt]{%
    \popxboldls\fontsize{7}{7}\selectfont\color{#3}\MakeUppercase{#1}};%
}
\newcommand{\popsquiggle}[1]{%
  \tikz[baseline=-0.4ex]{
    \draw[#1,line width=2.6pt,line cap=round]
      plot [smooth] coordinates {(0,0.09) (0.34,0.01) (0.68,0.21) (1.02,0.01) (1.36,0.10)};
  }%
}
% Mot-clé surligné (marqueur orange sous le texte)
\newcommand{\cle}[1]{{\popxbold\color{popdark}#1}}

% ============================================================
% En-tête / pied
% ============================================================
\pagestyle{fancy}
\fancyhf{}
\renewcommand{\headrulewidth}{0pt}
\renewcommand{\footrulewidth}{0pt}
\lhead{\raisebox{-0.15ex}{\poplogo\fontsize{13}{13}\selectfont\color{popink}OnBuch\color{poporange}+}}
\rhead{\poppill{\DOCMATIERE\ \textbullet\ \DOCNIVEAU\ \textbullet\ Leçon \DOCLECON}{white}{popink}}
\lfoot{\popxbold\fontsize{7.6}{7.6}\selectfont\color{popmuted}\MakeUppercase{Cours signé OnBuch+}\ \textbullet\ {\popsemi\DOCTITRE}}
\rfoot{\tikz[baseline=-0.6ex]\node[rounded corners=8.5pt,fill=popink,minimum height=17pt,minimum width=17pt,inner xsep=5pt,inner sep=0]{\popxbold\fontsize{8}{8}\selectfont\color{popcream}\thepage\,/\,\pageref*{LastPage}};}

% ============================================================
% Titres
% ============================================================
\newcommand{\chaptitre}[2]{%
  \begin{tcolorbox}[enhanced,colback=poporange,colframe=popink,boxrule=2.6pt,arc=11pt,
      drop shadow=popink,left=15pt,right=15pt,top=12pt,bottom=11pt,before skip=3pt,after skip=18pt]
    \poppill{#2}{popink}{popcream}\par\vspace{9pt}
    {\raggedright\hyphenpenalty=10000\exhyphenpenalty=10000\popdisplay\fontsize{22}{24}\selectfont\color{popcream}\MakeUppercase{#1}\par}
  \end{tcolorbox}
}
% Section numérotée : pastille orange avec le numéro + titre Archivo
\newcounter{popsec}
\newcommand{\coursec}[1]{%
  \stepcounter{popsec}\setcounter{popsub}{0}%
  \par\vspace{16pt}\noindent
  \begin{minipage}{\linewidth}
  \tikz[baseline=-0.7ex]\node[draw=popink,line width=1.6pt,fill=poporange,rounded corners=4pt,minimum size=22pt,inner sep=0]{\popdisplay\fontsize{12}{12}\selectfont\color{popcream}\thepopsec};%
  \hspace{8pt}{\popdisplay\fontsize{15}{17}\selectfont\color{popink}\MakeUppercase{#1}}\par
  \vspace{4pt}\noindent{\color{poporange}\rule{36pt}{3.6pt}}
  \end{minipage}\par\nopagebreak\vspace{8pt}
}
\newcounter{popsub}
\newcommand{\courssub}[1]{%
  \stepcounter{popsub}%
  \par\vspace{9pt}\noindent{\popxbold\fontsize{11.5}{13}\selectfont\color{popink}{\color{poporange}\thepopsec.\thepopsub}\hspace{6pt}#1}\par\nopagebreak\vspace{4pt}
}

% ============================================================
% Boîtes pédagogiques — même recette : fond blanc, bordure épaisse colorée,
% ombre dure encre, badge plein. Titre optionnel [#1].
% ============================================================
\tcbset{popbase/.style={breakable,enhanced,colback=white,boxrule=2.3pt,arc=9pt,
  shadow={3pt}{-3pt}{0pt}{popink},left=14pt,right=14pt,top=11pt,bottom=11pt,before skip=11pt,after skip=15pt,
  fontupper=\popsemi\fontsize{10.6}{14.5}\selectfont\color{popink},
  fontlower=\popsemi\fontsize{10.6}{14.5}\selectfont\color{popink}}}
% Coche dessinée (le glyphe ✓ n'existe pas dans les polices du document)
\newcommand{\popcheck}[1]{\tikz[baseline=-0.5ex]\draw[#1,line width=1.6pt,line cap=round,line join=round](-0.13,0.0)--(-0.04,-0.09)--(0.13,0.1);}
\providecommand{\coche}{\popcheck{popgreen}}
\providecommand{\resultat}[1]{\par\smallskip\noindent\fcolorbox{popgreen}{popgreenL}{\parbox{\dimexpr\linewidth-2\fboxsep-2\fboxrule\relax}{\textbf{Résultat.} #1}}\par\smallskip}
\newcommand{\popboxhead}[3]{\poppill{#1}{#2}{white}\ifx\relax#3\relax\else\hspace{6pt}{\popxbold\fontsize{10.6}{12}\selectfont\color{popink}#3}\fi\par\vspace{7pt}}

\newtcolorbox{aretenir}[1][]{popbase,colframe=popgreen,before upper={\popboxhead{À retenir}{popgreen}{#1}}}
\newtcolorbox{exemplebox}[1][]{popbase,colframe=poporange,before upper={\popboxhead{Exemple}{poporange}{#1}}}
\newtcolorbox{definition}[1][]{popbase,colframe=popblue,before upper={\popboxhead{Définition}{popblue}{#1}}}
\newtcolorbox{propriete}[1][]{popbase,colframe=poppurple,before upper={\popboxhead{Propriété}{poppurple}{#1}}}
\newtcolorbox{methode}[1][]{popbase,colframe=popgold,before upper={\popboxhead{Méthode}{popgold}{#1}}}
\newtcolorbox{attention}[1][]{popbase,colframe=poppink,before upper={\popboxhead{Attention}{poppink}{#1}}}
\newtcolorbox{experience}[1][]{popbase,colframe=popblue,colback=popblueL,before upper={\popboxhead{Expérience}{popblue}{#1}}}
% « Le savais-tu ? » : carte violette pleine (contexte, culture, Cameroun)
\newtcolorbox{savaistu}[1][]{popbase,colback=poppurple,colframe=popink,
  fontupper=\popsemi\fontsize{10.4}{14.2}\selectfont\color{white},
  before upper={\poppill{Le savais-tu\,?}{white}{poppurple}\ifx\relax#1\relax\else\hspace{6pt}{\popxbold\color{white}#1}\fi\par\vspace{7pt}}}
% Exercice résolu : énoncé en haut, \tcblower puis la solution sur fond crème
\newcounter{popexo}\newcounter{popexr}
\newtcolorbox{exoresolu}[1][]{popbase,colframe=poporange,colbacklower=popcream,
  segmentation style={popink,line width=1.2pt,dashed},
  before upper={\stepcounter{popexr}\popboxhead{Exercice résolu \thepopexr}{poporange}{#1}},
  before lower={\poppill{Solution}{popink}{popcream}\par\vspace{6pt}}}
% Exercice (non résolu en place) — la correction est dans la partie Corrigés
\newtcolorbox{exercice}[1][]{popbase,colframe=popink,
  before upper={\stepcounter{popexo}\popboxhead{Exercice \thepopexo}{popink}{#1}}}
% Corrigé
\newtcolorbox{corrige}[1][]{popbase,colframe=popgreen,colback=popgreenL,
  before upper={\popboxhead{Corrigé}{popgreen}{#1}}}
% Objectifs / prérequis / bilan
\newtcolorbox{objectifs}{popbase,colframe=popink,colback=poporangeL,
  before upper={\popboxhead{Objectifs de la leçon}{popink}{}}}
\newtcolorbox{prerequis}{popbase,colframe=popink,colback=popblueL,
  before upper={\popboxhead{Prérequis}{popblue}{}}}
\newtcolorbox{bilan}{popbase,colframe=popink,colback=white,boxrule=2.8pt,
  before upper={\popboxhead{Fiche bilan}{poporange}{L'essentiel à connaître pour l'examen}}}
% Figure encadrée avec légende
\newtcolorbox{popfigure}[1][]{enhanced,breakable=false,colback=white,colframe=popline,boxrule=1.4pt,arc=8pt,
  left=10pt,right=10pt,top=10pt,bottom=8pt,before skip=10pt,after skip=12pt,halign=center,
  lower separated=false,halign lower=center,
  fontlower=\popsemi\fontsize{9}{11.5}\selectfont\color{popmuted},
  #1}
\newcommand{\legende}[1]{\par\vspace{6pt}{\popsemi\fontsize{9}{11.5}\selectfont\color{popmuted}#1}}

% ============================================================
% Signature OnBuch+
% ============================================================
\newcommand{\poplogoinline}[1]{{\poplogo\fontsize{#1}{#1}\selectfont\color{popink}OnBuch\color{poporange}+}}
\newcommand{\signatureonbuch}{%
  \begin{tcolorbox}[enhanced,colback=popink,colframe=popink,boxrule=2.2pt,arc=10pt,
      drop shadow=poporange,left=14pt,right=14pt,top=10pt,bottom=10pt,before skip=0pt,after skip=0pt]
    \tikz[baseline=-0.7ex]\node[circle,fill=popgreen,draw=popcream,line width=1.3pt,minimum size=22pt,inner sep=0]{\popcheck{white}};%
    \hspace{9pt}\begin{minipage}[c]{0.62\linewidth}
      {\popxbold\fontsize{12}{13}\selectfont\color{popcream}Signé\ \textbullet\ {\poplogo OnBuch\color{poporange}+}}\par\vspace{2pt}
      {\raggedright\popsemi\fontsize{8}{10.5}\selectfont\color{popline}\MakeUppercase{Équipe pédagogique OnBuch+}\\\MakeUppercase{Conforme au programme officiel MINESEC}\par}
    \end{minipage}\hfill
    {\poplogo\fontsize{22}{22}\selectfont\color{popcream}OnBuch\color{poporange}+}
  \end{tcolorbox}%
}

% ============================================================
% Page de garde
% #1 titre · #2 matière · #3 niveau · #4 module · #5 n° leçon · #6 durée ·
% #7 description / objectifs (texte ou itemize)
% ============================================================
\newcommand{\pagegarde}[7]{%
  \thispagestyle{empty}
  \newgeometry{a4paper,margin=1.7cm,top=1.6cm,bottom=1.4cm}%
  \begin{tikzpicture}[remember picture,overlay]
    \fill[poporange!18] ([xshift=-3.2cm,yshift=-3.6cm]current page.north east) circle (4.6cm);
    \fill[poppurple!14] ([xshift=2.4cm,yshift=7.6cm]current page.south west) circle (3.8cm);
  \end{tikzpicture}%
  {\poplogo\fontsize{30}{30}\selectfont\color{popink}OnBuch\color{poporange}+}\hfill
  \raisebox{0.6ex}{\poppill{Cours signé \textbullet\ Premium}{popink}{popcream}}\par
  \vspace{1pt}\popsquiggle{poppurple}\par
  \vspace{8pt}
  \begin{tcolorbox}[enhanced,colback=poporange,colframe=popink,boxrule=2.8pt,arc=13pt,
      drop shadow=popink,left=18pt,right=18pt,top=12pt,bottom=13pt,after skip=10pt]
    \poppill{#2 \textbullet\ #3}{popink}{popcream}\hspace{5pt}\poppill{#4}{white}{popink}\par\vspace{8pt}
    {\popdisplay\fontsize{14}{14}\selectfont\color{popink}LEÇON #5}\par\vspace{4pt}
    {\raggedright\hyphenpenalty=10000\exhyphenpenalty=10000\popdisplay\fontsize{25}{27}\selectfont\color{popcream}\MakeUppercase{#1}\par}
    \vspace{8pt}
    {\popxbold\fontsize{9.5}{9.5}\selectfont\color{popink}\MakeUppercase{Durée conseillée : #6}}
  \end{tcolorbox}
  \begin{tcolorbox}[enhanced,colback=white,colframe=popink,boxrule=2.2pt,arc=10pt,
      drop shadow=popink,left=16pt,right=16pt,top=10pt,bottom=10pt,after skip=8pt]
    \poppill{Dans cette leçon}{poppurple}{white}\par\vspace{5pt}
    \popsemi\fontsize{9.3}{12.6}\selectfont\color{popink}#7
  \end{tcolorbox}
  \noindent\begin{minipage}[t]{0.487\linewidth}
    \begin{tcolorbox}[enhanced,colback=popgreen,colframe=popink,boxrule=2.2pt,arc=10pt,
        drop shadow=popink,left=11pt,right=11pt,top=10pt,bottom=10pt,height=3.1cm,valign=center]
      \tcbox[colback=white,colframe=popink,boxrule=1.3pt,arc=5pt,boxsep=0pt,nobeforeafter,
        left=3pt,right=3pt,top=3pt,bottom=3pt,valign=center]{\qrcode[height=1.9cm]{https://play.google.com/store/apps/details?id=com.luvvix.onbuch&hl=fr}}%
      \hspace{8pt}\begin{minipage}[c]{0.5\linewidth}
        {\popdisplay\fontsize{11}{12.5}\selectfont\color{white}\MakeUppercase{Télécharge OnBuch}}\par\vspace{4pt}
        {\raggedright\popsemi\fontsize{8.4}{11}\selectfont\color{white}Gratuit \textbullet\ Google Play\\Tuteur IA, annales, concours, résultats\par}
      \end{minipage}
    \end{tcolorbox}
  \end{minipage}\hfill
  \begin{minipage}[t]{0.487\linewidth}
    \begin{tcolorbox}[enhanced,colback=poppurple,colframe=popink,boxrule=2.2pt,arc=10pt,
        drop shadow=popink,left=11pt,right=11pt,top=10pt,bottom=10pt,height=3.1cm,valign=center]
      \tikz[baseline=-0.7ex]\node[circle,fill=white,draw=popink,line width=1.3pt,minimum size=1.6cm,inner sep=0]{\popdisplay\fontsize{24}{24}\selectfont\color{poppurple}+};%
      \hspace{8pt}\begin{minipage}[c]{0.6\linewidth}
        {\popdisplay\fontsize{11}{12.5}\selectfont\color{white}\MakeUppercase{Passe à OnBuch+}}\par\vspace{4pt}
        {\raggedright\popsemi\fontsize{8.4}{11}\selectfont\color{white}Tous les cours signés, Tuteur IA illimité, fiches PDF — onglet OnBuch+ de l'app\par}
      \end{minipage}
    \end{tcolorbox}
  \end{minipage}
  \vspace{8pt}
  \popcontenu
  \vfill
  \signatureonbuch
  \restoregeometry
  \newpage
}

% Rangée « Ce cours contient » (page de garde)
\newcommand{\poptile}[3]{% #1 couleur #2 gros texte #3 légende
  \begin{tcolorbox}[enhanced,colback=#1,colframe=popink,boxrule=2pt,arc=9pt,shadow={2.5pt}{-2.5pt}{0pt}{popink},
      left=6pt,right=6pt,top=9pt,bottom=9pt,halign=center,valign=center,height=1.75cm,width=\linewidth,nobeforeafter]
    {\popdisplay\fontsize{12.5}{13}\selectfont\color{white}\MakeUppercase{#2}}\par\vspace{3pt}
    {\centering\popsemi\fontsize{7.8}{9.5}\selectfont\color{white}#3\par}
  \end{tcolorbox}}
\newcommand{\popcontenu}{%
  \par\noindent{\popxboldls\fontsize{7.5}{7.5}\selectfont\color{popmuted}CE COURS SIGNÉ CONTIENT}\par\vspace{7pt}
  \noindent\begin{minipage}[t]{0.235\linewidth}\poptile{popblue}{Cours}{complet, illustré, pas à pas}\end{minipage}\hfill
  \begin{minipage}[t]{0.235\linewidth}\poptile{poporange}{Méthodes}{et exercices résolus}\end{minipage}\hfill
  \begin{minipage}[t]{0.235\linewidth}\poptile{poppink}{Exercices}{type examen + corrigés}\end{minipage}\hfill
  \begin{minipage}[t]{0.235\linewidth}\poptile{popgreen}{Fiche bilan}{l'essentiel à retenir}\end{minipage}\par}

% Sommaire
\newtcolorbox{sommaire}{enhanced,colback=white,colframe=popink,boxrule=2.2pt,arc=10pt,
  drop shadow=popink,left=18pt,right=18pt,top=13pt,bottom=13pt,before skip=6pt,after skip=20pt,
  before upper={\poppill{Sommaire}{poppurple}{white}\par\vspace{9pt}\popsemi\fontsize{10.6}{14.5}\selectfont\color{popink}}}

% Signature de fin de cours — poussée tout en bas de la dernière page
\newcommand{\findecours}{%
  \par\vfill\nopagebreak
  \begin{center}\popsquiggle{poporange}\end{center}\vspace{4pt}
  \signatureonbuch
}
\AtBeginDocument{\def\llbracket{\mathopen{[\mkern-3.5mu[}}\def\rrbracket{\mathclose{]\mkern-3.5mu]}}} % intervalles d'entiers (glyphe ⟦ absent de la police)
% Glyphes absents de Fira Math : redéfinis à partir de glyphes disponibles
\AtBeginDocument{%
  \def\setminus{\mathbin{\backslash}}\let\smallsetminus\setminus
  \def\vdots{\mathord{\vbox{\baselineskip3.2pt\lineskiplimit0pt\kern4pt\hbox{.}\hbox{.}\hbox{.}}}}%
  \def\ddots{\mathinner{\mkern1mu\raise6.5pt\hbox{.}\mkern2mu\raise3.6pt\hbox{.}\mkern2mu\raise0.7pt\hbox{.}\mkern1mu}}%
  \def\triangle{\mathord{\scalebox{0.9}{$\symup{\Delta}$}}}%
  \def\nabla{\mathord{\raisebox{\depth}{\scalebox{1}[-1]{$\symup{\Delta}$}}}}%
  \def\checkmark{\popcheck{popdark}}%
  \def\star{\mathbin{\ast}}\def\bigstar{\mathord{\ast}}%
  \def\diamond{\mathbin{\scalebox{0.6}{$\Diamond$}}}%
  \def\bigcup{\mathop{\mathchoice{\raisebox{-0.3em}{\scalebox{1.7}{$\cup$}}}{\scalebox{1.25}{$\cup$}}{\cup}{\cup}}\displaylimits}%
  \def\bigcap{\mathop{\mathchoice{\raisebox{-0.3em}{\scalebox{1.7}{$\cap$}}}{\scalebox{1.25}{$\cap$}}{\cap}{\cap}}\displaylimits}%
  \def\lesseqqgtr{\mathrel{\lessgtr}}\def\bigoplus{\mathop{\oplus}\displaylimits}%
}
\providecommand{\ssi}{si et seulement si} % abréviation fréquente
\AtBeginDocument{\providecommand{\wideparen}{\overparen}} % arc de cercle

%% FIGFIX-BEGIN v5 — correctifs globaux des figures (docs/onbuch-generation/tools/figfix)
\usepackage{adjustbox}\usepackage{etoolbox}\tcbuselibrary{raster}
% toute figure TikZ/pgfplots plus large que la ligne est réduite à la largeur disponible
\BeforeBeginEnvironment{tikzpicture}{\begin{adjustbox}{max width=\linewidth}}
\AfterEndEnvironment{tikzpicture}{\end{adjustbox}}
% légendes pgfplots lisibles par-dessus les courbes
\pgfplotsset{every axis/.append style={legend style={fill=white,fill opacity=0.92,text opacity=1,draw=black!15,inner sep=3pt}}}
% étiquettes d'arêtes (syntaxe edge["..."]) avec fond blanc
\tikzset{every edge quotes/.append style={fill=white,inner sep=1.2pt}}
%% FIGFIX-END
\begin{document}

\pagegarde{La morale}{Philosophie}{1re Tech}{Philosophie – classe de Première technique}{N07}{3 séances (fiche de progression harmonisée)}{Cette leçon interroge la notion de morale : sa définition, son fondement et ses grandes doctrines. À travers les morales de l'intérêt (épicurisme, eudémonisme, utilitarisme) et la morale du devoir, l'élève apprend à analyser et à juger les conduites humaines dans des situations concrètes de la vie camerounaise.
\vspace{4pt}
\begin{multicols}{2}\small
\begin{itemize}
\item À la fin de cette leçon, je sais définir la morale et la distinguer des notions voisines (mœurs, éthique, droit)
\item À la fin de cette leçon, je sais identifier et discuter les différents fondements proposés pour la morale (Dieu, la nature, la société, la raison)
\item À la fin de cette leçon, je sais expliquer la doctrine épicurienne et l'illustrer par des exemples de la vie courante
\item À la fin de cette leçon, je sais présenter l'eudémonisme d'Aristote et la notion de juste milieu
\item À la fin de cette leçon, je sais exposer l'utilitarisme de Bentham et de Mill et le critère du plus grand bonheur du plus grand nombre
\item À la fin de cette leçon, je sais décrire la morale kantienne du devoir et formuler l'impératif catégorique
\end{itemize}\end{multicols}}

\begin{sommaire}
\begin{multicols}{2}
\begin{enumerate}
\item Définition et fondement de la morale
\item Les morales de l'intérêt : présentation générale
\item L'épicurisme : la morale du plaisir
\item L'eudémonisme : la morale du bonheur
\item L'utilitarisme : la morale de l'utilité
\item La morale du devoir : Kant
\item Synthèse et application aux situations concrètes
\item Activité d'intégration
\item Méthodes et exercices résolus
\item Exercices
\item Corrigés des exercices
\item Fiche bilan
\end{enumerate}
\end{multicols}
\end{sommaire}

\begin{objectifs}
\begin{itemize}
\item À la fin de cette leçon, je sais définir la morale et la distinguer des notions voisines (mœurs, éthique, droit)
\item À la fin de cette leçon, je sais identifier et discuter les différents fondements proposés pour la morale (Dieu, la nature, la société, la raison)
\item À la fin de cette leçon, je sais expliquer la doctrine épicurienne et l'illustrer par des exemples de la vie courante
\item À la fin de cette leçon, je sais présenter l'eudémonisme d'Aristote et la notion de juste milieu
\item À la fin de cette leçon, je sais exposer l'utilitarisme de Bentham et de Mill et le critère du plus grand bonheur du plus grand nombre
\item À la fin de cette leçon, je sais décrire la morale kantienne du devoir et formuler l'impératif catégorique
\item À la fin de cette leçon, je sais comparer les morales de l'intérêt et la morale du devoir dans un tableau synthétique
\item À la fin de cette leçon, je sais appliquer ces doctrines à des situations concrètes camerounaises et justifier mon jugement moral
\item À la fin de cette leçon, je sais rédiger une argumentation structurée (thèse, arguments, exemples, conclusion) sur un sujet de morale
\end{itemize}
\end{objectifs}

\begin{prerequis}
\begin{itemize}
\item Notion de valeurs et de normes sociales (vie en société, collège)
\item Distinction entre le permis et le défendu dans la famille et la communauté
\item Notions de bonheur, plaisir et devoir telles qu'employées dans le langage courant
\item Capacité à lire et commenter un texte philosophique court (travail de Seconde)
\end{itemize}
\end{prerequis}

\newpage

\coursec{Définition et fondement de la morale}

\courssub{Situation d'accroche et problématique}

À Yaoundé, au cœur du marché Mokolo, un jeune commerçant vient de vendre un sac de riz. Le client, pressé et distrait, lui a tendu un billet de trop : 2 000 FCFA de plus que le prix convenu. Le client s'éloigne déjà dans la foule. Le jeune commerçant hésite. Doit-il courir après lui pour rendre la monnaie exacte, ou garder discrètement les 2 000 FCFA ?

Autour de lui, les avis fusent. Son voisin de boutique lui glisse : « Fais ce qui t'arrange, personne ne t'a vu. » Sa grand-mère, à qui il raconte la scène le soir, lui répond : « Fais ton devoir, même si personne ne regarde. » Un troisième ami lui conseille : « Demande-toi ce qui rendra tout le monde plus heureux. »

Ces trois conseils ne sont pas anodins : ils incarnent trois grandes conceptions de la morale que l'humanité a élaborées au fil des siècles — la morale de l'intérêt, la morale du devoir, la morale de l'utilité. Mais avant de les étudier, une question plus fondamentale s'impose : \emph{qu'est-ce que la morale ?} Et surtout : \emph{sur quoi repose-t-elle ?} Pourquoi devrais-je obéir à la morale, alors que rien ne m'y oblige physiquement ? Telle est la \cle{problématique} de cette première section : définir la morale, la distinguer de ce qui lui ressemble (les mœurs, le droit, la religion), puis interroger son \cle{fondement}, c'est-à-dire la source de l'obligation morale.

\courssub{Étymologie : morale et éthique, une même réalité}

Commençons par les mots eux-mêmes, car ils portent déjà un enseignement. Le mot \emph{morale} vient du latin \emph{mos, moris} (au pluriel \emph{mores}), qui signifie « coutume », « manière d'agir », « habitude de vie ». Le mot \emph{éthique}, lui, vient du grec \emph{ethos}, qui signifie exactement la même chose : la coutume, le caractère, la manière habituelle de se conduire.

\begin{savaistu}[Morale et éthique : deux mots, un seul sens]
Le mot « éthique » vient du grec \emph{ethos} et le mot « morale » du latin \emph{mores} : les deux mots signifient exactement la même chose, à savoir \emph{les coutumes}. C'est pourquoi, en philosophie, « morale » et « éthique » désignent la même réalité et peuvent être employés l'un pour l'autre. On parle indifféremment de « réflexion morale » ou de « réflexion éthique » (par exemple : l'éthique médicale, la morale professionnelle).
\end{savaistu}

Cette étymologie révèle une intuition profonde : à l'origine, la morale est affaire de \emph{conduite}, de manière d'agir au quotidien, et non d'abord de théorie abstraite. La morale naît de la question toute simple que chacun se pose, comme notre commerçant de Mokolo : \emph{que dois-je faire ?}

\courssub{Définition de la morale}

De cette étymologie, nous pouvons tirer une définition rigoureuse.

\begin{definition}[La morale]
La \cle{morale} est l'ensemble des \emph{règles de conduite} et des \emph{valeurs} (le \emph{bien} et le \emph{mal}) qui orientent les actions humaines. Elle permet de juger ce qui est bien (ce qui doit être fait) et ce qui est mal (ce qui doit être évité), et elle s'adresse à la conscience de chacun sous la forme d'un devoir : « tu dois », « tu ne dois pas ».
\end{definition}

Analysons cette définition. La morale comporte trois éléments :

\begin{itemize}
\item des \textbf{règles de conduite} : des prescriptions (« sois honnête », « ne mens pas », « rends ce qui ne t'appartient pas ») qui guident l'action ;
\item des \textbf{valeurs} : le \emph{bien} et le \emph{mal}, c'est-à-dire les critères grâce auxquels on juge les actions, les intentions et les personnes ;
\item une \textbf{exigence} : la morale ne se contente pas de décrire ce que les hommes font ; elle dit ce qu'ils \emph{devraient} faire. C'est la dimension du \emph{devoir}, de l'obligation.
\end{itemize}

Ainsi, quand le commerçant se demande s'il doit rendre les 2 000 FCFA, il fait une expérience morale : il éprouve un conflit entre ce qu'il désire (garder l'argent) et ce qu'il sent être son devoir (être honnête). Cette voix intérieure qui juge nos actes, on l'appelle la \cle{conscience morale}.

\courssub{Distinctions essentielles : morale, mœurs, droit, religion}

\textbf{Morale et mœurs.} Il faut d'abord distinguer la morale des \emph{mœurs}.

\begin{definition}[Les mœurs]
Les \cle{mœurs} sont l'ensemble des usages et coutumes \emph{effectivement pratiqués} dans une société donnée, à une époque donnée : manières de se saluer, de s'habiller, de se marier, de manger, d'enterrer les morts, etc.
\end{definition}

La différence est capitale : les mœurs \emph{décrivent} ce qui se fait ; la morale \emph{prescrit} ce qui doit être fait. Les mœurs varient d'un peuple à l'autre et d'une époque à l'autre (relativité des coutumes), tandis que la morale prétend à une certaine universalité : dire « il ne faut pas mentir » vaut pour tous les hommes, pas seulement pour un groupe. On peut d'ailleurs juger les mœurs d'une société au nom de la morale : une coutume répandue (par exemple la lépreuse exclusion dont sont victimes les personnes albinos dans certaines régions) peut être condamnée moralement, même si elle est effectivement pratiquée.

\textbf{Morale et droit.} La seconde distinction, classique au Probatoire, oppose la morale et le \emph{droit}.

\begin{attention}[Ne pas confondre morale et droit]
Erreur fréquente : assimiler la morale au droit. Le \emph{droit} est l'ensemble des lois édictées par l'État ; sa contrainte est \emph{extérieure} et sa sanction est publique (amende, prison, tribunal). La \emph{morale}, elle, s'adresse à la conscience intérieure : sa sanction est le remords, la honte, le sentiment de culpabilité. Une action peut être légale sans être morale, et inversement.
\end{attention}

\begin{exemplebox}[Immoral et/ou illégal : des cas camerounais]
Au Cameroun, jeter des ordures sur la voie publique à Douala peut être \emph{à la fois} immoral (manque de civisme, mépris de la santé de tous) \emph{et} illégal (contravention, amende prévue par la réglementation municipale). En revanche, mentir à un ami est \emph{immoral} (c'est trahir sa confiance) mais \emph{pas illégal} : aucun tribunal ne vous condamnera pour cela. Inversement, certaines actions légales peuvent choquer la conscience morale : exploiter la misère d'autrui dans le respect strict de la loi reste moralement condamnable.
\end{exemplebox}

\begin{center}
\begin{tabularx}{\linewidth}{|l|X|X|}\hline
\rowcolor{popblueL} \textbf{Critère} & \textbf{Morale} & \textbf{Droit} \\ \hline
\textbf{Source} & La conscience, la raison, les valeurs transmises & L'État (lois votées, décrets, règlements) \\ \hline
\textbf{Sanction} & Intérieure : remords, culpabilité, estime de soi & Extérieure : amende, prison, dommages et intérêts \\ \hline
\textbf{Domaine} & Toute la conduite, y compris les intentions secrètes et la vie privée & Les actes extérieurs qui intéressent l'ordre social \\ \hline
\textbf{Portée} & Universelle (vaut pour tous les hommes) & Limitée à un territoire et à une époque \\ \hline
\end{tabularx}
\end{center}

\textbf{Morale et religion.} Enfin, la morale n'est pas la religion. La religion lie l'homme à Dieu par la foi, le culte et les rites ; la morale concerne les rapports entre les hommes et s'adresse à la conscience. Certes, les religions proposent des règles morales (les Dix Commandements, par exemple), mais on peut être moral sans être croyant, et l'histoire montre malheureusement qu'on peut se dire croyant sans être moral. La question du rapport entre morale et religion se posera de nouveau à propos du fondement.

\courssub{Le problème du fondement : pourquoi dois-je obéir à la morale ?}

Nous voici au cœur de la section. La morale me dit : « tu dois être honnête ». Mais une question surgit immédiatement : \emph{pourquoi ?} Qui me l'ordonne ? Sur quoi repose cette obligation ? C'est le problème du \cle{fondement} de la morale : chercher le fondement, c'est chercher la \emph{source}, le \emph{socle}, la \emph{justification dernière} de l'obligation morale.

Ce problème n'est pas théorique : il se pose concrètement. Le voisin du commerçant de Mokolo dit : « personne ne t'a vu ». Autrement dit : si nul ne peut me punir, pourquoi être honnête ? Si la morale ne repose sur rien de solide, alors l'intérêt personnel devient la seule loi. Quatre grandes réponses ont été proposées à la question du fondement.

\textbf{1. Le fondement religieux : la morale vient de Dieu.} Pour les croyants, la morale est fondée sur la volonté divine : le bien est ce que Dieu commande, le mal ce qu'il interdit. Les Dix Commandements donnés à Moïse (« Tu ne tueras point », « Tu ne voleras point », « Honore ton père et ta mère ») en sont l'exemple le plus célèbre. L'obligation morale est alors absolue, car elle émane d'un être parfait et tout-puissant qui voit tout, même ce qui est caché : « personne ne t'a vu » ne vaut plus, puisque Dieu voit.

Mais cette réponse soulève une difficulté décisive : \emph{que dire à celui qui ne croit pas ?} Si la morale repose uniquement sur Dieu, alors l'incroyant n'a aucune raison d'être moral — ce qui est faux, car il existe d'excellents hommes sans religion. De plus, les religions diffèrent entre elles : la morale deviendrait-elle multiple ?

\textbf{2. Le fondement naturel : la morale est inscrite dans la nature humaine.} Une seconde réponse affirme que la morale est écrite dans la nature même de l'homme. C'est la thèse de la \emph{loi naturelle} : l'homme, par sa seule nature raisonnable, connaît spontanément le bien et le mal. Les \emph{stoïciens} (philosophes grecs et romains comme Épictète et Marc Aurèle) soutiennent que vivre moralement, c'est vivre « conformément à la nature » : la nature humaine, étant raisonnable et sociale, nous porte naturellement vers la justice et l'entraide. De même, ne sentons-nous pas spontanément, devant un enfant maltraité, que c'est mal, sans qu'on nous l'ait appris ? Cette répulsion immédiate serait la trace de la loi morale inscrite en nous.

\textbf{3. Le fondement social : la morale est un produit de la société.} Le sociologue Émile \emph{Durkheim} (1858-1917) défend une troisième thèse : la morale est une création de la société. C'est la société qui, pour assurer sa cohésion, inculque à chaque individu des règles : respecter les anciens, ne pas voler, tenir sa parole. L'individu reçoit la morale de son groupe (famille, village, école, nation) comme il reçoit sa langue. La force de l'obligation vient du fait que la société est « plus grande » que l'individu et le précède.

Mais cette thèse conduit à une critique redoutable : le \emph{relativisme}. Si la morale est un produit social, alors chaque société a sa morale, et aucune n'est supérieure aux autres. Ce qui est bien ici serait mal ailleurs ; ce qui était bien hier (l'esclavage, admis pendant des siècles) serait devenu mal aujourd'hui. Or la conscience morale proteste : l'esclavage n'a jamais été juste, même quand toutes les sociétés le pratiquaient. Le fondement social explique les mœurs, mais semble incapable de fonder le devoir véritable.

\textbf{4. Le fondement rationnel : la raison seule fonde la morale.} Enfin, le philosophe Emmanuel \emph{Kant} (1724-1804) soutient que la morale ne vient ni de Dieu, ni de la nature, ni de la société, mais de la \emph{raison} elle-même. C'est parce que je suis un être raisonnable que je peux me donner à moi-même une loi universelle : « agis de telle sorte que ta conduite puisse valoir pour tous les hommes ». La morale devient alors \emph{autonome} : l'homme n'obéit plus à une loi étrangère, il obéit à la loi que sa propre raison se prescrit. Cette thèse, appelée \emph{morale du devoir}, fera l'objet de la section 6 de notre leçon.

\begin{popfigure}
\begin{center}
\begin{tikzpicture}[
  grow=down,
  level distance=1.6cm,
  level 1/.style={sibling distance=3.1cm},
  every node/.style={draw, rounded corners, align=center, font=\small},
  edge from parent/.style={draw=popmuted, thick}
]
\node[fill=popblue!80, text=white, font=\bfseries] {Fondements de la morale}
  child { node[fill=popgoldL, align=center]{Dieu\\ \emph{(religion)}}
    child { node[fill=white, font=\scriptsize, align=center]{Ex. : les Dix\\ Commandements} }
  }
  child { node[fill=popgreenL, align=center]{Nature\\ \emph{(loi naturelle)}}
    child { node[fill=white, font=\scriptsize, align=center]{Ex. : stoïcisme,\\ vivre selon la nature} }
  }
  child { node[fill=poporangeL, align=center]{Société\\ \emph{(Durkheim)}}
    child { node[fill=white, font=\scriptsize, align=center]{Ex. : règles\\ inculquées par le groupe} }
  }
  child { node[fill=poppurpleL, align=center]{Raison\\ \emph{(Kant)}}
    child { node[fill=white, font=\scriptsize, align=center]{Ex. : la loi que\\ je me donne moi-même} }
  };
\end{tikzpicture}
\end{center}
\legende{Les quatre grands fondements proposés pour la morale : Dieu (fondement religieux), la Nature (loi naturelle), la Société (fondement social, Durkheim) et la Raison (fondement rationnel, Kant).}
\end{popfigure}

\courssub{Conclusion provisoire : hétéronomie et autonomie}

Ce débat sur le fondement peut se résumer par une grande opposition. D'un côté, les morales dites \cle{hétéronomes} (du grec \emph{heteros}, « autre », et \emph{nomos}, « loi ») : la loi morale vient \emph{d'ailleurs} que de l'homme — de Dieu, de la nature ou de la société. L'homme reçoit la loi. De l'autre côté, les morales dites \cle{autonomes} (\emph{autos}, « soi-même ») : la loi morale est celle que l'homme se donne \emph{lui-même}, par sa raison. C'est la position de Kant.

\begin{aretenir}[L'essentiel de la section]
\begin{itemize}
\item La morale (latin \emph{mores}) et l'éthique (grec \emph{ethos}) désignent la même réalité : les coutumes, la manière d'agir.
\item La morale est l'ensemble des règles de conduite et des valeurs (bien/mal) qui orientent les actions humaines ; elle prescrit ce qui doit être fait.
\item Elle se distingue des mœurs (usages effectifs), du droit (contrainte extérieure sanctionnée par l'État) et de la religion (foi et culte).
\item Le problème du fondement : pourquoi obéir à la morale ? Quatre réponses : Dieu, la nature, la société (Durkheim), la raison (Kant).
\item Opposition finale : morales hétéronomes (loi reçue d'ailleurs) contre morales autonomes (loi que la raison se donne).
\end{itemize}
\end{aretenir}

\begin{exoresolu}[Appliquer les distinctions à un cas concret]
\textbf{Énoncé.} À Bafoussam, un élève découvre que son camarade a copié l'épreuve de mathématiques grâce à un téléphone caché. Il hésite à le dénoncer. Identifiez dans cette situation : (a) ce qui relève de la morale ; (b) ce qui relève du droit ou du règlement ; (c) le fondement moral invoqué si l'élève se dit : « même si personne ne le sait jamais, tricher est mal ».
\tcblower
\textbf{Solution détaillée.} (a) \emph{La morale} : tricher est une faute morale, car c'est mentir sur son travail et manquer d'honnêteté envers soi-même et envers les autres candidats ; la conscience de l'élève témoin le juge intérieurement (remords possible s'il se tait). (b) \emph{Le droit / le règlement} : la fraude à un examen officiel est sanctionnée extérieurement par le règlement scolaire et, au Baccalauréat, par la réglementation de l'OBC (annulation de l'épreuve, exclusion) : c'est une contrainte publique, non une simple affaire de conscience. (c) \emph{Le fondement} : dire « même si personne ne le sait, c'est mal », c'est refuser le critère de l'intérêt et de la sanction extérieure ; c'est invoquer une loi intérieure universelle, donc un fondement \emph{rationnel} (kantien) ou, selon le cas, \emph{naturel} (la conscience inscrite dans la nature humaine) — en tout cas une morale qui ne dépend pas d'être vu ou puni.
\end{exoresolu}

\begin{exercice}[Pour s'entraîner]
1. Définis la morale et distingue-la des mœurs en t'appuyant sur un exemple tiré de la vie quotidienne camerounaise. 2. Un chauffeur de taxi-moto grille un feu rouge la nuit, sans danger et sans policier : a-t-il commis une faute morale, une faute juridique, les deux ? Justifie. 3. « La morale change avec les sociétés, donc elle n'a aucun fondement solide. » Discute cette affirmation en mobilisant la thèse de Durkheim et sa critique.
\end{exercice}

\coursec{Les morales de l'intérêt : présentation générale}

Nous avons vu dans la section précédente que la morale est l'ensemble des règles et des valeurs qui orientent nos actions vers le bien, et que toute doctrine morale repose sur un \cle{fondement} : ce qui justifie d'agir bien plutôt que mal. Mais sur quoi fonder la morale ? Deux grandes réponses s'affrontent dans l'histoire de la philosophie. Pour certains penseurs, la morale repose sur le \cle{devoir}, c'est-à-dire sur une obligation qui s'impose à nous indépendamment de nos désirs et de nos avantages. Pour d'autres, au contraire, la morale se fonde sur l'\cle{intérêt} : on agit bien parce que bien agir nous procure quelque chose — le plaisir, le bonheur ou l'utilité. C'est cette seconde famille de doctrines que nous allons étudier maintenant, avant de voir ensuite comment Kant lui oppose la morale du devoir.

\courssub{Qu'est-ce qu'une morale de l'intérêt ?}

Commençons par une expérience de pensée très simple. Demandez à un élève de Première technique pourquoi il travaille sérieusement à l'atelier. Il répondra peut-être : « Pour réussir mon Probatoire et trouver un bon emploi. » Demandez à un commerçant de Douala pourquoi il est honnête avec ses clients. Il répondra peut-être : « Parce qu'un client satisfait revient, et cela fait vivre ma boutique. » Dans les deux cas, l'action bonne (travailler, être honnête) est accomplie \emph{en vue d'un avantage} pour celui qui agit. C'est exactement la logique des morales de l'intérêt : le bien n'est pas recherché pour lui-même, mais parce qu'il sert une fin qui nous profite.

\begin{definition}[Morale de l'intérêt]
Une \cle{morale de l'intérêt} est une doctrine qui fait du \textbf{bonheur} ou de l'\textbf{utilité} le \emph{fondement} et la \emph{fin} de l'action morale. Autrement dit : agir moralement, c'est agir en vue d'une fin qui profite à l'agent (ou à la communauté) — le plaisir, le bonheur, l'avantage. Le critère du bien n'est pas le devoir abstrait, mais la \textbf{conséquence heureuse} de l'action.
\end{definition}

Le principe commun à toutes ces doctrines peut se formuler ainsi : \emph{agir bien, c'est agir en vue d'une fin qui nous profite}. Le mot « intérêt » vient du latin \emph{inter-esse}, « être parmi, être au milieu de » : mon intérêt, c'est ce qui me concerne, ce qui me touche, ce qui m'importe. Une morale de l'intérêt affirme donc que l'homme n'agit jamais sans raison, et que la raison d'agir moralement doit être cherchée dans ce que l'action rapporte à celui qui agit.

\begin{attention}[Ne pas confondre « intérêt » et « égoïsme »]
C'est le piège le plus fréquent dans les copies. Dire qu'une morale est une morale de l'\cle{intérêt} ne signifie pas qu'elle est \cle{égoïste}. L'égoïsme ne recherche que \emph{mon} bonheur personnel, au mépris des autres. Or l'utilitarisme, par exemple, est une morale de l'intérêt, mais de l'intérêt \textbf{de tous} : il commande de rechercher le plus grand bonheur du plus grand nombre. Il faut donc distinguer l'\textbf{intérêt égoïste} (mon seul bonheur) et l'\textbf{intérêt élargi} (le bonheur de tous, dans lequel le mien est inclus). Une morale de l'intérêt peut donc être profondément généreuse.
\end{attention}

\courssub{Les trois formes de la morale de l'intérêt}

Au programme de Première technique, trois grandes doctrines de l'intérêt sont étudiées. Elles partagent le même principe (la morale vise une fin avantageuse), mais diffèrent par la \emph{nature} de cette fin et par \emph{celui qui en profite}.

\textbf{1. L'épicurisme : la morale du plaisir.} Pour Épicure (341-270 av. J.-C.), philosophe grec fondateur de l'école du Jardin à Athènes, la fin de la vie morale est le \cle{plaisir}, entendu non comme débauche, mais comme absence de souffrance du corps (\emph{aponie}) et tranquillité de l'âme (\emph{ataraxie}). Le sage épicurien choisit avec soin les plaisirs naturels et nécessaires (manger, boire, l'amitié) et renonce aux plaisirs vains qui engendrent des souffrances. Ici, l'intérêt recherché est d'abord \emph{individuel} : c'est ma propre tranquillité.

\textbf{2. L'eudémonisme : la morale du bonheur.} Le mot vient du grec \emph{eudaimonia}, « bonheur, vie accomplie ». Pour Aristote (384-322 av. J.-C.), tous les hommes recherchent le \cle{bonheur} comme fin suprême : on veut être heureux non pour autre chose, mais pour lui-même. Le bonheur véritable s'obtient par la pratique des vertus (courage, justice, tempérance) dans une vie achevée. L'eudémonisme occupe une position intermédiaire : le bonheur est personnel, mais il ne peut s'atteindre qu'en vivant vertueusement \emph{avec les autres}, dans la cité.

\textbf{3. L'utilitarisme : la morale de l'utilité collective.} Développé par les philosophes anglais Jeremy Bentham (1748-1832) et John Stuart Mill (1806-1873), l'\cle{utilitarisme} affirme qu'une action est bonne si elle produit \emph{le plus grand bonheur pour le plus grand nombre}. Le critère moral est l'\cle{utilité} : on évalue une action à ses conséquences sur l'ensemble de la société. Ici, l'intérêt est pleinement \emph{élargi} : mon propre bonheur ne compte pas plus que celui d'un autre.

\begin{popfigure}
\begin{center}
\begin{tikzpicture}[scale=1.0]
% Axe principal
\draw[->,very thick,popdark] (0,0) -- (12.5,0);
% Graduations
\foreach \x in {1.5,6,10.5} {\draw[thick,popdark] (\x,-0.15) -- (\x,0.15);}
% Points doctrines
\fill[poporange] (1.5,0) circle (0.14);
\fill[poppurple] (6,0) circle (0.14);
\fill[popblue] (10.5,0) circle (0.14);
% Étiquettes doctrines
\node[above,poporange,font=\bfseries] at (1.5,0.25) {Épicurisme};
\node[below,poporange,font=\small] at (1.5,-0.25) {le plaisir};
\node[above,poppurple,font=\bfseries] at (6,0.25) {Eudémonisme};
\node[below,poppurple,font=\small] at (6,-0.25) {le bonheur};
\node[above,popblue,font=\bfseries] at (10.5,0.25) {Utilitarisme};
\node[below,popblue,font=\small] at (10.5,-0.25) {l'utilité collective};
% Extrémités de l'axe
\node[left,popmuted,font=\small\itshape,align=right] at (0,0.9) {Intérêt individuel};
\node[left,popmuted,font=\small\itshape,align=right] at (0,0.5) {(mon bonheur)};
\node[right,popmuted,font=\small\itshape,align=left] at (12.5,0.9) {Intérêt collectif};
\node[right,popmuted,font=\small\itshape,align=left] at (12.5,0.5) {(le bonheur de tous)};
% Flèche de sens
\draw[->,thick,popgreen] (2.2,1.5) -- (9.8,1.5) node[midway,above,font=\small\itshape,popgreen] {élargissement croissant de l'intérêt};
\end{tikzpicture}
\end{center}
\legende{Les trois morales de l'intérêt sur un axe allant de l'intérêt individuel à l'intérêt collectif : l'épicurisme vise d'abord le plaisir personnel du sage, l'eudémonisme le bonheur accompli dans la vie en communauté, l'utilitarisme le bonheur du plus grand nombre.}
\end{popfigure}

Ce schéma montre l'essentiel : les trois doctrines appartiennent à la même famille (la fin recherchée est toujours avantageuse), mais elles se situent à des degrés différents sur l'axe qui va de \emph{moi seul} à \emph{tous}. Nous les étudierons en détail dans les sections suivantes.

\courssub{Comment analyser une doctrine morale ?}

Face à n'importe quelle doctrine morale, à l'examen comme dans la vie, il faut savoir la décomposer méthodiquement. C'est un savoir-faire exigible : « rédiger et justifier une démarche ».

\begin{methode}[Analyser une doctrine morale en quatre questions]
\begin{enumerate}
\item \textbf{Identifier la fin poursuivie} : que cherche-t-on en agissant bien ? (le plaisir ? le bonheur ? l'utilité ? l'accomplissement du devoir ?)
\item \textbf{Identifier le fondement} : sur quoi repose l'obligation d'agir ainsi ? (la nature humaine ? le calcul des conséquences ? la raison ?)
\item \textbf{Dégager le critère de l'action bonne} : comment reconnaît-on qu'une action est morale ? (elle procure le plaisir durable ? elle rend heureux ? elle maximise le bonheur collectif ?)
\item \textbf{Illustrer par un exemple de la vie courante} : appliquer la doctrine à une situation concrète (un choix professionnel, un acte d'honnêteté, une décision publique) pour vérifier sa cohérence.
\end{enumerate}
\end{methode}

\begin{exemplebox}[Application de la méthode à l'utilitarisme]
Prenons l'exemple d'un maire camerounais qui hésite entre construire un stade de prestige ou un centre de santé dans sa commune.
\begin{enumerate}
\item \emph{Fin poursuivie} : le plus grand bonheur du plus grand nombre.
\item \emph{Fondement} : l'utilité, c'est-à-dire le calcul des conséquences heureuses pour la population.
\item \emph{Critère} : l'action bonne est celle qui produit le maximum de bien-être collectif.
\item \emph{Application} : le centre de santé soignera chaque année des milliers d'habitants, tandis que le stade ne servira qu'occasionnellement. L'utilitariste choisit donc le centre de santé, même si le stade flatterait davantage sa gloire personnelle.
\end{enumerate}
\end{exemplebox}

\courssub{Exercice résolu : reconnaître une morale de l'intérêt}

\begin{exoresolu}[Identifier la doctrine à partir d'une situation]
\textbf{Énoncé.} Trois élèves expliquent pourquoi ils ont aidé un camarade malade à rattraper ses cours :
\begin{itemize}
\item Amina : « Je l'ai aidé parce que cela m'a procuré une vraie joie tranquille, sans trouble ; l'amitié est le plus sûr des plaisirs. »
\item Boris : « Je l'ai aidé parce qu'être généreux fait partie d'une vie réussie et pleinement heureuse ; on ne peut pas s'épanouir seul. »
\item Carine : « Je l'ai aidé parce que si tout le monde s'entraide, toute la classe progressera, et c'est le bien du plus grand nombre. »
\end{itemize}
Attribuez à chaque élève la doctrine morale qui correspond à son propos, en justifiant.
\tcblower
\textbf{Solution détaillée.}
\begin{itemize}
\item \textbf{Amina relève de l'épicurisme.} La fin qu'elle invoque est le plaisir, compris comme joie tranquille « sans trouble » : c'est l'ataraxie d'Épicure. Le fait qu'elle valorise l'amitié comme « le plus sûr des plaisirs » confirme l'analyse, car Épicure faisait de l'amitié un pilier de la vie heureuse. L'intérêt visé est individuel : sa propre tranquillité.
\item \textbf{Boris relève de l'eudémonisme.} Il parle de « vie réussie », de bonheur pleinement accompli, et il lie ce bonheur à la pratique d'une vertu (la générosité) et à la vie avec autrui. C'est exactement la conception aristotélicienne : le bonheur comme accomplissement de soi par la vertu, dans la communauté.
\item \textbf{Carine relève de l'utilitarisme.} Son critère est explicite : « le bien du plus grand nombre ». Elle évalue son action à ses conséquences collectives (le progrès de toute la classe), ce qui est la définition même de l'utilité comme critère moral.
\end{itemize}
On remarque que les trois élèves accomplissent la \emph{même action} (aider un camarade), mais la \emph{justifient} différemment : c'est bien le fondement qui distingue les doctrines, non le geste extérieur.
\end{exoresolu}

\courssub{Problématisation : l'intérêt peut-il fonder la morale ?}

Avant d'étudier chaque doctrine en détail, posons dès maintenant la question critique qui traversera toute l'unité : \emph{peut-on fonder la morale sur l'intérêt sans la dissoudre dans le calcul égoïste ?}

Le reproche est puissant. Si j'agis bien \emph{parce que cela m'est avantageux}, alors ma conduite ressemble à un calcul commercial : je suis honnête tant que l'honnêteté « rapporte ». Mais que se passe-t-il le jour où le mensonge ou la tricherie deviennent plus profitables ? Le commerçant honnête par intérêt trompera ses clients dès qu'il pourra le faire impunément. L'élève qui travaille par intérêt trichera à l'examen si la fraude est plus sûre que le travail. La morale fondée sur l'intérêt semble donc fragile : elle dépend des circonstances, et elle peut justifier l'immoralité dès que celle-ci devient avantageuse.

C'est précisément ce reproche que Kant adressera aux morales de l'intérêt : elles rendent la morale \emph{conditionnelle} (« agis bien \emph{si} tu veux être heureux »), alors qu'une véritable obligation morale devrait être \emph{inconditionnelle} (« agis bien, parce que c'est ton devoir »). La morale du devoir, que nous étudierons à la fin de cette leçon, se présentera comme la réponse à cette difficulté.

Mais les partisans de l'intérêt ont des arguments à opposer. D'abord, disent-ils, il est \emph{réaliste} de partir de ce que les hommes désirent réellement — le bonheur — plutôt que d'exiger d'eux un sacrifice sans raison. Ensuite, l'intérêt bien compris peut coïncider avec la vertu : l'épicurien découvre que les plaisirs désordonnés ruinent le bonheur, l'eudémoniste que nul ne peut être heureux en étant injuste, l'utilitariste que mon bonheur dépend du bonheur des autres. L'intérêt élargi pourrait donc conduire aux mêmes conduites que le devoir, par un chemin différent.

\begin{aretenir}[L'essentiel de la section]
\begin{itemize}
\item Une \cle{morale de l'intérêt} fait du bonheur ou de l'utilité le fondement et la fin de l'action morale : agir bien, c'est agir en vue d'une fin qui profite.
\item Trois formes au programme : l'\textbf{épicurisme} (le plaisir, intérêt individuel), l'\textbf{eudémonisme} (le bonheur accompli par la vertu, position intermédiaire), l'\textbf{utilitarisme} (l'utilité, intérêt collectif).
\item Intérêt $\neq$ égoïsme : l'utilitarisme recherche l'intérêt \emph{de tous}.
\item Méthode d'analyse : fin poursuivie, fondement, critère de l'action bonne, exemple concret.
\item Question critique : l'intérêt peut-il fonder une obligation stable, ou la morale risque-t-elle de se dissoudre dans le calcul égoïste ? C'est à cette objection que répondra la morale du devoir de Kant.
\end{itemize}
\end{aretenir}

Nous pouvons maintenant examiner chacune de ces doctrines en détail, en commençant par la plus ancienne et la plus mal comprise : l'épicurisme, morale du plaisir.

\coursec{L'épicurisme : la morale du plaisir}

Nous venons de voir que les morales de l'intérêt partagent une intuition commune : l'action morale n'est jamais désintéressée, car l'homme agit toujours en vue de son propre bien, de son bonheur, de son avantage. Mais une question surgit immédiatement : quel est donc ce bien suprême que tout homme recherche ? À cette question, l'épicurisme apporte la réponse la plus directe et la plus provocante : c'est le \cle{plaisir}. Cette doctrine, née en Grèce au III\textsuperscript{e} siècle avant Jésus-Christ, est la première grande morale de l'intérêt que nous devons étudier en détail. Encore faut-il la comprendre correctement, car elle est sans doute la philosophie la plus malmenée par le langage courant.

\courssub{Épicure et l'école du Jardin : présentation du philosophe}

\cle{Épicure} (341-270 av. J.-C.) est un philosophe grec, né sur l'île de Samos, qui s'installe à Athènes vers 306 av. J.-C. Il y fonde une école de philosophie appelée le \cle{Jardin}, parce que les cours se donnaient dans le jardin de sa maison. Ce détail n'est pas anodin : contrairement aux écoles prestigieuses de son temps, le Jardin d'Épicure accueillait tout le monde, y compris les femmes et les esclaves, ce qui était révolutionnaire dans la Grèce antique. La philosophie n'y était pas un savoir réservé à une élite, mais un art de vivre, une médecine de l'âme accessible à tous.

Épicure écrit énormément (on lui attribue environ trois cents ouvrages), mais presque tout est perdu. Nous connaissons sa pensée grâce à trois lettres conservées par Diogène Laërce, dont la célèbre \emph{Lettre à Ménécée}, véritable petit manuel de sagesse, et grâce au poème du philosophe latin Lucrèce, \emph{De la nature des choses}.

\begin{popfigure}
\begin{center}
\begin{tikzpicture}[x=1cm,y=1cm]
\draw[popline, thick, -{Stealth}] (-0.5,0) -- (12.5,0);
\foreach \x/\lab in {0/{-500}, 2.5/{-400}, 5/{-341}, 7.5/{-270}, 10/{0}}{
  \draw[popline] (\x,0.12) -- (\x,-0.12);
  \node[below, popmuted, font=\small] at (\x,-0.15) {\lab};
}
\fill[poporange] (5,0) circle (2.5pt);
\node[above, popdark, font=\small\bfseries] at (5,0.35) {Naissance d'Épicure (-341)};
\fill[poporange] (7.5,0) circle (2.5pt);
\node[above, popdark, font=\small\bfseries] at (7.5,0.35) {Mort d'Épicure (-270)};
\draw[poporange, very thick] (5,0.05) -- (7.5,0.05);
\node[below, poporange, font=\small\itshape] at (6.25,-0.9) {Vie d'Épicure};
\draw[popblue, very thick, -{Stealth}] (5.6,0.55) -- (5.6,1.15);
\node[above, popblue, font=\small] at (5.6,1.15) {Fondation du Jardin à Athènes (vers -306)};
\node[popmuted, font=\small\itshape, anchor=west] at (-0.5,1.9) {Antiquité grecque};
\end{tikzpicture}
\end{center}
\legende{Frise chronologique situant Épicure dans l'Antiquité grecque (échelle indicative).}
\end{popfigure}

\courssub{La thèse centrale : le plaisir est le souverain bien}

Commençons par une intuition simple. Observez un nourrisson : dès la naissance, sans qu'on le lui apprenne, il recherche ce qui lui fait plaisir (le lait, la chaleur, les bras de sa mère) et il fuit ce qui lui fait mal (la faim, le froid, la douleur). Aucun adulte ne lui a enseigné cela. Épicure part de cette observation : la nature elle-même nous montre que le plaisir est le bien premier et que la douleur est le mal premier. Tout être vivant poursuit spontanément le plaisir et fuit la souffrance.

De cette constatation, Épicure tire sa thèse morale fondamentale, qu'il formule dans la \emph{Lettre à Ménécée} : « le plaisir est le commencement et la fin de la vie heureuse ». Autrement dit, le plaisir est à la fois le \emph{point de départ} de toute nos actions (c'est lui que nous cherchons d'abord) et leur \emph{but ultime} (c'est en lui que consiste la vie heureuse). Le plaisir est donc le \cle{souverain bien}, c'est-à-dire le bien suprême auquel tous les autres biens se subordonnent.

\begin{definition}[L'épicurisme]
L'\cle{épicurisme} est la doctrine morale fondée par Épicure, qui fait du \cle{plaisir}, compris comme \emph{absence de souffrance} (du corps et de l'âme), le souverain bien et la fin dernière de toutes nos actions. Vivre moralement, c'est vivre agréablement ; mais vivre agréablement exige la prudence, la sobriété et la sagesse.
\end{definition}

\begin{attention}[Erreur fréquente]
L'erreur la plus répandue consiste à réduire l'épicurisme à la recherche des plaisirs grossiers : boire, manger sans fin, multiplier les aventures. C'est exactement le contraire de la pensée d'Épicure ! Le philosophe recommande au contraire la \cle{sobriété} : il vivait de pain, d'eau et d'un peu de fromage, et affirmait qu'un tel régime suffit au bonheur. Dans une dissertation, écrire « l'épicurisme est la doctrine de la débauche » est une faute grave qui révèle une méconnaissance du cours.
\end{attention}

\courssub{La clarification majeure : le vrai plaisir est un plaisir stable}

Voici le cœur de la doctrine, et sa partie la plus subtile. Épicure distingue deux sortes de plaisirs :

\begin{itemize}
\item le \cle{plaisir en mouvement} (ou plaisir kinétique) : c'est le plaisir vif mais bref que procure la satisfaction d'un désir (manger quand on a faim, boire quand on a soif, danser, faire la fête). Ce plaisir est instable : il disparaît dès que le désir est satisfait, et il peut même engendrer de nouvelles souffrances (la gueule de bois après la fête, les dettes après les dépenses) ;
\item le \cle{plaisir stable} (ou plaisir katastématique) ou plaisir en repos : c'est l'état durable où le corps ne souffre pas et où l'âme n'est pas troublée. C'est ce plaisir-là, et non le premier, qu'Épicure déclare être le vrai souverain bien.
\end{itemize}

Ce plaisir stable se décompose en deux états que le sage doit atteindre :

\begin{itemize}
\item l'\cle{aponie} : l'absence de douleur du corps (du grec \emph{a-}, privation, et \emph{ponos}, douleur) ;
\item l'\cle{ataraxie} : l'absence de trouble de l'âme.
\end{itemize}

\begin{definition}[L'ataraxie]
L'\cle{ataraxie} (du grec \emph{ataraxia}, absence de trouble) est l'état de tranquillité parfaite de l'âme, dans lequel celle-ci n'est agitée ni par la peur (de la mort, des dieux, de l'avenir), ni par les désirs insatisfaits, ni par les passions. C'est cet état de sérénité que le sage épicurien recherche comme le plaisir le plus haut et le plus durable.
\end{definition}

On comprend alors le paradoxe épicurien : le philosophe du plaisir est en réalité un philosophe de la \emph{modération}. Pourquoi ? Parce que les plaisirs intenses et coûteux rendent dépendant, créent des besoins toujours nouveaux, et finissent par troubler l'âme. Celui qui s'habitue au luxe souffre dès que le luxe lui manque ; celui qui se contente de peu ne manque de rien. Épicure écrit : « le pain et l'eau procurent le plaisir le plus vif, dès lors qu'on les porte à des lèvres affamées ». Le plaisir le plus sûr n'est donc pas celui qu'on ajoute, mais celui qu'on obtient en supprimant la souffrance.

\courssub{La classification des désirs : le guide pratique de la sagesse}

Comment savoir quels désirs satisfaire et lesquels combattre ? Épicure propose une classification célèbre, véritable outil pratique de la sagesse. Il distingue trois catégories de désirs :

\begin{enumerate}
\item les désirs \cle{naturels et nécessaires} : manger, boire, dormir, se protéger du froid. Sans leur satisfaction, le corps souffre et la vie elle-même est menacée. Il faut les satisfaire, et ils sont faciles à satisfaire ;
\item les désirs \cle{naturels mais non nécessaires} : les mets raffinés, les boissons coûteuses, les plaisirs variés. Ils sont naturels (ils prolongent les désirs précédents), mais on peut vivre sans eux. Le sage les satisfait avec modération, sans en devenir esclave ;
\item les désirs \cle{vains} (ou non naturels et non nécessaires) : la gloire, la richesse sans limite, le pouvoir, l'immortalité. Ces désirs n'ont pas de fondement naturel : ils naissent de l'opinion des autres et de l'imagination. Surtout, ils sont \emph{insatiables} : on n'en voit jamais le bout. Le sage s'en débarrasse complètement.
\end{enumerate}

\begin{popfigure}
\begin{center}
\begin{tikzpicture}[
  case/.style={draw=popline, rounded corners=4pt, fill=#1, text width=3.6cm, minimum height=4.6cm, align=center, inner sep=6pt},
  titre/.style={font=\small\bfseries, popdark},
  ex/.style={font=\footnotesize\itshape, popmuted}
]
\node[case=popgreenL, anchor=north west, align=center] at (0,0){
  {\bfseries Désirs naturels\\et nécessaires}\\[4pt]
  \footnotesize À satisfaire : sans eux, le corps souffre.\\[6pt]
  \itshape Ex. camerounais : manger son ndolé ou son eru, boire de l'eau, dormir, se loger décemment.
};
\node[case=popgoldL, anchor=north west, align=center] at (4.35,0){
  {\bfseries Désirs naturels\\mais non nécessaires}\\[4pt]
  \footnotesize À modérer : agréables, mais on peut s'en passer.\\[6pt]
  \itshape Ex. camerounais : le champagne au mariage, les mets raffinés du restaurant, le téléphone haut de gamme.
};
\node[case=poppinkL, anchor=north west, align=center] at (8.7,0){
  {\bfseries Désirs vains\\(non naturels)}\\[4pt]
  \footnotesize À éliminer : insatiables, ils troublent l'âme.\\[6pt]
  \itshape Ex. camerounais : courir après la gloire, « écraser » les autres par l'étalage de richesse, vouloir posséder sans limite.
};
\end{tikzpicture}
\end{center}
\legende{La classification épicurienne des désirs, avec un exemple camerounais pour chaque catégorie.}
\end{popfigure}

La \cle{sagesse épicurienne} consiste donc en une véritable diététique des désirs : satisfaire simplement les désirs naturels et nécessaires, modérer les désirs naturels non nécessaires, et extirper de son âme les désirs vains. Celui qui applique cette méthode devient maître de son bonheur, car il ne dépend plus de la fortune ni du regard des autres.

\begin{exemplebox}[La sagesse épicurienne au quotidien]
Se contenter d'un repas simple — un bon ndolé et de l'eau fraîche — partagé joyeusement avec des amis, plutôt que de s'endetter pour organiser un festin ostentatoire destiné à impressionner le quartier : voilà une attitude parfaitement épicurienne. Le premier choix procure un plaisir stable et sans lendemain douloureux ; le second achète un plaisir bref au prix de mois de troubles financiers, donc d'angoisses, donc de perte de l'ataraxie.
\end{exemplebox}

\courssub{L'amitié : le plus sûr des plaisirs}

On pourrait croire que la recherche de son propre plaisir conduit à l'isolement. C'est l'inverse qu'enseigne Épicure : l'\cle{amitié} occupe une place centrale dans sa doctrine. Il déclare que « de tout ce que la sagesse procure pour le bonheur de la vie entière, le plus grand de beaucoup est la possession de l'amitié ».

Pourquoi l'amitié est-elle si précieuse ? Pour deux raisons. D'abord, elle est en elle-même l'un des plaisirs les plus doux et les plus stables : la conversation entre amis, la confiance, la joie partagée ne coûtent rien et ne troublent pas l'âme. Ensuite, l'amitié est le meilleur \emph{rempart contre les malheurs} : l'ami vous secourt dans la maladie, vous aide dans la pauvreté, vous console dans le deuil. Celui qui a de vrais amis n'a pas besoin d'accumuler des richesses pour se sentir en sécurité. C'est pourquoi le Jardin d'Épicure fonctionnait comme une communauté d'amis vivant ensemble, sobrement, loin des agitations de la vie politique.

\begin{savaistu}[Un mot qui a trahi son maître]
Aujourd'hui, le mot « épicurien » désigne couramment un amateur de bonne chère, de vin et de festins. Cet usage trahit presque totalement la pensée d'Épicure, qui fut un homme d'une frugalité exemplaire, se nourrissant de pain, d'eau et d'un peu de fromage, et écrivant à un ami : « envoie-moi un petit pot de fromage, afin que je puisse faire bombance quand je le voudrai ». Le luxe d'Épicure, c'était un pot de fromage ! Voilà qui devrait nous rendre prudents face aux étiquettes que le langage courant colle aux philosophies.
\end{savaistu}

\courssub{Application aux situations concrètes : le « mbenguiste » et le sage}

Appliquons maintenant la doctrine à une situation bien connue de la vie sociale camerounaise. Considérons le personnage du « \cle{mbenguiste} » ou du « sapeur » : celui qui court après les signes extérieurs de richesse — vêtements de marque, téléphone dernier cri, voiture de prestige, dépenses spectaculaires dans les cérémonies — souvent au prix de dettes écrasantes et d'une dépendance totale au regard des autres.

Analysons ce cas avec les outils d'Épicure :

\begin{itemize}
\item le désir de paraître riche est un désir \emph{vain} : il ne répond à aucun besoin naturel, il naît de l'opinion, et il est insatiable (il y aura toujours quelqu'un de plus riche à impressionner) ;
\item sa satisfaction procure un plaisir \emph{en mouvement}, bref et fragile : l'admiration du public dure un soir, les mensualités du crédit durent des années ;
\item ce mode de vie détruit l'ataraxie : peur de perdre la face, angoisse des dettes, jalousie, dépendance à l'égard du jugement d'autrui.
\end{itemize}

À l'opposé, le sage épicurien camerounais se contente de peu : un travail honnête, une nourriture simple, un toit, et surtout des amis véritables. Son bonheur ne dépend ni de la fortune ni de l'applaudissement. Il est libre. Épicure dirait du mbenguiste qu'il est riche en apparence et pauvre en réalité, car « rien ne suffit à celui pour qui le suffisant est peu ».

\begin{exoresolu}[Analyser une situation à la lumière de l'épicurisme]
Énoncé : Paul, jeune fonctionnaire à Yaoundé, emprunte une forte somme pour acheter une voiture de luxe afin d'être respecté dans son quartier. Il passe ensuite ses week-ends à craindre les accidents, les vols et les échéances du crédit, et néglige ses amis d'enfance qu'il juge « trop modestes ». À l'aide de la classification épicurienne des désirs, analysez le comportement de Paul et dites ce que lui conseillerait Épicure.
\tcblower
Solution détaillée :
\begin{enumerate}
\item \textbf{Identifier la nature du désir.} Le désir d'être respecté par l'étalage d'une voiture de luxe n'est ni naturel ni nécessaire : on peut vivre et être heureux sans elle. C'est un désir \emph{vain}, fondé sur l'opinion des autres et insatiable (après la voiture, il faudra la villa, puis une voiture plus chère encore).
\item \textbf{Évaluer le plaisir obtenu.} Le respect du quartier procure un plaisir en mouvement, bref et instable, acheté au prix de souffrances durables : dettes, angoisses, peurs quotidiennes. Le compte du plaisir est négatif : Paul perd l'aponie (soucis matériels) et l'ataraxie (troubles de l'âme).
\item \textbf{Repérer l'erreur sur l'amitié.} En dédaignant ses amis modestes, Paul détruit le plus sûr des plaisirs et le meilleur rempart contre les malheurs, selon Épicure.
\item \textbf{Formuler le conseil épicurien.} Épicure conseillerait à Paul de renoncer aux désirs vains (le paraître), de se contenter d'un moyen de transport simple et adapté à ses revenus, et de cultiver l'amitié véritable. Il retrouverait ainsi la tranquillité de l'âme, seul fondement du plaisir stable.
\end{enumerate}
\end{exoresolu}

\courssub{Discussion critique : le plaisir peut-il fonder le devoir envers autrui ?}

L'épicurisme est séduisant par sa simplicité et sa sagesse pratique. Mais il appelle une critique décisive, qui prépare la suite de notre leçon : si le plaisir est le souverain bien, que faire lorsque mon plaisir entre en conflit avec le bien d'autrui ?

Posons le problème concrètement. Si un mensonge me procure un avantage sans me causer de trouble, l'épicurisme peut-il me l'interdire ? Si je dois choisir entre sauver un inconnu au péril de ma tranquillité et rester tranquillement chez moi, la recherche de l'ataraxie ne me commande-t-elle pas l'égoïsme ? Certes, Épicure répond que l'injustice trouble toujours l'âme (peur d'être découvert, remords) et que l'amitié exige la fidélité. Mais cette réponse reste fragile : elle fait dépendre la justice d'un calcul de plaisirs, et un calcul peut parfois pencher du côté de l'injustice discrète.

Autrement dit, le plaisir, même raffiné et stable, semble incapable de fonder un \cle{devoir} véritable envers autrui, c'est-à-dire une obligation inconditionnelle qui s'imposerait à moi même contre mon intérêt. Le risque d'\cle{égoïsme} — chacun cherchant d'abord sa propre tranquillité — demeure la faiblesse de toute morale de l'intérêt. C'est précisément pour répondre à cette difficulté que Kant construira une morale du devoir, fondée non sur le plaisir mais sur la raison : ce sera l'objet de la dernière partie de notre leçon.

\begin{aretenir}[L'essentiel de l'épicurisme]
\begin{itemize}
\item Épicure (341-270 av. J.-C.), fondateur du Jardin à Athènes, soutient que « le plaisir est le commencement et la fin de la vie heureuse ».
\item Le vrai plaisir n'est pas la débauche mais le plaisir \emph{stable} : l'aponie (absence de douleur du corps) et l'ataraxie (absence de trouble de l'âme).
\item La sagesse repose sur la classification des désirs : satisfaire les désirs naturels et nécessaires, modérer les désirs naturels non nécessaires, éliminer les désirs vains.
\item L'amitié est le plus sûr des plaisirs et le meilleur rempart contre les malheurs.
\item Limite de la doctrine : le plaisir fonde difficilement le devoir envers autrui ; risque d'égoïsme.
\end{itemize}
\end{aretenir}

\coursec{L'eudémonisme : la morale du bonheur}

Après avoir examiné l'épicurisme, qui identifie le bien au plaisir entendu comme absence de douleur (\emph{ataraxia}), nous nous tournons vers une doctrine plus ambitieuse : l'eudémonisme. Si Épicure invite le sage à se retirer du tumulte du monde pour y trouver la tranquillité, Aristote, au contraire, ancre la morale dans la \cle{vie politique} et l'\cle{activité}. Pour lui, le bonheur n'est pas un état passif que l'on subit, mais une \cle{activité} que l'on déploie tout au long d'une existence complète.

\courssub{Étymologie et thèse fondamentale : le bonheur comme fin suprême}

Le terme « eudémonisme » vient du grec \emph{eu} (bien) et \emph{daimôn} (génie, destin), d'où \emph{eudaimonia} : littéralement « bonne condition du génie », « vie bien dirigée », « félicité ». Contrairement à notre mot « bonheur » qui suggère souvent une chance passagère (\emph{bon-heur}), \emph{eudaimonia} désigne une \textbf{réussite de vie}, une existence accomplie dans sa totalité.

Aristote (384-322 av. J.-C.), dans l'\emph{Éthique à Nicomaque}, ouvre son enquête par une observation frappante : \emph{« Tout art, toute enquête, de même que toute action et tout choix, semble tendre vers quelque bien ; c'est avec raison qu'on a déclaré que le bien est ce vers quoi tout tend. »} Tous les hommes désirent être heureux. Mais si l'on demande en quoi consiste le bonheur, les réponses divergent : les uns disent le plaisir, d'autres l'honneur, d'autres la richesse. Aristote montre que ces biens sont \cle{choisis en vue d'autre chose} (on veut de l'argent pour acheter, on veut l'honneur pour être reconnu), alors que le bonheur est \cle{choisi pour lui-même} et jamais en vue d'autre chose. C'est la \cle{fin suprême} (\emph{telos}), la fin dernière qui donne sens à toutes nos actions intermédiaires.

\begin{definition}[L'eudémonisme]
Doctrine morale qui fait du \textbf{bonheur} (\emph{eudaimonia}), entendu comme vie accomplie et réussite de l'existence, la \cle{fin dernière} de l'action humaine. Toute action morale vise, médiatement ou immédiatement, cette fin suprême.
\end{definition}

\begin{aretenir}[Le bonheur, fin dernière et autosuffisante]
\begin{itemize}
    \item \textbf{Finalité :} on le désire pour lui-même, non pour autre chose.
    \item \textbf{Autosuffisance :} il rend la vie désirable et ne manque de rien (\emph{autarkês}).
    \item \textbf{Activité :} ce n'est pas un état (comme dormir), mais une \emph{energeia} (activité déployée) de l'âme.
\end{itemize}
\end{aretenir}

\courssub{Qu'est-ce que le bonheur ? Ni plaisir, ni honneur, ni richesse}

Aristote écarte successivement les candidats vulgaires au bonheur :
\begin{itemize}
    \item \textbf{La vie de plaisir} : elle est « servile » et « bestiale », esclave des sens, instable.
    \item \textbf{La vie politique} (recherche des honneurs) : l'honneur dépend de ceux qui le donnent, non de celui qui le reçoit ; il est trop extérieur et trop précaire.
    \item \textbf{La vie d'affaires} (recherche de la richesse) : l'argent n'est qu'un moyen, un instrument ; on ne le désire que pour s'en servir.
\end{itemize}
Le vrai bonheur est \textbf{« activité de l'âme conforme à la vertu »} (\emph{psuchês energeia kat' aretên}). L'âme (\emph{psuchê}) est le principe de vie ; sa fonction propre (\emph{ergon}) est l'activité rationnelle. Être heureux, c'est donc \textbf{fonctionner excellemment en tant qu'être rationnel}, c'est-à-dire vivre selon la raison et la vertu, et cela sur une \textbf{vie entière} : \emph{« une hirondelle ne fait pas le printemps »}, un jour ou un court moment ne rend pas heureux.

\begin{definition}[Le bonheur chez Aristote]
Le bonheur (\emph{eudaimonia}) est une \textbf{activité de l'âme conforme à la vertu parfaite}, exercée au long d'une \textbf{vie complète}, pourvue de \textbf{biens extérieurs} suffisants (santé, amis, ressources minimales).
\end{definition}

\courssub{La vertu (arêtê) : une disposition acquise par l'habitude}

Si le bonheur est activité vertueuse, qu'est-ce que la vertu (\emph{arêtê}) ? Ce n'est ni un sentiment (colère, peur), ni une faculté (capacité de ressentir), mais une \cle{disposition} (\emph{hexis}) : un état stable de l'âme qui nous fait choisir l'acte bon. Surtout, \textbf{la vertu ne vient pas de la nature} : « Rien de ce qui existe par nature ne peut être modifié par l'habitude » (\emph{Éthique à Nicomaque}, II, 1). On ne naît pas vertueux, on le \textbf{devient} par l'\cle{éducation} et l'\cle{exercice répété} (\emph{ethos}, d'où le mot « éthique »).

\begin{propriete}[Devenir vertueux par l'exercice]
\textbf{Thèse :} « Les vertus, nous les acquérons en les exerçant, comme on acquiert les arts. Car ce que nous devons faire après l'avoir appris, nous l'apprenons en le faisant : on devient bâtisseur en bâtissant, cithariste en jouant de la cithare ; de même, on devient juste en faisant des actes justes, tempérant en faisant des actes tempérants, courageux en faisant des actes courageux. » (Aristote, \emph{Éthique à Nicomaque}, II, 1).
\begin{itemize}
    \item \textbf{Acquisition :} par répétition d'actes conformes à la règle droite (\emph{orthos logos}).
    \item \textbf{Consolidation :} l'habitude (\emph{ethos}) transforme l'effort initial en disposition stable et finalement agréable.
    \item \textbf{Rôle de l'éducateur :} le législateur et le précepteur doivent habituer le jeune aux actes vertueux avant même qu'il n'en comprenne pleinement la raison.
\end{itemize}
\end{propriete}

\courssub{La doctrine du juste milieu (médiété) : la structure de la vertu}

C'est le cœur de l'éthique aristotélicienne. La vertu est un \textbf{juste milieu} (\emph{mesotês}) entre deux vices : l'un par \textbf{excès}, l'autre par \textbf{défaut}. Mais ce milieu n'est pas une moyenne arithmétique fixe (ce serait absurde de prescrire la même ration alimentaire à tous) ; il est \textbf{relatif à nous} (\emph{pros hêmas}), déterminé par la \cle{raison} et plus précisément par la \cle{prudence} (\emph{phronêsis}), intelligence pratique qui discerne, dans chaque situation concrète, ce qui convient.

\begin{definition}[La vertu comme juste milieu]
La vertu est \textbf{une disposition à choisir}, consistant en un \textbf{milieu relatif à nous}, \textbf{déterminé par la raison} (et comme le déterminerait l'homme prudent), \textbf{milieu entre deux vices}, l'un par excès, l'autre par défaut.
\end{definition}

\begin{attention}[Le juste milieu n'est pas la médiocrité !]
\begin{itemize}
    \item \textbf{Piège 1 :} confondre « milieu » et « moyenne arithmétique ». Le courage n'est pas « avoir 50\,\% de peur ». C'est avoir \emph{la peur qu'il faut, quand il faut, comme il faut, pour la raison qu'il faut}.
    \item \textbf{Piège 2 :} croire que le milieu est le même pour tous. Le repas du lutteur (excès pour le sédentaire) est son juste milieu.
    \item \textbf{Piège 3 :} penser que toute action admet un milieu. \textbf{Non :} l'adultère, le vol, le meurtre sont \emph{mauvais par nature} ; on ne peut pas les commettre « avec modération ».
\end{itemize}
\end{attention}

\begin{popfigure}
\begin{tikzpicture}[scale=1.1, every node/.style={font=\small}]
    % Ligne principale
    \draw[popline, thick] (-5.5,0) -- (5.5,0);
    % Graduations
    \foreach \x/\label in {-4/Vice par défaut, 0/Vertu (juste milieu), 4/Vice par excès} {
        \draw[popdark, thick] (\x,-0.15) -- (\x,0.15);
        \node[align=center, popdark] at (\x, -0.7) {\label};
    }
    % Exemple 1 : Courage (ligne du haut)
    \draw[popblue, very thick] (-4,1.2) -- (4,1.2);
    \node[popblue, font=\bfseries] at (-4,1.7) {Lâcheté};
    \node[popblue, font=\bfseries] at (0,1.7) {COURAGE};
    \node[popblue, font=\bfseries] at (4,1.7) {Témérité};
    \draw[popblue, ->] (-4,1.2) -- (-0.5,1.2);
    \draw[popblue, ->] (4,1.2) -- (0.5,1.2);
    % Exemple 2 : Générosité (ligne du bas)
    \draw[poporange, very thick] (-4,-1.2) -- (4,-1.2);
    \node[poporange, font=\bfseries] at (-4,-1.7) {Avarice};
    \node[poporange, font=\bfseries] at (0,-1.7) {GÉNÉROSITÉ};
    \node[poporange, font=\bfseries] at (4,-1.7) {Prodigalité};
    \draw[poporange, ->] (-4,-1.2) -- (-0.5,-1.2);
    \draw[poporange, ->] (4,-1.2) -- (0.5,-1.2);
    % Indication des pôles
    \node[popmuted] at (-5.5, 0.4) {Défaut};
    \node[popmuted] at (5.5, 0.4) {Excès};
\end{tikzpicture}
\legende{Schéma du juste milieu aristotélicien : la vertu (courage, générosité) est un point d'équilibre entre un vice par défaut et un vice par excès, déterminé par la prudence pour chaque situation concrète.}
\end{popfigure}

\begin{center}
\begin{tabularx}{\linewidth}{|>{\bfseries}l|X|X|X|}
\hline
\rowcolor{popblueL}
\textbf{Vice par défaut} & \textbf{Vertu (juste milieu)} & \textbf{Vice par excès} & \textbf{Domaine concerné} \\
\hline
Lâcheté & \cle{Courage} & Témérité & Peur et confiance face au danger \\
\hline
Avarice & \cle{Générosité} & Prodigalité & Don et dépense d'argent \\
\hline
Insensibilité & \cle{Tempérance} & Intempérance & Plaisirs corporels (manger, boire) \\
\hline
Petitesse d'âme & \cle{Magnanimité} & Vanité & Honneurs et estime de soi \\
\hline
\end{tabularx}
\end{center}

\courssub{Conditions du bonheur : vertu, prudence et biens extérieurs}

La vertu \emph{éthique} (courage, tempérance, justice) règle nos passions et nos actions. Mais pour choisir le juste milieu, il faut la \cle{vertu dianoétique} (intellectuelle) qu'est la \textbf{prudence} (\emph{phronêsis}) : ce n'est ni la science théorique (\emph{epistêmê}), ni l'habileté technique (\emph{technê}), mais la \textbf{sagacité du bien vivre}, la capacité de délibérer correctement sur ce qui est bon pour soi et pour la cité.

De plus, Aristote refuse tout idéalisme excessif : \textbf{la vertu seule ne suffit pas au bonheur}. Il faut un \textbf{minimum de biens extérieurs} : santé, amis, ressources matérielles suffisantes. « Car il est impossible, ou du moins difficile, de faire de belles actions si l'on n'a pas les instruments » (\emph{Éthique à Nicomaque}, I, 8). Le bonheur est donc \textbf{fragile} : il dépend en partie de la fortune, mais l'homme vertueux supporte les revers avec noblesse.

\begin{aretenir}[Les trois piliers du bonheur aristotélicien]
\begin{enumerate}
    \item \textbf{Vertu éthique} (caractère) : disposition à bien agir (juste milieu).
    \item \textbf{Prudence} (intelligence pratique) : discernement du juste milieu \emph{ici et maintenant}.
    \item \textbf{Biens extérieurs} (fortune) : conditions matérielles minimales pour agir.
\end{enumerate}
Le tout déployé sur une \textbf{vie complète}.
\end{aretenir}

\courssub{Application camerounaise : doser et discerner}

L'eudémonisme n'est pas une théorie abstraite ; il offre une grille de lecture concrète pour l'action quotidienne au Cameroun.

\begin{exemplebox}[L'élève de Première technique : travail et détente]
Un élève prépare le Probatoire.
\begin{itemize}
    \item \textbf{Excès (excès de zèle) :} il étudie 16 heures par jour, dort 4 heures, coupe tout lien social $\rightarrow$ épuisement, anxiété, échec probable.
    \item \textbf{Défaut (paresse) :} il ne révise pas, sort tout le temps, compte sur la chance $\rightarrow$ ignorance, regret, échec certain.
    \item \textbf{Juste milieu (persévérance réglée par la prudence) :} il planifie 3 heures de révision efficace le soir, 1 heure de sport ou de détente, 8 heures de sommeil. Il \emph{agit} avec régularité (habitude vertueuse) et \emph{discerne} (prudence) les priorités du programme. C'est la voie de la réussite scolaire.
\end{itemize}
\end{exemplebox}

\begin{exemplebox}[Le commerçant de Bafoussam : la générosité en acte]
\textbf{Situation :} M. Kamga, commerçant au Marché A de Bafoussam, réalise un bénéfice net de \num{500000} FCFA en fin d'année. Sa communauté (famille élargie, association de développement du village) sollicite son aide.
\begin{itemize}
    \item \textbf{Avarice (vice par défaut) :} il garde tout, « l'argent ne se partage pas ». Il s'isole, perd le respect, et la solidarité ne jouera pas pour lui en cas de coup dur.
    \item \textbf{Prodigalité (vice par excès) :} il donne tout pour « faire l'homme ». Sa famille manque de tout, son capital de travail est englouti $\rightarrow$ faillite ; il devient une charge pour les autres.
    \item \textbf{Générosité (vertu, juste milieu) :} il donne \textbf{10\,\%} (\num{50000} FCFA) : \num{30000} pour le fonds villageois (toiture de l'école), \num{20000} pour les frais de scolarité de neveux orphelins. Il garde \num{450000} FCFA pour réapprovisionner sa boutique, payer le loyer et nourrir sa famille.
\end{itemize}
\textbf{Analyse :} le montant (10\,\%) n'est pas une règle fixe. Si l'année est mauvaise (bénéfice de \num{100000} FCFA), donner \num{50000} serait prodigue ; si l'année est exceptionnelle (\num{5000000} FCFA), donner seulement \num{50000} serait avare. La \cle{prudence} calcule le « relatif à nous ».
\end{exemplebox}

\courssub{Problématisation critique : le bonheur personnel suffit-il à fonder l'obligation morale ?}

L'eudémonisme d'Aristote est une morale de \textbf{l'excellence} et de la \textbf{responsabilité}. Il évite le subjectivisme du plaisir (épicurisme) et le rigorisme abstrait du devoir (Kant, section suivante). Pourtant, une objection majeure se dresse :

\begin{attention}[La critique de l'égoïsme structurel]
Si le bonheur \emph{propre} est la fin suprême, alors \textbf{toute action vertueuse risque de devenir un moyen pour mon bonheur}.
\begin{itemize}
    \item Je suis juste \emph{afin d'être heureux} (car l'injustice trouble l'âme).
    \item Je suis courageux \emph{afin de réaliser mon excellence}.
    \item L'autre (l'ami, le citoyen, l'étranger) devient-il alors un \textbf{instrument} de ma félicité ?
\end{itemize}
Aristote répond que l'homme est un \emph{animal politique} : la vertu \emph{est} constitutive du bonheur, elle n'en est pas l'instrument extérieur. Aimer son ami pour lui-même (\emph{philia}) fait partie du bonheur. Mais la question reste ouverte : \textbf{le fondement de l'obligation envers autrui est-il assez fort si mon bonheur est l'horizon indépassable ?} C'est pour répondre à cette exigence d'universalité et de désintéressement que Kant posera le \textbf{devoir} comme unique mobile moral.
\end{attention}

\begin{savaistu}[Aristote au Cameroun : la sagesse de la mesure dans les chefferies]
La pensée du « juste milieu » résonne avec la sagesse traditionnelle des chefferies de l'Ouest et du Nord-Ouest camerounais. Le \emph{Fo} ou \emph{Fon} (chef) doit être \textbf{généreux} (redistribuer les récoltes, financer les fêtes) sans être \textbf{prodigue} (vider les greniers royaux), et sans être \textbf{avare} (thésauriser au détriment du peuple). Sa légitimité repose sur cette \emph{phronêsis} politique : discerner, chaque saison, la part juste. L'eudémonisme n'est pas une pensée « occidentale » ; il formalise une sagesse universelle de la \textbf{mesure}.
\end{savaistu}

\begin{exoresolu}[Analyser une situation à la lumière du juste milieu]
\textbf{Énoncé :} Paul, chef d'équipe sur un chantier à Douala, découvre que son meilleur ami et collègue, Jacques, a volé des sacs de ciment. Paul hésite : le dénoncer (rigueur) ou se taire (amitié).
\textbf{Question :} en vous référant à la doctrine du juste milieu, expliquez quelle serait l'attitude vertueuse de Paul et pourquoi ce n'est ni de la lâcheté, ni de la délation excessive.

\tcblower
\textbf{Solution détaillée :}
\begin{enumerate}
    \item \textbf{Identifier la vertu en jeu :} il s'agit de la \textbf{justice} (respect de la loi et du bien commun) tempérée par l'\textbf{amitié} et la \textbf{bienveillance}. La vertu visée est l'\textbf{équité} (\emph{epieikeia}) : corriger la loi là où elle est trop générale pour le cas particulier.
    \item \textbf{Identifier les deux excès :}
    \begin{itemize}
        \item \textbf{Vice par défaut (lâcheté, complicité) :} se taire par peur de perdre l'ami ou par « solidarité » mal comprise. Paul devient complice, trahit sa fonction de chef, encourage l'impunité $\rightarrow$ injustice pour l'entreprise et les autres ouvriers.
        \item \textbf{Vice par excès (rigueur excessive, délation froide) :} dénoncer immédiatement Jacques à la direction sans lui parler, sans chercher à comprendre (dettes, maladie dans la famille), en appliquant le règlement à la lettre sans humanité. C'est une « justice sans miséricorde » qui détruit le lien social et la confiance.
    \end{itemize}
    \item \textbf{Déterminer le juste milieu (relatif à nous, par la prudence) :}
    \begin{itemize}
        \item Paul convoque Jacques \textbf{en privé} (respect de la dignité).
        \item Il \textbf{confronte} les faits : « J'ai vu, on m'a signalé. » (vérité).
        \item Il \textbf{écoute} la raison (prudence : quel est le contexte ?).
        \item Il \textbf{exige la réparation} : restitution ou retenue sur salaire, engagement de ne plus recommencer.
        \item \textbf{Sanction proportionnée :} avertissement formel, mutation d'équipe si nécessaire, mais pas licenciement immédiat ni poursuites (sauf récidive).
    \end{itemize}
    \item \textbf{Conclusion :} l'acte vertueux n'est ni la tolérance passive (vice par défaut), ni la répression aveugle (vice par excès). C'est une \textbf{décision raisonnée, contextuelle, qui vise la correction du fautif et la préservation du bien commun}. Paul agit en homme prudent (\emph{phronimos}), non en robot ni en complice.
\end{enumerate}
\end{exoresolu}

\begin{exercice}[Entraînement : repérer le vice et la vertu]
Pour chacune des situations ci-dessous, identifiez : (a) la vertu concernée, (b) le vice par défaut, (c) le vice par excès, (d) l'acte du « juste milieu » probable.
\begin{enumerate}
    \item Un député vote une loi impopulaire mais nécessaire pour l'avenir du pays, malgré les menaces sur sa réélection.
    \item Une mère donne tout son salaire à son fils majeur qui refuse de travailler, au détriment de ses propres soins médicaux.
    \item Un ingénieur refuse de signer un rapport de conformité falsifié, mais propose une solution technique alternative pour tenir le délai.
    \item Un étudiant refuse tout loisir pendant 3 mois avant l'examen, dort 3 heures par nuit et fait une dépression.
\end{enumerate}
\end{exercice}

\begin{corrige}[Exercice n°1]
\begin{enumerate}
    \item \textbf{Vertu :} courage civique, intégrité. \textbf{Défaut :} lâcheté (voter ce qui est populaire pour se faire réélire). \textbf{Excès :} témérité, intransigeance stérile (refuser tout compromis, faire tomber le gouvernement pour un principe mineur). \textbf{Juste milieu :} voter la loi, mais l'expliquer pédagogiquement aux électeurs et accepter le risque électoral.
    \item \textbf{Vertu :} générosité, responsabilité parentale. \textbf{Défaut :} dureté, abandon (couper les vivres brutalement). \textbf{Excès :} prodigalité, complicité infantilisante (tout donner). \textbf{Juste milieu :} fixer un ultimatum (3 mois), financer une formation, loger temporairement, mais exiger une contribution ou un loyer symbolique.
    \item \textbf{Vertu :} honnêteté professionnelle, justice. \textbf{Défaut :} complicité, corruption passive (signer le faux rapport). \textbf{Excès :} dénonciation publique immédiate sans chercher de solution (rigueur stérile). \textbf{Juste milieu :} refuser de signer (intégrité), alerter la hiérarchie par écrit (protection), proposer l'alternative technique (attitude constructive).
    \item \textbf{Vertu :} tempérance, étude équilibrée. \textbf{Défaut :} paresse, négligence (aucune révision). \textbf{Excès :} excès ascétique, surmenage autodestructeur. \textbf{Juste milieu :} planning réaliste (par exemple 4 heures par jour), pauses sportives, 7 heures de sommeil, vie sociale légère. La prudence dit : « Mieux vaut 4 heures efficaces que 16 heures en état d'épuisement. »
\end{enumerate}
\end{corrige}

\coursec{L'utilitarisme : la morale de l'utilité}

\courssub{De l'eudémonisme à l'utilitarisme : du bonheur individuel au bonheur collectif}

Dans la section précédente, nous avons vu que l'eudémonisme fait du \emph{bonheur} la fin suprême de l'action humaine : pour Aristote, vivre bien, c'est accomplir son excellence propre et atteindre ainsi la félicité. Mais cette morale reste centrée sur l'individu : c'est \emph{mon} bonheur, \emph{ma} vie accomplie qui sont en jeu. Une question surgit alors, d'une redoutable actualité : que faire lorsque les bonheurs individuels entrent en conflit ? Le bonheur du commerçant qui vend des produits frelatés s'oppose au bonheur des consommateurs empoisonnés ; le bonheur d'un village s'oppose parfois à celui de la nation entière. Faut-il une morale capable d'\emph{additionner} les bonheurs, de les comparer, de trancher ?

C'est précisément la réponse de l'\cle{utilitarisme}, né en Angleterre aux XVIII\textsuperscript{e} et XIX\textsuperscript{e} siècles, dans un pays en pleine révolution industrielle, confronté à la misère ouvrière, à l'insalubrité des villes et à la nécessité de réformer les lois, les prisons, l'école. Deux penseurs en sont les fondateurs : \textbf{Jeremy Bentham} (1748-1832), juriste et réformateur, et son disciple \textbf{John Stuart Mill} (1806-1873), philosophe et économiste. Tous deux veulent faire de la morale une \emph{science} : non plus une affaire de sentiments vagues ou de commandements reçus, mais un calcul rationnel au service de la société.

\begin{definition}[L'utilitarisme]
L'\cle{utilitarisme} est une doctrine morale qui juge les actions et les institutions d'après leur \emph{utilité}, c'est-à-dire leur capacité à produire \og le plus grand bonheur du plus grand nombre \fg{}. Une action est moralement bonne si elle augmente la somme totale de bonheur dans la société ; elle est mauvaise si elle la diminue.
\end{definition}

\courssub{Le principe d'utilité}

L'intuition de départ est d'une simplicité frappante. Pourquoi désapprouvons-nous le vol ? Parce qu'il cause plus de souffrance (à la victime, à la société qui se méfie) que de plaisir (au voleur). Pourquoi admirons-nous le médecin qui se dévoue pendant une épidémie ? Parce que son action produit un bonheur immense pour un grand nombre. Notre jugement moral spontané est déjà, souvent, un calcul des conséquences. Bentham se contente de le formuler rigoureusement.

\begin{definition}[Le principe d'utilité]
Le \cle{principe d'utilité} est le critère selon lequel une action est \emph{bonne} si elle augmente le bonheur global et \emph{mauvaise} si elle le diminue. Bentham l'appelle aussi \og principe du plus grand bonheur \fg{}. Il s'applique aussi bien aux actions privées qu'aux lois et aux institutions politiques.
\end{definition}

Ce principe repose sur une constatation que Bentham expose dès l'ouverture de ses \emph{Principes de la morale et de la législation} (1789) : \og la nature a placé l'humanité sous le gouvernement de deux maîtres souverains, la peine et le plaisir \fg{}. Tout ce que nous faisons, tout ce que nous devons faire, se règle sur eux. La morale n'a donc pas à inventer des valeurs mystérieuses : elle doit \emph{mesurer} les plaisirs et les peines que produisent nos actions.

\begin{propriete}[Impartialité du calcul utilitariste]
Dans le calcul utilitariste, \og chacun compte pour un, et personne ne compte pour plus d'un \fg{} (formule de Bentham reprise par Mill). Le bonheur du ministre ne pèse pas plus lourd que celui du paysan de l'Extrême-Nord ; le mien ne pèse pas plus lourd que celui d'un inconnu. L'utilitarisme dépasse donc radicalement l'\emph{égoïsme} : ce n'est pas \emph{mon} utilité qui compte, mais l'utilité de \emph{tous} ceux que mon action affecte.
\end{propriete}

\courssub{Bentham : le calcul des plaisirs ou \og arithmétique morale \fg{}}

Pour que le principe d'utilité devienne opératoire, il faut pouvoir \emph{comparer} les plaisirs. Bentham propose donc une véritable \cle{arithmétique morale} : toute action doit être évaluée en pesant les plaisirs qu'elle promet contre les peines qu'elle risque de causer, selon sept critères.

\begin{aretenir}[Les sept critères du calcul benthamien]
Pour évaluer un plaisir (ou une peine), Bentham retient :
\begin{enumerate}
\item l'\textbf{intensité} : quelle est sa force ?
\item la \textbf{durée} : combien de temps dure-t-il ?
\item la \textbf{certitude} : quelle est la probabilité qu'il se produise ?
\item la \textbf{proximité} : dans combien de temps se produira-t-il ?
\item la \textbf{fécondité} : a-t-il des chances d'être suivi d'autres plaisirs ?
\item la \textbf{pureté} : a-t-il des chances de n'être pas suivi de peines ?
\item l'\textbf{étendue} : combien de personnes sont concernées ?
\end{enumerate}
Le dernier critère, l'étendue, est décisif : c'est lui qui fait passer la morale de l'individu à la société.
\end{aretenir}

\begin{popfigure}
\begin{center}
\begin{tikzpicture}[scale=1.0,>=Stealth]
% Socle et colonne
\draw[fill=popmuted!40] (-0.5,0) rectangle (0.5,0.25);
\draw[fill=popmuted!40] (-0.12,0.25) rectangle (0.12,2.6);
% Flèche centrale
\draw[fill=popdark] (-0.35,2.6) -- (0.35,2.6) -- (0,2.25) -- cycle;
% Bras de la balance
\draw[line width=2pt,popdark] (-3.2,2.6) -- (3.2,2.6);
% Chaînes et plateau gauche (plaisirs)
\draw[popblue] (-3.0,2.6) -- (-3.6,1.5);
\draw[popblue] (-3.0,2.6) -- (-2.4,1.5);
\draw[fill=popblueL,draw=popblue,line width=1pt] (-3.9,1.5) arc (180:360:0.9 and 0.35) -- (-2.1,1.5);
\node[popblue,font=\bfseries] at (-3.0,0.75) {PLAISIRS};
\node[popblue,font=\small,align=center] at (-3.0,0.25) {intensité, durée,\\fécondité, étendue\ldots};
% Chaînes et plateau droit (peines)
\draw[poporange] (3.0,2.6) -- (2.4,1.5);
\draw[poporange] (3.0,2.6) -- (3.6,1.5);
\draw[fill=poporangeL,draw=poporange,line width=1pt] (2.1,1.5) arc (180:360:0.9 and 0.35) -- (3.9,1.5);
\node[poporange,font=\bfseries] at (3.0,0.75) {PEINES};
\node[poporange,font=\small,align=center] at (3.0,0.25) {souffrances causées,\\risques, contraintes\ldots};
% Verdict
\node[align=center,font=\small\itshape] at (0,3.3) {Si les plaisirs l'emportent : l'action est bonne};
\end{tikzpicture}
\end{center}
\legende{La balance de l'arithmétique morale de Bentham : toute action est jugée en pesant la somme des plaisirs qu'elle produit contre la somme des peines qu'elle engendre, pour l'ensemble des personnes concernées.}
\end{popfigure}

\courssub{Mill : la distinction qualitative des plaisirs}

Le calcul de Bentham a un point faible, que ses adversaires ne manquent pas de souligner : si seule la \emph{quantité} de plaisir compte, alors le plaisir grossier de l'ivrogne vaut autant que celui du poète, pourvu qu'il soit aussi intense. La morale deviendrait, disait-on, \og une doctrine digne des pourceaux \fg{}.

John Stuart Mill répond à cette objection dans son ouvrage \emph{L'Utilitarisme} (1861) en introduisant une distinction \emph{qualitative} : les plaisirs ne se valent pas seulement par leur quantité, mais aussi par leur \emph{nature}. Les plaisirs de l'esprit — l'intelligence, la culture, l'amitié, la dignité — sont \emph{supérieurs} aux plaisirs du corps, et celui qui a connu les deux préfère toujours les premiers. D'où sa formule célèbre :

\begin{aretenir}[La hiérarchie des plaisirs selon Mill]
\og Mieux vaut un Socrate mécontent qu'un imbécile satisfait. \fg{} Autrement dit : un plaisir noble, même accompagné d'insatisfaction, vaut mieux qu'un contentement bestial. L'imbécile est satisfait parce qu'il ne connaît que son petit plaisir ; Socrate, qui connaît les deux, ne renoncerait à la vie de l'esprit pour rien au monde.
\end{aretenir}

Cette correction est capitale : elle montre que l'utilitarisme n'est pas une invitation à la jouissance facile, mais qu'il peut exiger l'effort, l'éducation, le sacrifice — car ce sont eux qui produisent les bonheurs les plus hauts et les plus durables.

\courssub{Une morale sociale et politique}

La grande originalité de l'utilitarisme, comparée à l'épicurisme et à l'eudémonisme, est sa portée \emph{politique}. Si le critère du bien est le plus grand bonheur du plus grand nombre, alors ce sont d'abord les \emph{lois} et les \emph{institutions} qui doivent être jugées : une loi est bonne si elle maximise le bonheur collectif, mauvaise si elle le diminue. Bentham a ainsi milité pour la réforme des prisons, l'éducation pour tous, la simplification du droit. L'utilitarisme devient un instrument de réforme sociale : on ne demande plus \og cette loi est-elle traditionnelle ? est-elle voulue par Dieu ? \fg{}, mais \og cette loi rend-elle les citoyens plus heureux ? \fg{}.

\begin{exemplebox}[Le calcul utilitariste appliqué à une décision publique]
Un gouvernement hésite entre deux projets de même coût : construire un hôpital régional ou un grand stade. Le calcul utilitariste compare le \emph{bonheur total} produit par chaque option, en tenant compte de l'intensité, de la durée et de l'étendue (nombre de bénéficiaires).
\end{exemplebox}

\begin{popfigure}
\begin{center}
\begin{tikzpicture}
\begin{axis}[
  ybar, bar width=1.1cm,
  width=0.85\linewidth, height=6.2cm,
  symbolic x coords={Hôpital, Stade},
  xtick=data,
  ymin=0, ymax=100,
  ylabel={Bonheur total estimé (unités d'utilité)},
  ylabel style={font=\small},
  nodes near coords,
  every node near coord/.append style={font=\small\bfseries},
  axis line style={popdark},
  tick label style={font=\small},
  enlarge x limits=0.45,
]
\addplot[fill=popblue,draw=popblue!70!black] coordinates {(Hôpital,86) (Stade,41)};
\end{axis}
\end{tikzpicture}
\end{center}
\legende{Illustration du calcul utilitariste : l'hôpital procure un bonheur intense (santé, vies sauvées), durable, certain et étendu (toute une région, toute l'année) ; le stade procure un plaisir réel mais intermittent et moins vital. Le principe d'utilité commande donc de construire l'hôpital (valeurs fictives à visée pédagogique).}
\end{popfigure}

\begin{exoresolu}[Vaccination obligatoire : un raisonnement utilitariste]
Énoncé. Une campagne de vaccination obligatoire contre une épidémie coûte 2 milliards de FCFA à l'État et contraint quelques milliers de récalcitrants qui ne veulent pas se faire vacciner. Elle épargne environ 50\,000 vies. Un utilitariste jugera-t-il cette politique bonne ? Justifier.
\tcblower
Solution détaillée. Appliquons les critères de Bentham.
\begin{enumerate}
\item \emph{Étendue} : des millions de personnes sont protégées, directement (les vaccinés) et indirectement (l'épidémie ne se propage plus).
\item \emph{Intensité} : le bien évité est la mort — le pire des maux ; le mal imposé est une contrainte légère (une piqûre, une obligation).
\item \emph{Fécondité} : la campagne produit des biens durables — économie préservée, hôpitaux désengorgés, enfants non orphelins.
\item \emph{Peines} : le coût financier (2 milliards de FCFA) et la frustration des récalcitrants sont réels mais faibles en regard des 50\,000 vies sauvées.
\end{enumerate}
Conclusion : le bonheur total augmente massivement. L'utilitariste juge donc la vaccination obligatoire \emph{moralement bonne}, même si elle contraint une minorité — car chacun compte pour un, et 50\,000 vies pèsent infiniment plus lourd que quelques milliers de mécontentements.
\end{exoresolu}

\begin{exemplebox}[Autres applications camerounaises]
\begin{itemize}
\item \emph{Le barrage hydroélectrique} : la construction d'un barrage impose le déplacement d'un village de quelques centaines de personnes, mais alimente en électricité des millions de foyers, usines et hôpitaux. L'utilitariste peut juger le projet bon \emph{à condition} que les déplacés soient indemnisés et relogés — sinon leur peine, ignorée, fausserait le calcul.
\item \emph{La fermeture d'un marché insalubre} : fermer un marché où se vendent des denrées avariées prive des centaines de commerçants de leur revenu, mais protège des dizaines de milliers de consommateurs contre les intoxications et les épidémies. Le bonheur du plus grand nombre l'emporte — mais un bon utilitariste exigera un marché de remplacement.
\end{itemize}
\end{exemplebox}

\begin{attention}[Erreur fréquente : l'utilitarisme ne justifie pas n'importe quoi]
On croit souvent, à tort, que l'utilitarisme autorise tout ce qui est \og efficace \fg{} ou tout ce qui arrange le pouvoir. C'est faux à deux titres. D'une part, le calcul doit être \emph{impartial} : il porte sur le bonheur de \textbf{tous} ceux qui sont affectés, pas sur celui d'un seul groupe (le parti au pouvoir, l'ethnie majoritaire, les riches). D'autre part, il doit être \emph{complet} : oublier la souffrance d'une minorité dans le calcul, c'est cesser d'être utilitariste pour devenir simplement tyrannique. Un raisonnement utilitariste malhonnête n'est pas un raisonnement utilitariste.
\end{attention}

\begin{savaistu}[Le corps \og utile \fg{} de Bentham]
Fidèle à sa doctrine jusqu'au bout, Jeremy Bentham a demandé à sa mort (1832) que son corps soit disséqué publiquement — pour faire progresser la science médicale — puis embaumé et exposé. Son \og auto-icône \fg{}, vêtu de ses propres habits, est encore visible aujourd'hui dans une vitrine de l'University College de Londres. Ultime geste utilitariste : même mort, Bentham voulait encore être \emph{utile} à l'instruction du plus grand nombre.
\end{savaistu}

\courssub{Critiques de l'utilitarisme}

L'utilitarisme est séduisant par sa rationalité et sa générosité, mais il se heurte à trois objections majeures, qu'il faut savoir mobiliser dans une dissertation.

\begin{enumerate}
\item \textbf{La difficulté de mesurer le bonheur.} Comment additionner des plaisirs ? Combien d'unités de bonheur vaut une vie sauvée, une œuvre d'art, une soirée en famille ? Le calcul benthamien suppose une monnaie commune du bonheur qui n'existe peut-être pas. Mill lui-même, en distinguant les qualités de plaisirs, a reconnu que la pure quantité ne suffit pas.
\item \textbf{Le risque de sacrifier les minorités.} Si seul compte le bonheur du \emph{plus grand nombre}, ne peut-on pas justifier l'oppression d'une petite minorité, dès lors que la majorité en retire un avantage ? Faut-il expulser un village pour un barrage, exproprier des paysans pour une autoroute, au simple motif que des millions en profiteront ? L'utilitarisme semble manquer d'une protection absolue des droits individuels.
\item \textbf{La fin justifie-t-elle les moyens ?} Si une injustice ponctuelle (condamner un innocent pour apaiser une foule, mentir \og pour le bien du pays \fg{}) produisait un bonheur total supérieur, l'utilitarisme devrait l'approuver. Or notre conscience morale proteste : certaines choses — tuer un innocent, torturer — semblent mauvaises \emph{en elles-mêmes}, quelles qu'en soient les conséquences. C'est précisément l'objection que la morale du devoir, celle de Kant, développera : la valeur d'une action ne réside pas dans ses résultats, mais dans l'intention et le respect du devoir.
\end{enumerate}

\begin{exercice}[Pour s'entraîner]
La mairie d'une ville camerounaise veut raser un quartier précaire de 500 familles pour construire un centre des affaires qui, affirme-t-elle, créera 20\,000 emplois. 1) Construisez le raisonnement utilitariste en faveur du projet. 2) Montrez, à l'aide des critiques de l'utilitarisme, ce que ce raisonnement risque d'oublier. 3) Formulez votre conclusion personnelle argumentée.
\end{exercice}

\begin{corrige}[Exercice]
1) \emph{Raisonnement utilitariste} : le centre des affaires procure un bonheur de grande \emph{étendue} (20\,000 salariés et leurs familles, soit peut-être 100\,000 personnes), de grande \emph{fécondité} (revenus, écoles, santé financés par ces emplois) et de longue \emph{durée}. La peine des 500 familles déplacées est réelle mais quantitativement inférieure. Le bonheur total semble donc augmenter : le projet paraît bon.

2) \emph{Critiques} : d'abord, le calcul est-il honnête ? Le bonheur promis est \emph{incertain} (critère de certitude) tandis que la misère des expulsés est immédiate et sûre. Ensuite, le risque de sacrifier une minorité pauvre au profit d'intérêts puissants est manifeste : les 500 familles perdent leur toit, leur quartier, leurs liens — des biens que l'argent ne mesure pas (difficulté de mesurer le bonheur). Enfin, exproprier sans consentement ni relogement décent, c'est traiter des personnes comme de simples moyens : la fin économique ne justifie pas tous les moyens.

3) \emph{Conclusion possible} : un utilitarisme rigoureux et impartial n'approuverait le projet qu'à des conditions strictes — relogement digne, indemnisation juste, emplois réellement accessibles aux déplacés. À défaut, le calcul est truqué au profit d'un groupe, et cesse d'être moral. On mesure ainsi la force de l'utilitarisme (il oblige à compter tout le monde) et sa limite (il ne garantit à lui seul ni la justice ni la dignité de chacun) — ce qui annonce la nécessité d'une morale du devoir.
\end{corrige}

\begin{aretenir}[L'essentiel de la section]
\begin{itemize}
\item L'utilitarisme (Bentham, Mill) juge les actions et les lois d'après leur \emph{utilité} : est bonne celle qui produit \og le plus grand bonheur du plus grand nombre \fg{}.
\item Bentham propose un \emph{calcul des plaisirs} selon sept critères (intensité, durée, certitude, proximité, fécondité, pureté, étendue) ; Mill ajoute une hiérarchie \emph{qualitative} : les plaisirs de l'esprit valent mieux que ceux du corps.
\item Le calcul est \emph{impartial} : chacun compte pour un, personne ne compte pour plus d'un — l'utilitarisme n'est donc pas un égoïsme.
\item Appliqué aux politiques publiques (vaccination, barrages, salubrité), il est un puissant instrument de décision, mais il bute sur trois critiques : l'impossibilité de mesurer le bonheur, le danger de sacrifier les minorités, et la tentation de justifier de mauvais moyens par une bonne fin.
\end{itemize}
\end{aretenir}

\coursec{La morale du devoir : Kant}

Après avoir examiné les morales de l'intérêt — où la valeur d'une action se mesure à son utilité ou au plaisir qu'elle procure — nous abordons une conception radicalement différente : la \textbf{morale du devoir}. Ici, ce qui compte, ce n'est pas la conséquence de l'acte, mais l'intention qui le guide. Emmanuel Kant (1724–1804), professeur à l'université de Königsberg, en est le représentant majeur. Son projet : fonder la morale sur la raison pure, indépendante de tout désir, de tout intérêt, de toute circonstance empirique.

\courssub{Emmanuel Kant et le projet d'une morale rationnelle}

Né en 1724 à Königsberg (aujourd'hui Kaliningrad, Russie), Kant y vécut toute sa vie, menant une existence d'une régularité proverbiale. Ses ouvrages fondateurs pour l'éthique sont le \emph{Fondement de la métaphysique des mœurs} (1785) et la \emph{Critique de la raison pratique} (1788). Il y pose une question décisive : \og Quel est le principe suprême de la morale ? \fg{} Sa réponse exclut toute référence au bonheur, au plaisir ou à l'utilité : la morale ne peut reposer sur des mobiles empiriques, variables et contingents. Elle doit trouver son fondement dans la \cle{raison pure pratique}, c'est-à-dire dans la capacité de l'homme à se donner à lui-même sa loi.

\begin{savaistu}[La régularité légendaire de Kant]
Kant n'a jamais quitté sa ville natale de Königsberg. Ses promenades quotidiennes étaient si ponctuelles que les habitants réglaient leurs montres sur son passage. Cette discipline extérieure reflète sa conception intérieure de la morale : l'\cle{autonomie}, c'est-à-dire l'obéissance à une loi qu'on se donne à soi-même par la raison, non par contrainte extérieure ni par inclination.
\end{savaistu}

\courssub{La bonne volonté : unique bien sans restriction}

Kant ouvre le \emph{Fondement} par une affirmation célèbre :

\begin{quote}
\og Il n'y a rien au monde, ni même en général hors du monde, qui puisse être tenu sans restriction pour bon, sinon seulement une \cle{bonne volonté}. \fg
\end{quote}

L'intelligence, le courage, la richesse, la santé, le bonheur même — tous ces biens peuvent devenir mauvais si la volonté qui les utilise n'est pas bonne. Seule la \cle{bonne volonté} est bonne \emph{en soi}, sans condition. Elle ne l'est pas par ce qu'elle \emph{réalise} (ses effets), mais par ce qu'elle \emph{est} : une volonté déterminée par le respect de la loi morale.

\begin{definition}[Devoir (sens kantien)]
Nécessité d'agir par \textbf{respect pour la loi morale}, indépendamment de toute inclination, de tout intérêt ou de toute conséquence attendue. Le devoir est la forme même de l'action morale : agir \emph{par devoir}, non simplement \emph{conformément au devoir}.
\end{definition}

\courssub{Agir par devoir vs agir conformément au devoir}

C'est la distinction capitale de l'éthique kantienne. Deux actions peuvent être \emph{extérieurement identiques} (même geste, même résultat) et pourtant avoir une valeur morale radicalement différente selon le \textbf{mobile} qui les anime.

\begin{itemize}
\item \textbf{Agir conformément au devoir} : l'action respecte la règle morale, mais son mobile est un intérêt, une inclination, une crainte (ex. : ne pas voler par peur de la prison, rendre la monnaie pour garder sa clientèle).
\item \textbf{Agir par devoir} : l'action est accomplie \emph{uniquement} parce qu'elle est commandée par la loi morale. Le mobile est le respect pur de la loi. L'agent se dit : \og Je fais cela \emph{parce que c'est mon devoir}, et non \og Je fais cela \emph{afin que}... \fg
\end{itemize}

\begin{exemplebox}[Le marchand honnête — exemple kantien]
Un marchand ne trompe pas ses clients : il rend la monnaie exacte, ne falsifie pas ses poids. \textbf{Mais} il agit ainsi \emph{par calcul} : s'il trompait, il perdrait sa réputation et sa clientèle. Son action est \textbf{conforme au devoir}, mais elle n'a \textbf{aucune valeur morale propre} pour Kant, car son mobile est l'intérêt personnel (conserver ses clients), non le respect de la loi morale.
\end{exemplebox}

\begin{popfigure}
\begin{tikzpicture}[
  node distance=1.2cm and 1.5cm,
  box/.style={draw=popdark, thick, rounded corners, fill=popcream, minimum width=3.2cm, minimum height=1.1cm, align=center, font=\small},
  arrow/.style={-Stealth, thick, draw=popblue},
  arrowalt/.style={-Stealth, thick, draw=poporange},
  labelnode/.style={font=\small, text width=3.5cm, align=center}
]
% Nœuds sources
\node[box, fill=poporangeL, align=center] (interet) at (0,0){Intérêt personnel\\ (garder la clientèle)};
\node[box, fill=popblueL, align=center] (respect) at (0,-3){Respect de la loi morale\\ (honnêteté par devoir)};
% Nœud action commune
\node[box, fill=popgreenL, minimum width=4cm] (action) at (5.5,-1.5) {Action : \textbf{rendre la monnaie exacte}};
% Flèches
\draw[arrow] (interet.east) -- node[above, labelnode, align=center]{Agir \textbf{conformément}\\ au devoir} (action.west);
\draw[arrowalt] (respect.east) -- node[below, labelnode] {Agir \textbf{par} devoir} (action.west);
% Étiquettes de valeur morale
\node[labelnode, text=poporange, font=\small\bfseries] at (5.5,1.2) {Valeur morale : \textcolor{poporange}{nulle}};
\node[labelnode, text=popblue, font=\small\bfseries] at (5.5,-4.2) {Valeur morale : \textcolor{popblue}{authentique}};
\end{tikzpicture}
\legende{Schéma comparatif : une même action extérieure (rendre la monnaie) peut procéder de deux mobiles radicalement différents. Seule l'action motivée par le respect de la loi morale (par devoir) possède une valeur morale véritable chez Kant.}
\end{popfigure}

\begin{attention}[Piège fréquent : confondre conformité et moralité]
\emph{Erreur} : \og Si je fais ce qui est juste, mon action est morale. \fg
\emph{Correction} : Pour Kant, la moralité ne se juge pas à l'apparence (le \textbf{fait}), mais à l'intention (le \textbf{mobile}). Un élève qui rend un portefeuille trouvé à la gare de Ngaoundéré \emph{en espérant une récompense} agit conformément au devoir, mais non par devoir. Son acte n'a pas de valeur morale au sens strict.
\end{attention}

\begin{exemplebox}[Deux élèves à la gare de Ngaoundéré]
Deux élèves trouvent un portefeuille à la gare de Ngaoundéré.
\begin{itemize}
\item \textbf{Élève A} : le rend au guichet \emph{en se disant} : \og Le propriétaire me donnera peut-être de l'argent. \fg{} Mobile : intérêt / espoir de récompense. Action \textbf{conforme au devoir}, sans valeur morale.
\item \textbf{Élève B} : le rend \emph{en se disant} : \og C'est mon devoir de le rendre, qu'il y ait récompense ou non. \fg{} Mobile : respect de la loi morale. Action \textbf{par devoir} — seule celle-ci est moralement bonne pour Kant.
\end{itemize}
\end{exemplebox}

\courssub{L'impératif catégorique : la loi morale universelle}

Comment connaître cette loi morale ? Kant répond : par la \cle{raison pure}. Elle nous commande sous la forme d'un \textbf{impératif} (un \og il faut \fg). Mais tous les impératifs ne se valent pas.

\begin{center}
\begin{tabularx}{\linewidth}{|l|X|X|}
\hline
\rowcolor{popblueL}
\textbf{Critère} & \textbf{Impératif hypothétique} & \textbf{Impératif catégorique} \\
\hline
Forme & \og Si tu veux $X$, fais $Y$ \fg & \og Fais $Y$ \fg (sans condition) \\
\hline
Condition & Dépend d'un \textbf{but} désiré (bonheur, réussite, plaisir...) & \textbf{Inconditionnel} : vaut pour tout être rationnel, quels que soient ses désirs \\
\hline
Exemple & \og Si tu veux réussir ton Probatoire, révise tes cours. \fg & \og Ne mens pas. \fg (valable même si mentir t'avantage) \\
\hline
Fondement & Empirique (basé sur l'expérience, les inclinations) & \textbf{A priori} (basé sur la raison pure) \\
\hline
Portée & Technique / pragmatique & \textbf{Moral} \\
\hline
\end{tabularx}
\end{center}

\begin{definition}[Impératif catégorique]
Commandement \textbf{inconditionnel} de la raison pure qui exige d'agir selon une \cle{maxime} (règle subjective d'action) que l'on peut vouloir devenir \textbf{loi universelle}, et qui traite l'humanité toujours comme une \textbf{fin}, jamais comme un simple \textbf{moyen}.
\end{definition}

Kant donne plusieurs formulations de l'impératif catégorique ; les deux principales sont :

\begin{enumerate}
\item \textbf{Formule de la loi universelle} : \og Agis seulement d'après la maxime telle que tu puisses vouloir en même temps qu'elle devienne une \textbf{loi universelle}. \fg
\item \textbf{Formule de l'humanité comme fin} : \og Agis de telle sorte que tu traites l'humanité, en ta personne et en la personne de tout autre, toujours en même temps comme une \textbf{fin}, et \textbf{jamais simplement comme un moyen}. \fg
\end{enumerate}

\courssub{Tester sa maxime : la méthode de l'universalisation}

Une \cle{maxime} est la règle subjective sur laquelle j'agis (ex. : \og Je mentirai pour sortir d'embarras \fg). Pour savoir si elle est morale, Kant propose un test en trois temps.

\begin{methode}[Tester une maxime à la manière de Kant]
\begin{enumerate}
\item \textbf{Formuler ma maxime} explicitement : \og Je ferai $X$ dans la circonstance $C$ dans le but $B$. \fg
\item \textbf{L'universaliser} : \og Et si \emph{tout le monde} agissait sur cette maxime ? \fg (Imaginer un monde où cette maxime est loi universelle.)
\item \textbf{Vérifier la contradiction} : 
\begin{itemize}
\item \textbf{Contradiction dans la conception} : la maxime, universalisée, se détruit elle-même (ex. : le mensonge universel rend la communication impossible, donc le mensonge lui-même inefficace).
\item \textbf{Contradiction dans la volonté} : on ne peut pas \emph{vouloir} rationnellement un tel monde (ex. : refuser toute aide aux autres : on pourrait en avoir besoin soi-même un jour).
\end{itemize}
\item \textbf{Vérifier le respect des personnes} : cette maxime traite-t-elle autrui comme une fin ou comme un simple moyen ?
\end{enumerate}
Si la maxime passe le test : elle est un \textbf{devoir parfait} (ex. : ne pas mentir, ne pas voler). Si elle échoue seulement par contradiction dans la volonté : c'est un \textbf{devoir imparfait} (ex. : se perfectionner, secourir autrui — on ne peut pas le faire en tout temps, mais on ne doit pas s'en dispenser).
\end{methode}

\begin{exoresolu}[Appliquer le test : « Je ferai une fausse promesse pour obtenir un prêt »]
\textbf{Énoncé} : Un étudiant a besoin d'argent pour payer ses frais d'inscription. Il songe à promettre de rembourser dans un mois, alors qu'il sait qu'il ne le pourra pas. Tester cette maxime.
\\
\tcblower
\textbf{Solution détaillée} :
\begin{enumerate}
\item \textbf{Maxime} : \og Je ferai une fausse promesse d'emprunt pour obtenir de l'argent quand j'en ai besoin. \fg
\item \textbf{Universalisation} : Imaginons un monde où \emph{tout le monde} fait de fausses promesses d'emprunt dès qu'il a besoin d'argent.
\item \textbf{Contradiction dans la conception} : Dans un tel monde, \textbf{personne ne croirait plus aux promesses d'emprunt}. L'institution même de la promesse s'effondrerait. Or, pour que ma fausse promesse fonctionne, il faut que l'autre \emph{croie} à ma promesse. Universaliser ma maxime détruit la condition de sa propre réussite. \textbf{Contradiction logique.}
\item \textbf{Respect des personnes} : Je traite le prêteur comme un \textbf{simple moyen} pour obtenir de l'argent, en le trompant sur mon intention réelle. Je ne respecte pas sa rationalité ni sa liberté de consentir en connaissance de cause.
\item \textbf{Conclusion} : La maxime est \textbf{immorale}. Il y a un \textbf{devoir parfait} de ne pas faire de fausses promesses.
\end{enumerate}
\end{exoresolu}

\courssub{Autonomie, dignité et le royaume des fins}

Pour Kant, l'être humain n'est pas seulement soumis à la loi morale : il en est \textbf{l'auteur}. C'est l'\cle{autonomie} (du grec \emph{autos} = soi, \emph{nomos} = loi) : \og La volonté est une loi pour elle-même. \fg L'\cle{hétéronomie}, au contraire, c'est obéir à une loi venue de l'extérieur (désirs, coutumes, autorité, recherche du bonheur).

Puisque tout être rationnel se donne la loi morale à lui-même, il possède une \cle{dignité} (valeur infinie, incomparable), non un \textbf{prix} (valeur marchande, échangeable). D'où la deuxième formulation : traiter l'humanité comme une \textbf{fin en soi}.

\begin{aretenir}[Autonomie vs Hétéronomie]
\begin{itemize}
\item \textbf{Autonomie} : Je suis la loi morale \emph{parce que ma raison me la dicte}. Je suis libre (au sens moral) et responsable. C'est la condition de la moralité.
\item \textbf{Hétéronomie} : J'agis sous l'empire de mobiles extérieurs (plaisir, peur, coutume, autorité, intérêt). Même si l'action est conforme au devoir, elle n'est pas \emph{morale} au sens strict.
\end{itemize}
L'idéal kantien : le \textbf{royaume des fins} — une communauté d'êtres rationnels qui se traitent mutuellement comme des fins, se donnant des lois universelles par leur raison commune.
\end{aretenir}

\courssub{Applications camerounaises concrètes}

La morale du devoir n'est pas une abstraction : elle éclaire des situations vécues au Cameroun.

\begin{exemplebox}[Fonctionnaire et pot-de-vin]
Un fonctionnaire au ministère des Finances à Yaoundé reçoit une proposition de \og facilitation \fg (pot-de-vin) pour accélérer un dossier.
\begin{itemize}
\item S'il refuse \emph{par peur de la CONAC} (Commission Nationale Anti-Corruption) : il agit \textbf{conformément au devoir} (mobile : crainte de la sanction). Valeur morale : nulle.
\item S'il refuse \emph{en se disant} : \og La corruption traite les citoyens comme des moyens ; elle viole la loi universelle d'honnêteté ; je ne peux pas vouloir qu'elle devienne une règle générale \fg : il agit \textbf{par devoir}. Valeur morale : authentique.
\end{itemize}
\end{exemplebox}

\begin{exemplebox}[Enseignant et correction des copies]
Un enseignant de Première technique à Douala corrige les copies du Probatoire. Un parent lui propose de l'argent pour majorer la note de son enfant.
\begin{itemize}
\item Refus par peur du conseil de discipline : \textbf{conformité au devoir}.
\item Refus parce que \og Majorer une note par favoritisme, c'est traiter les autres élèves comme des moyens, c'est violer l'universalité de la justice \fg : \textbf{action par devoir}.
\end{itemize}
Seul le second motive une estime de soi morale (le \emph{respect de soi} kantien).
\end{exemplebox}

\courssub{Critiques et limites de la morale kantienne}

Aucune théorie n'est à l'abri d'objections. Les principales critiques adressées à Kant :

\begin{itemize}
\item \textbf{Rigidité / Le mensonge au meurtrier} : Kant affirme qu'on ne doit \textbf{jamais} mentir, même à un meurtrier qui demande où se cache votre ami. Benjamin Constant objecte : \og Dire la vérité au meurtrier, c'en faire le complice. \fg Kant répond que les conséquences ne nous appartiennent pas ; seul le devoir compte. Cette rigidité heurte l'intuition morale commune.
\item \textbf{Conflits de devoirs} : Que faire si deux devoirs parfaits s'opposent ? (Ex. : ne pas mentir vs protéger une vie innocente). Kant nie qu'il y ait de vrais conflits de devoirs (le devoir le plus strict l'emporte), mais la situation concrète laisse souvent perplexe.
\item \textbf{Formalisme} : L'impératif catégorique ne donne pas de \textbf{contenu matériel} précis (quoi faire exactement dans telle situation complexe). Il ne dit que \emph{comment} tester sa maxime. Certains y voient une \og morale vide \fg (Hegel).
\item \textbf{Négligence des inclinations} : Kant ne dit pas qu'il faut \emph{supprimer} les inclinations, mais qu'elles ne doivent pas être le \emph{mobile} de l'action morale. Pourtant, une morale sans place pour la sympathie, l'amour, l'amitié semble inhumaine à certains (Schopenhauer, Nietzsche).
\item \textbf{Autonomie et condition sociale} : L'autonomie suppose une rationalité pleinement exercée. Mais l'aliénation sociale, la pauvreté, l'ignorance peuvent-elles réduire la capacité d'autonomie ? La morale du devoir ne risque-t-elle pas d'ignorer les conditions matérielles de l'agir moral ?
\end{itemize}

\begin{attention}[Ne pas caricaturer Kant]
Kant ne dit \textbf{pas} : \og Il faut être sans cœur, sans émotion, robotique. \fg Il dit : \og L'émotion ne doit pas être le \emph{mobile déterminant} de l'action morale. \fg On peut \emph{éprouver} de la compassion en secourant quelqu'un ; ce qui compte, c'est que la raison dise : \og Je le fais \emph{parce que c'est mon devoir} \fg — et non \og Je le fais \emph{parce que j'en ai envie} \fg.
\end{attention}

\courssub{Synthèse : ce qu'il faut retenir de Kant}

\begin{propriete}[Points clés de la morale du devoir (Kant)]
\begin{enumerate}
\item \textbf{Bonne volonté} : seul bien sans restriction ; source de la valeur morale.
\item \textbf{Devoir} : nécessité d'agir par respect pour la loi morale (mobile pur).
\item \textbf{Conforme au devoir / Par devoir} : distinction cruciale — seule l'intention (maxime) compte.
\item \textbf{Impératif catégorique} : commandement inconditionnel de la raison (vs hypothétique).
\item \textbf{Deux formulations majeures} : (1) Universalité de la maxime ; (2) Humanité comme fin, jamais comme simple moyen.
\item \textbf{Autonomie} : l'être rationnel se donne sa loi — fondement de la \textbf{dignité}.
\item \textbf{Devoirs parfaits} (interdits stricts : ne pas mentir, ne pas tuer, ne pas voler) et \textbf{devoirs imparfaits} (obligations larges : se perfectionner, aider autrui).
\item \textbf{Application} : tester sa maxime par l'universalisation et le respect des personnes.
\end{enumerate}
\end{propriete}

\begin{exercice}[Entraînement : tester une maxime]
Un élève de Première technique, en retard à l'école, envisage de dire au surveillant : \og Mon taxi a eu un accident \fg (alors qu'il a simplement traîné). 
\begin{enumerate}
\item Formuler sa maxime.
\item L'universaliser.
\item Dire si elle passe le test de l'impératif catégorique (contradiction dans la conception ou dans la volonté ? Respect des personnes ?).
\item Conclure : y a-t-il un devoir parfait ou imparfait ici ?
\end{enumerate}
\end{exercice}

\begin{corrige}[Exercice : le mensonge du retardataire]
\begin{enumerate}
\item \textbf{Maxime} : \og Je mentirai sur la cause de mon retard pour éviter une sanction quand je suis en tort. \fg
\item \textbf{Universalisation} : Monde où tout élève ment sur ses retards pour éviter la sanction.
\item \textbf{Test} : 
\begin{itemize}
\item \textit{Contradiction dans la conception} : Si tous mentent, le surveillant ne croira plus aucune excuse. L'excuse (le mensonge) devient inefficace. La maxime se détruit elle-même. \textbf{Contradiction logique.}
\item \textit{Respect des personnes} : L'élève traite le surveillant comme un \textbf{simple moyen} (le tromper pour éviter la sanction), non comme une fin rationnelle.
\end{itemize}
\item \textbf{Conclusion} : Devoir \textbf{parfait} de ne pas mentir. L'action est interdite, quelles que soient les conséquences (sanction, colère des parents).
\end{enumerate}
\end{corrige}

\courssub{Transition vers la synthèse de la leçon}

Nous avons maintenant parcouru les trois grandes familles de morales au programme :
\begin{itemize}
\item Les \textbf{morales de l'intérêt} (Épicurisme, Eudémonisme, Utilitarisme) : la fin détermine le bien (plaisir, bonheur, utilité).
\item La \textbf{morale du devoir} (Kant) : l'intention (bonne volonté, respect de la loi) détermine le bien, indépendamment des conséquences.
\end{itemize}

La dernière section de cette leçon (\textbf{Section 7 : Synthèse et application aux situations concrètes}) mettra en dialogue ces approches à travers des cas pratiques camerounais, et vous entraînera à la rédaction d'une dissertation et d'un commentaire de texte philosophiques — les deux épreuves du Probatoire.

\coursec{Synthèse et application aux situations concrètes}

Nous venons d'étudier la morale du devoir de Kant, qui oppose à toutes les morales de l'intérêt une exigence radicale : agir par devoir, et non par inclination. Il est temps maintenant de rassembler les quatre doctrines étudiées dans cette leçon, de les comparer rigoureusement, puis de les mettre à l'épreuve de situations concrètes de la vie quotidienne. C'est en effet l'un des savoir-faire exigés au Probatoire et au Baccalauréat technique : appliquer les notions de l'unité à des situations de la vie courante et justifier sa réponse.

\courssub{Tableau comparatif des quatre doctrines morales}

Commençons par l'essentiel : chaque doctrine morale répond à trois questions fondamentales. Quelle \cle{fin} l'homme doit-il poursuivre ? Sur quel \cle{fondement} repose la morale ? Quel \cle{critère} permet de dire qu'une action est bonne ? Le tableau suivant résume les réponses.

\begin{center}
\rowcolors{1}{popcream}{white}
\begin{tabularx}{\linewidth}{|l|X|X|X|X|}
\hline
\rowcolor{popblueL} \textbf{Doctrine} & \textbf{Fin poursuivie} & \textbf{Fondement} & \textbf{Critère de l'action bonne} & \textbf{Critique principale} \\
\hline
\textbf{Épicurisme} (Épicure) & Le \cle{plaisir}, entendu comme absence de douleur (\emph{ataraxie}) et de trouble & La nature : tout être vivant recherche le plaisir et fuit la douleur ; la raison calcule les conséquences & L'action est bonne si elle procure un plaisir durable (katastématique) et raisonnable, sans douleur ni trouble ultérieur & Risque d'égoïsme : la justice et l'amitié ne sont valorisées que pour l'utilité qu'elles procurent à l'agent ; certains plaisirs restent exclus \\
\hline
\textbf{Eudémonisme} (Aristote) & Le \cle{bonheur} (\emph{eudaimonia}), souverain bien, vie accomplie selon la vertu & La nature humaine : l'homme est un être raisonnable dont la fonction propre (\emph{ergon}) est l'activité rationnelle vertueuse & L'action est bonne si elle est conforme à la vertu (juste milieu) et contribue à une vie heureuse dans la durée & Le bonheur reste indéterminé (subjectivité) ; la vertu peut sembler élitiste ; risque de justifier l'égoïsme sous couvert d'excellence \\
\hline
\textbf{Utilitarisme} (Bentham, Mill) & La plus grande \cle{utilité} : le plus grand bonheur du plus grand nombre & Le calcul rationnel des conséquences (calcul hédonistique quantitatif chez Bentham, qualitatif chez Mill) & L'action est bonne si ses conséquences maximisent le bonheur général (somme plaisirs moins peines) & Impossibilité pratique de tout calculer ; l'utilité peut justifier le sacrifice injuste d'une minorité (tyrannie de la majorité) \\
\hline
\textbf{Morale du devoir} (Kant) & La bonne \cle{volonté} : agir par devoir, par respect pour la loi morale, non par intérêt & La \cle{raison} pure pratique, autonome, indépendante de toute expérience (a priori) & L'action est bonne si sa maxime est universalisable sans contradiction (impératif catégorique) et si l'homme est traité comme une fin, jamais comme un simple moyen & Rigorisme : abstraction faite des conséquences réelles ; conflits de devoirs possibles ; exigence héroïque difficile à tenir \\
\hline
\end{tabularx}
\end{center}

\begin{aretenir}[Le point commun des morales de l'intérêt]
L'épicurisme, l'eudémonisme et l'utilitarisme sont des \cle{morales de l'intérêt} (ou téléologiques) : chez elles, la \textbf{fin} (plaisir, bonheur, utilité) \emph{précède} la règle et la \emph{justifie}. On obéit à la règle morale \emph{parce que} cela rend heureux (ou procure du plaisir). La morale du devoir (déontologique) inverse ce rapport : chez Kant, c'est la \textbf{règle} (la loi morale) qui précède tout, et la recherche du bonheur n'est légitime que si elle respecte d'abord le devoir. On n'agit pas bien \emph{pour} être heureux ; on agit bien \emph{parce que} c'est notre devoir.
\end{aretenir}

\begin{popfigure}
\begin{tikzpicture}[node distance=0.8cm and 1.2cm,
etape/.style={rectangle, rounded corners, draw=popblue, fill=popblueL, text width=4.8cm, align=center, font=\small, thick},
question/.style={diamond, draw=poporange, fill=poporangeL, aspect=2, align=center, font=\small, inner sep=2pt},
verdict/.style={rectangle, rounded corners, draw=popgreen, fill=popgreenL, text width=3.8cm, align=center, font=\small, thick},
fleche/.style={-{Stealth[length=2.5mm]}, thick, popdark}]
\node[etape, align=center] (e1){\textbf{1. Décrire l'action}\\ Que fait-on exactement ? Dans quelles circonstances ? Avec quelle intention (maxime) ?};
\node[etape, below=of e1, align=center] (e2){\textbf{2. Identifier la fin visée}\\ Plaisir personnel (Épicure) ? Bonheur vertueux (Aristote) ? Utilité générale (Bentham/Mill) ? Respect du devoir (Kant) ?};
\node[etape, below=of e2, align=center] (e3){\textbf{3. Appliquer le critère propre à chaque doctrine}\\ Plaisir durable et sans trouble ? Vertu et juste milieu ? Conduite maximisant l'utilité ? Maxime universalisable + humanité comme fin ?};
\node[question, below=of e3, align=center] (q){Les verdicts\\ convergent-ils ?};
\node[verdict, below left=1cm and 2.5cm of q] (v1) {\textbf{Oui} : le jugement moral est solide, facile à justifier et à communiquer};
\node[verdict, below right=1cm and 2.5cm of q] (v2) {\textbf{Non} : confronter les arguments, peser les valeurs en jeu, puis justifier son jugement personnel};
\draw[fleche] (e1) -- (e2);
\draw[fleche] (e2) -- (e3);
\draw[fleche] (e3) -- (q);
\draw[fleche] (q) -- node[above left, font=\scriptsize, pos=0.5]{oui} (v1);
\draw[fleche] (q) -- node[above right, font=\scriptsize, pos=0.5]{non} (v2);
\end{tikzpicture}
\legende{Arbre de décision : comment juger moralement une action ? On part de la description précise de l'acte et de son intention, on identifie la fin poursuivie, on applique successivement le critère de chaque doctrine, puis on compare les verdicts avant de conclure.}
\end{popfigure}

\courssub{Méthode : résoudre un cas de conscience}

\begin{methode}[Résoudre un cas de conscience en cinq étapes]
\begin{enumerate}
\item \textbf{Décrire la situation} : qui agit, que fait-il, dans quelles circonstances, avec quelle intention (maxime) ?
\item \textbf{Identifier les doctrines mobilisables} : quelles morales peuvent juger cette action ? (Au minimum les quatre du programme).
\item \textbf{Appliquer le critère de chacune} : que dirait Épicure (plaisir durable) ? Aristote (vertu/bonheur) ? Bentham/Mill (utilité générale) ? Kant (universalisation/humanité comme fin) ?
\item \textbf{Comparer les verdicts} : convergent-ils ou s'opposent-ils ? Sur quel point précis (intention vs conséquence, individu vs collectivité) ?
\item \textbf{Justifier son propre jugement} : prendre position avec des arguments philosophiques, sans se contenter d'une opinion non étayée.
\end{enumerate}
\end{methode}

\begin{attention}[Piège fréquent au Bac]
Ne présentez jamais une doctrine comme « vraie » et les autres comme « fausses ». Chaque thèse repose sur des arguments sérieux qu'il faut d'abord \emph{exposer} fidèlement. La discussion (vos critiques, votre position) ne vient qu'ensuite. Un correcteur sanctionne l'élève qui caricature Kant (rigidité inhumaine) ou qui traite Épicure de « défenseur de la débauche » : l'épicurisme est une morale du plaisir \emph{réglé par la raison} (calcul des conséquences), non de l'excès. De même, l'utilitarisme n'est pas le « chacun pour soi », mais le calcul du bien commun.
\end{attention}

\courssub{Étude de cas guidée n°1 : le sujet du devoir trouvé sur un téléphone}

\begin{exoresolu}[Le sujet du devoir surveillé]
\textbf{Énoncé.} Un élève de Première technique trouve, sur le téléphone d'un ami, le sujet du prochain devoir surveillé de philosophie. Il hésite à le recopier et à le partager. Appliquer la méthode du cas de conscience : que diraient Épicure, Aristote, Bentham et Kant ?
\tcblower
\textbf{Solution détaillée.}

\emph{Étape 1 — Description.} L'élève peut obtenir une bonne note sans travail personnel, au détriment de ses camarades qui, eux, préparent honnêtement le devoir. L'intention (maxime) est : « Profiter d'une fuite pour réussir sans effort. »

\emph{Étape 2 — Doctrines mobilisables.} Les quatre doctrines de la leçon peuvent juger cette action.

\emph{Étape 3 — Application des critères.}
\begin{itemize}
\item \textbf{Épicure} : la triche procure un plaisir immédiat (la note, la fierté), mais elle engendre des troubles durables : peur d'être découvert, culpabilité, risque de sanction lourde, perte d'estime de soi. Le plaisir véritable étant l'absence de trouble (\emph{ataraxie}), Épicure conseillerait de refuser : ce plaisir-là coûte plus cher qu'il ne rapporte.
\item \textbf{Aristote} : le bonheur suppose l'exercice des vertus, notamment l'honnêteté, la tempérance et le courage de travailler. Une réussite obtenue par la ruse ne relève pas de l'excellence humaine (\emph{aretè}) ; elle éloigne du bonheur véritable qui demande de l'effort.
\item \textbf{Bentham / Mill} : calculons l'utilité. Le plaisir de l'élève est faible, éphémère et incertain ; la peine collective est grande si la triche se généralise (dévaluation des diplômes, injustice envers les honnêtes élèves, perte de confiance dans l'institution). L'action diminue le bonheur général : elle est condamnable.
\item \textbf{Kant} : la maxime « tricher quand on le peut pour réussir » est-elle universalisable ? Non : si tout le monde triche, l'évaluation perd tout sens, le devoir surveillé ne mesure plus rien, et la maxime se détruit elle-même (contradiction dans la conception). De plus, l'élève utilise ses camarades et son professeur comme de simples moyens pour sa réussite. L'action est contraire au devoir.
\end{itemize}

\emph{Étape 4 — Comparaison.} Les quatre verdicts convergent ici : il ne faut pas tricher, mais pour des raisons différentes (prudence/absence de trouble, vertu/excellence, utilité collective, respect de la loi morale).

\emph{Étape 5 — Jugement personnel.} On peut conclure que la convergence des quatre doctrines rend le jugement particulièrement solide : la triche est moralement condamnable, quelle que soit la fin poursuivie. L'honnêteté s'impose comme une exigence rationnelle partagée.
\end{exoresolu}

\courssub{Étude de cas guidée n°2 : l'infirmier qui détourne des médicaments}

\begin{exemplebox}[Le détournement généreux : utilitarisme contre kantisme]
Un infirmier de l'hôpital de district d'une localité du Cameroun détourne une partie des médicaments du stock pour soigner gratuitement des malades pauvres qui ne peuvent pas payer.

\textbf{Analyse utilitariste.} Les conséquences immédiates sont favorables : des vies sont sauvées, des souffrances soulagées, et le stock détourné est faible au regard des besoins. Au nom du plus grand bonheur du plus grand nombre, l'utilitariste peut \emph{approuver} l'action, du moins tant que le stock réservé aux autres patients n'en souffre pas de manière critique.

\textbf{Analyse kantienne.} Le détournement est un vol, fût-il généreux. La maxime « détourner les biens confiés quand on juge la cause bonne » n'est pas universalisable : si chacun détournait selon son jugement personnel, toute confiance, toute gestion transparente et toute organisation hospitalière s'effondreraient (contradiction dans la volonté). Kant \emph{condamne} l'acte, quelles que soient ses conséquences heureuses, car il viole le devoir de probité et traite l'institution comme un simple moyen.

\textbf{Discussion.} Ce cas montre la force et la limite de chaque doctrine : l'utilitarisme tient compte des conséquences réelles et de la souffrance concrète, mais peut justifier l'illégalité et l'arbitraire ; le kantisme protège la règle, la confiance et l'universalité du droit, mais sa rigueur paraît inhumaine face à l'urgence vitale. Un jugement nuancé cherchera une troisième voie institutionnelle : l'infirmier devrait alerter sa hiérarchie, saisir les ONG ou plaider pour un fonds de solidarité officiel plutôt que de voler, ce qui préserve à la fois le devoir et l'efficacité.
\end{exemplebox}

\courssub{Étude de cas guidée n°3 : le pot-de-vin « parce que tout le monde le fait »}

Un candidat à un concours administratif verse un pot-de-vin en se justifiant ainsi : « tout le monde le fait ». Cette justification résiste-t-elle à l'examen kantien ?

\begin{exemplebox}[L'échec du test d'universalisation]
Soumettons la maxime « verser un pot-de-vin pour réussir » au test kantien : puis-je vouloir que \emph{tous} les candidats agissent ainsi ? Si tout le monde corrompt, le concours ne sélectionne plus les plus compétents mais les plus riches ou les plus corrompus ; il cesse d'être un concours méritocratique pour devenir une vente aux enchères. La maxime, une fois universalisée, \textbf{détruit l'institution même} dont le candidat veut profiter : elle est contradictoire dans sa conception. L'argument « tout le monde le fait » n'est donc pas une justification morale, mais l'aveu que l'on profite d'une faute collective pour servir son intérêt. De plus, le corrupteur et le corrompu se traitent mutuellement comme de simples moyens (l'un pour l'argent, l'autre pour le poste), ce qui viole le second impératif catégorique (humanité comme fin).
\end{exemplebox}

\courssub{Méthode de la dissertation sur un sujet de morale}

\begin{methode}[Disserter sur un sujet de morale]
\begin{enumerate}
\item \textbf{Analyser le sujet} : définir les termes clés (morale, devoir, bonheur, intérêt, conscience, obligation...) et repérer l'opposition ou la question centrale (ex: devoir vs inclination, fin vs moyen, conséquence vs intention).
\item \textbf{Dégager la problématique} : transformer le sujet en un problème philosophique. Exemple : « Faut-il agir par devoir ou par intérêt ? » devient : le bonheur peut-il fonder la morale, ou seul le devoir la rend-il possible ? La morale est-elle hétéronome (venue de l'extérieur/conséquences) ou autonome (venue de la raison) ?
\item \textbf{Construire le plan} : structure classique en trois parties. 
  \begin{itemize}
  \item \textbf{Partie I :} Les morales de l'intérêt (Épicure, Aristote, Bentham/Mill) — thèse : le bien (plaisir, bonheur, utilité) fonde le devoir.
  \item \textbf{Partie II :} La morale du devoir (Kant) — antithèse : le devoir (loi morale) fonde le bien, indépendamment des conséquences.
  \item \textbf{Partie III :} Discussion et dépassement — articulation : le devoir n'exclut pas le bonheur (Kant : « devoir de se rendre heureux » sous condition), mais le subordonne ; la conscience morale comme instance de jugement.
  \end{itemize}
\item \textbf{Argumenter} : chaque partie présente une thèse, ses arguments, un exemple concret (camerounais si possible), puis une critique interne ou externe.
\item \textbf{Conclure avec nuance} : répondre clairement à la problématique, sans caricaturer aucune doctrine, et ouvrir prudemment (ex: sur la conscience morale, le droit, la politique).
\end{enumerate}
\end{methode}

\begin{savaistu}[Le tableau comparatif, votre assurance-points]
Au Baccalauréat technique camerounais, les sujets de philosophie portent régulièrement sur le devoir, le bonheur, la conscience ou les fondements de la morale. L'élève qui maîtrise le tableau comparatif des quatre doctrines (fin, fondement, critère, critique) dispose immédiatement d'un plan structuré, d'auteurs à citer (Épicure, Aristote, Bentham, Mill, Kant) et d'arguments prêts à l'emploi : c'est plusieurs points sécurisés dès l'introduction et le développement.
\end{savaistu}

\courssub{Ouverture : vers le jugement moral personnel}

Faut-il choisir une fois pour toutes entre l'intérêt et le devoir ? Peut-être pas. Les quatre doctrines éclairent chacune un aspect indispensable de la vie morale : le souci légitime du bonheur (Épicure, Aristote), la responsabilité des conséquences collectives (Utilitarisme), le respect inconditionnel de la règle et de la dignité humaine (Kant). Reste alors une faculté que chacun porte en soi et que la philosophie doit affûter.

\begin{definition}[La conscience morale]
La \cle{conscience morale} est la faculté intérieure par laquelle l'homme juge ses propres actions et éprouve, selon le cas, satisfaction (paix) ou remords. Elle est comme un « tribunal intérieur » : nul n'a besoin d'un juge extérieur pour savoir, au fond de lui-même, s'il a bien ou mal agi, pourvu qu'il ait pris le temps de la réflexion rationnelle.
\end{definition}

Ainsi, au-delà des doctrines, la morale aboutit à la formation d'un \cle{jugement moral personnel}, éclairé par la raison, instruit par les grands philosophes, mais assumé en conscience. C'est toute l'ambition de la philosophie au lycée technique : non pas réciter des doctrines, mais apprendre à juger par soi-même avec rigueur et humanité.

\begin{exercice}[Pour s'entraîner]
Un chauffeur de taxi-moto, pressé de gagner sa course, grille un feu rouge la nuit, alors que la route lui semble déserte. Analysez cette action selon les quatre doctrines en suivant la méthode du cas de conscience (5 étapes), puis justifiez votre propre jugement en cinq lignes maximum.
\end{exercice}

\begin{corrige}[Exercice 1]
\textbf{Épicure} : le gain de temps est minime, le risque (accident grave, amende, perte du permis) est disproportionné ; le plaisir durable (\emph{ataraxie}) commande de s'arrêter. \textbf{Aristote} : la prudence (\emph{phronesis}) est la vertu cardinale ; griller le feu est un excès de témérité contraire au juste milieu. \textbf{Bentham/Mill} : si l'exemple se généralise, les accidents nocturnes se multiplient ; l'utilité générale exige le respect strict du code de la route. \textbf{Kant} : la maxime « passer au rouge quand la voie semble libre » n'est pas universalisable, car elle détruit la signification du code de la route (confiance réciproque) ; elle est contraire au devoir. Les quatre verdicts convergent : l'action est condamnable. Le jugement personnel ajoute que la règle protège précisément contre les erreurs d'appréciation subjectives (« la route \emph{semblait} déserte »).
\end{corrige}

\newpage

\coursec{Activité d'intégration}

\coursec{La morale}

\courssub{Objectifs de l'activité}

À l'issue de cette activité d'intégration, l'élève sera capable de :

\begin{itemize}
\item mobiliser les quatre doctrines morales étudiées dans la leçon (épicurisme, eudémonisme, utilitarisme, morale du devoir) pour analyser des situations concrètes de la vie courante ;
\item identifier, dans chaque situation, la \cle{fin poursuivie} par l'agent, le \cle{critère moral} invoqué et le \cle{jugement} qui en découle ;
\item comparer les verdicts produits par des doctrines différentes appliquées à un même cas ;
\item rédiger et justifier une conclusion personnelle argumentée, en mobilisant au moins deux notions du cours (\cle{devoir}, \cle{intérêt}, \cle{bonheur}, \cle{utilité}, \cle{universalisation}) ;
\item travailler en groupe, défendre une position en débat et accepter la confrontation des arguments.
\end{itemize}

\courssub{Situation}

\textbf{Le tribunal des morales.} Le lycée technique de Nkolndongo organise sa semaine de la citoyenneté. Le professeur de philosophie transforme la salle de Première technique en « tribunal des morales ». Trois affaires réelles, tirées de l'actualité et de la vie quotidienne camerounaises, doivent être jugées. Mais attention : il ne s'agit pas de dire si les actes sont légaux ou illégaux. Il s'agit de dire s'ils sont \emph{moralement bons ou mauvais}, et surtout \emph{pourquoi}. La classe est divisée en quatre groupes-experts : les \textbf{épicuriens}, les \textbf{aristotéliciens} (eudémonistes), les \textbf{utilitaristes} et les \textbf{kantiens}. Chaque groupe reçoit les trois dossiers suivants.

\begin{exemplebox}[Dossier n°1 : Le commerçant de Douala]
Un article du quotidien \emph{Mutations} rapporte les faits suivants. M. Etoundi, commerçant au marché Mboppi à Douala, découvre dans sa boutique une sacoche contenant \num{500000} FCFA, oubliée par un client transporteur. Après trois jours de recherche, il retrouve le client et lui restitue intégralement la somme. Interrogé par le journaliste, il déclare : « Rendre cet argent, c'est mon devoir. Mais je vous avoue aussi que depuis, ma réputation a grandi : les clients me font davantage confiance, et mon chiffre d'affaires a augmenté d'environ 20 \% en trois mois. » Un voisin de boutique affirme de son côté : « Il a bien fait, mais s'il avait gardé l'argent, personne ne l'aurait jamais su. »
\end{exemplebox}

\begin{exemplebox}[Dossier n°2 : La déviation routière]
Le ministère des Travaux publics présente un projet de déviation routière dans la région du Nord-Ouest. Les données chiffrées du dossier technique sont les suivantes :
\begin{center}
\begin{tabularx}{\linewidth}{|X|c|}
\hline
\rowcolor{popblueL} \textbf{Élément du projet} & \textbf{Donnée chiffrée} \\
\hline
Familles expropriées (perte des terres et des maisons) & 200 familles, soit environ \num{1200} personnes \\
\hline
Indemnisation prévue par famille & \num{800000} FCFA \\
\hline
Régions désenclavées par la nouvelle route & 3 régions \\
\hline
Population desservie & environ \num{450000} habitants \\
\hline
Baisse estimée du coût du transport des denrées & 30 \% \\
\hline
Emplois créés pendant les travaux & 600 emplois temporaires \\
\hline
\end{tabularx}
\end{center}
Les 200 familles expropriées protestent : « Personne ne nous a demandé notre avis. Nos terres sont notre vie. » Le préfet répond : « L'intérêt général prime : cette route sauvera des vies et nourrira des centaines de milliers de personnes. »
\end{exemplebox}

\begin{exemplebox}[Dossier n°3 : La dénonciation à l'examen]
Lors d'une composition du Probatoire technique blanc, Amina, élève de Première, surprend quatre de ses camarades en train de recopier des réponses sur des téléphones portables dissimulés. Elle hésite toute la nuit, puis dénonce la fraude au chef d'établissement. Les quatre fraudeurs sont exclus de la session. Ses camarades de classe la traitent de « rapporteuse » ; certains refusent de lui parler. Amina déclare : « Je savais que je perdrais des amis. Mais si tout le monde triche, le diplôme ne vaut plus rien pour personne, même pour ceux qui travaillent honnêtement. »
\end{exemplebox}

\courssub{Consignes}

\textbf{Travail en groupes-experts (séance 1 et début de la séance 2).}

\begin{enumerate}
\item Pour chacun des trois dossiers, identifiez précisément : (a) l'action à juger ; (b) la \cle{fin} poursuivie par l'agent d'après votre doctrine ; (c) le \cle{critère} moral que votre doctrine applique ; (d) le \cle{verdict} (moralement bon, moralement mauvais, ou nuancé) que votre doctrine prononce.
\item Rédigez pour chaque dossier un « avis d'expert » de 8 à 10 lignes, justifiant le verdict par au moins une citation ou une thèse de l'auteur de référence de votre groupe (Épicure, Aristote, Bentham ou Mill, Kant).
\item Construisez sur le tableau un tableau comparatif à quatre lignes (une par doctrine) et trois colonnes (une par dossier), en inscrivant dans chaque case le verdict et son critère en quelques mots.
\end{enumerate}

\textbf{Débat plénier (fin de la séance 2).}

\begin{enumerate}
\setcounter{enumi}{3}
\item Chaque groupe présente ses verdicts en 3 minutes. Les autres groupes peuvent poser une question critique par dossier.
\item Relevez au moins deux dossiers où les doctrines aboutissent à des verdicts \emph{opposés}. Expliquez d'où vient le désaccord : sur la fin de la morale ? sur le critère ? sur la place de l'intention ?
\end{enumerate}

\textbf{Conclusion personnelle écrite (séance 3).}

\begin{enumerate}
\setcounter{enumi}{5}
\item Rédigez individuellement une conclusion argumentée de 15 à 20 lignes : pour chaque dossier, indiquez la doctrine que \emph{vous} jugez la plus pertinente, et justifiez votre choix en mobilisant au moins deux notions du cours (\cle{devoir}, \cle{intérêt}, \cle{bonheur}, \cle{utilité}, \cle{universalisation}).
\end{enumerate}

\courssub{Ressources à mobiliser}

\begin{aretenir}[Les quatre doctrines en quatre formules]
\begin{itemize}
\item \textbf{Épicurisme} (Épicure) : le souverain bien est le \cle{plaisir} stable, c'est-à-dire l'absence de trouble (\emph{ataraxie}). Critère : l'action est bonne si elle procure un plaisir durable et écarte la douleur et l'inquiétude.
\item \textbf{Eudémonisme} (Aristote) : le souverain bien est le \cle{bonheur} (\emph{eudaimonia}), qui s'atteint par la pratique des vertus et une vie accomplie. Critère : l'action est bonne si elle réalise l'excellence de l'homme et construit une vie heureuse.
\item \textbf{Utilitarisme} (Bentham, Mill) : l'action est bonne si elle produit « le plus grand bonheur du plus grand nombre ». Critère : l'\cle{utilité}, c'est-à-dire le bilan des conséquences pour l'ensemble des personnes concernées.
\item \textbf{Morale du devoir} (Kant) : l'action est bonne si elle est faite par \cle{devoir} et non par inclination ou intérêt. Critère : l'\cle{impératif catégorique} — « Agis de telle sorte que la maxime de ta volonté puisse valoir comme principe d'une législation universelle » (\cle{universalisation}).
\end{itemize}
\end{aretenir}

\begin{attention}[Pièges fréquents]
Ne confondez pas plaisir immédiat et plaisir épicurien (Épicure refuse les plaisirs qui engendrent ensuite des douleurs). Ne réduisez pas Kant à « il faut obéir aux règles » : chez Kant, c'est l'\emph{intention} et la possibilité d'universaliser la maxime qui fondent la valeur morale, jamais les conséquences. Enfin, l'utilitarisme n'est pas « chacun son intérêt » : c'est l'intérêt du \emph{plus grand nombre}.
\end{attention}

\courssub{Démarche et corrigé commenté}

\begin{exoresolu}[Corrigé commenté du tribunal des morales]
Énoncé : produire les verdicts des quatre doctrines sur les trois dossiers, le tableau comparatif et un modèle de conclusion personnelle.
\tcblower
\textbf{Étape 1 — Dossier n°1 (le commerçant de Douala).}

\emph{Verdict épicurien} : moralement bon. En restituant l'argent, M. Etoundi s'épargne la crainte d'être découvert et le remords ; il jouit d'une conscience tranquille, condition de l'\emph{ataraxie}. Le plaisir stable l'emporte sur le plaisir immédiat de garder \num{500000} FCFA.

\emph{Verdict eudémoniste} : moralement bon. L'honnêteté est une \cle{vertu} ; en la pratiquant, le commerçant accomplit une vie excellente et construit son bonheur. Aristote dirait qu'il agit « comme l'homme vertueux agit ».

\emph{Verdict utilitariste} : moralement bon. Les conséquences profitent à tous : le client récupère son argent, la confiance commerciale augmente au marché Mboppi, le commerçant lui-même gagne 20 \% de chiffre d'affaires. Le bilan global est positif.

\emph{Verdict kantien} : moralement bon \emph{à condition que l'intention soit le devoir}. Kant distingue l'action \emph{conforme au devoir} (par intérêt : la réputation, le gain de 20 \%) et l'action faite \emph{par devoir}. La déclaration « rendre cet argent, c'est mon devoir » plaide pour la valeur morale pleine ; mais si le mobile était le profit, l'action serait seulement légale, sans valeur morale proprement dite. C'est le seul groupe qui soulève le problème de l'\cle{intention}.

\textbf{Étape 2 — Dossier n°2 (la déviation routière).}

\emph{Verdict utilitariste} : moralement bon. Le calcul des conséquences est net : \num{450000} habitants desservis, 30 \% de baisse du coût des denrées, 600 emplois, contre \num{1200} personnes expropriées mais indemnisées (\num{800000} FCFA par famille). « Le plus grand bonheur du plus grand nombre » est du côté du projet.

\emph{Verdict kantien} : moralement condamnable en l'état. Exproprier 200 familles « sans leur demander leur avis », c'est les traiter comme de simples \emph{moyens} au service de l'intérêt général, ce qui viole la deuxième formule de l'impératif catégorique : « Agis de telle sorte que tu traites l'humanité toujours en même temps comme une fin, et jamais simplement comme un moyen. » Le projet ne serait acceptable qu'avec consentement et juste indemnisation.

\emph{Verdict épicurien} : nuancé. Le projet supprime la tranquillité des 200 familles (douleur certaine) pour un bien-être collectif lointain ; l'épicurien exige au minimum que la douleur infligée soit compensée et limitée.

\emph{Verdict eudémoniste} : nuancé. Le bonheur d'une communauté suppose la justice ; sacrifier des familles sans débat corrompt la vie bonne de la cité, même si la prospérité générale s'en trouve améliorée.

\textbf{Étape 3 — Dossier n°3 (la dénonciation).}

\emph{Verdict kantien} : moralement bon. La maxime d'Amina est universalisable : « si tout le monde triche, le diplôme ne vaut plus rien » — la fraude universalisée se détruit elle-même, donc elle est immorale ; la dénonciation, elle, peut valoir comme loi universelle. Amina agit par devoir, au prix de ses amitiés : c'est précisément le signe d'une bonne volonté.

\emph{Verdict utilitariste} : moralement bon. Les conséquences profitent au plus grand nombre : la valeur du diplôme est protégée pour tous les candidats honnêtes ; seuls quatre fraudeurs sont sanctionnés.

\emph{Verdict épicurien} : mitigé. Amina perd la tranquillité (isolement, hostilité) ; mais garder le secret aurait pu nourrir un trouble durable. L'épicurien rappelle que l'amitié est un des plus grands biens : la dénonciation a un coût réel.

\emph{Verdict eudémoniste} : moralement bon. Amina fait preuve de courage et d'honnêteté, vertus qui construisent un caractère excellent, donc un bonheur véritable, supérieur à la popularité perdue.

\textbf{Étape 4 — Tableau comparatif (modèle attendu au tableau).}

\begin{center}
\begin{tabularx}{\linewidth}{|l|X|X|X|}
\hline
\rowcolor{popblueL} \textbf{Doctrine} & \textbf{Dossier 1} & \textbf{Dossier 2} & \textbf{Dossier 3} \\
\hline
Épicurisme & Bon (tranquillité) & Nuancé (douleur des familles) & Mitigé (amitié perdue) \\
\hline
Eudémonisme & Bon (vertu d'honnêteté) & Nuancé (justice de la cité) & Bon (courage, honnêteté) \\
\hline
Utilitarisme & Bon (bénéfices pour tous) & Bon (plus grand nombre) & Bon (diplôme protégé) \\
\hline
Kant & Bon si fait par devoir & Condamnable (personnes traitées comme moyens) & Bon (maxime universalisable) \\
\hline
\end{tabularx}
\end{center}

\textbf{Étape 5 — Modèle de conclusion personnelle (extrait).}

« Sur le dossier du commerçant, je retiens la morale kantienne : seule elle distingue l'honnêteté par devoir de l'honnêteté calculée, car le gain de 20 \% aurait pu être le vrai mobile. Sur le dossier routier, je juge l'utilitarisme plus pertinent pour évaluer le projet dans son ensemble (\num{450000} bénéficiaires contre \num{1200} expropriés), mais la critique kantienne m'oblige à exiger le consentement des familles : l'\cle{utilité} ne justifie pas tout. Sur le dossier d'Amina, la notion d'\cle{universalisation} tranche : la fraude généralisée détruit le diplôme, donc dénoncer était un \cle{devoir}, même si l'épicurisme a raison de rappeler le prix de l'amitié. Aucune doctrine ne suffit seule : le devoir juge les intentions, l'utilité juge les conséquences, et le \cle{bonheur} rappelle la fin dernière de toute morale. »
\end{exoresolu}

\courssub{Bilan de l'activité}

Cette activité a permis de mettre en évidence trois acquis essentiels de la leçon.

\begin{enumerate}
\item \textbf{Une même action peut recevoir des verdicts différents} selon la doctrine mobilisée, parce que chaque morale fixe une \cle{fin} différente (plaisir, bonheur, utilité, devoir) et donc un critère différent du bien et du mal. Le désaccord moral n'est pas un caprice : il est rationnellement explicable.
\item \textbf{Le clivage fondamental de la leçon est apparu concrètement} : les morales de l'\cle{intérêt} (épicurisme, eudémonisme, utilitarisme) jugent par la fin poursuivie et les conséquences, tandis que la morale du \cle{devoir} (Kant) juge par l'intention et l'universalisation de la maxime. Le dossier n°2 a montré la force et les limites de chaque camp : l'utilitarisme légitime l'intérêt général mais risque de sacrifier les minorités ; le kantisme protège la dignité de chacun mais peut paraître rigide.
\item \textbf{Le jugement moral est une démarche argumentée} : identifier la fin, appliquer un critère, justifier un verdict, puis rédiger une conclusion personnelle justifiée — exactement les savoir-faire exigés à la dissertation du Probatoire et du Baccalauréat technique.
\end{enumerate}

\begin{exercice}[Pour aller plus loin]
Sujet de dissertation : « Faut-il toujours agir par devoir, jamais par intérêt ? » Rédigez l'introduction (accroche, définition des termes, problématique, annonce du plan) en vous appuyant sur les trois dossiers du tribunal des morales.
\end{exercice}

\newpage

\coursec{Méthodes et exercices résolus}

\courssub{Savoir-faire 1 : Appliquer les notions de morale à une situation concrète}

\begin{methode}[Analyser une situation morale concrète]
Face à une situation de la vie courante (un cas de tricherie à l'examen, un mensonge, un don, une promesse), la démarche se fait en cinq étapes :
\begin{enumerate}
\item \textbf{Décrire la situation} : identifier l'action posée, l'agent (qui agit ?), les personnes concernées, les conséquences possibles.
\item \textbf{Dégager le problème moral} : formuler la question sous la forme « Cette action est-elle bonne ou mauvaise ? » ou « Que dois-je faire ? ».
\item \textbf{Choisir la grille d'analyse} : demander quel critère de jugement est mobilisé — le \cle{plaisir} (épicurisme), le \cle{bonheur} (eudémonisme), l'\cle{utilité} (utilitarisme), le \cle{devoir} (morale kantienne).
\item \textbf{Appliquer la théorie} : juger l'action selon le critère choisi. Par exemple, selon Kant, on teste l'universalité de la maxime ; selon l'utilitarisme, on compare le bonheur produit et le malheur causé.
\item \textbf{Conclure et nuancer} : donner un jugement clair, puis montrer qu'une autre théorie conduirait à un jugement différent. C'est ce contraste qui fait la valeur de la copie.
\end{enumerate}
\textbf{Réflexe à retenir} : une même action peut être jugée bonne par une morale de l'intérêt et mauvaise par la morale du devoir. Cite toujours au moins deux théories pour montrer la complexité du cas.
\end{methode}

\begin{exoresolu}[La tricherie à l'examen]
\textbf{Énoncé.} À quelques jours du Probatoire technique, un élève de Yaoundé découvre qu'un camarade a photographié le sujet de l'épreuve et le partage sur un groupe WhatsApp. Lui-même hésite à consulter le sujet. Analysez cette situation à l'aide de deux théories morales étudiées, puis concluez.
\tcblower
\textbf{Solution détaillée.}

\textbf{Étape 1 — Description.} L'agent est l'élève ; l'action envisagée est la consultation d'un sujet d'examen volé ; les personnes concernées sont l'élève lui-même, ses camarades (concurrence faussée) et l'institution scolaire.

\textbf{Étape 2 — Problème moral.} La consultation du sujet est-elle une action bonne ou mauvaise ?

\textbf{Étape 3 et 4 — Application de la morale du devoir (Kant).} Pour Kant, une action n'a de valeur morale que si elle est faite par \cle{devoir} et si sa maxime peut être universalisée. Testons la maxime : « Je peux tricher pour réussir ». Universalisée, elle devient : « Tout le monde peut tricher pour réussir ». Mais si tout le monde triche, l'examen ne mesure plus rien, le diplôme perd toute valeur, et l'action détruit elle-même sa condition. La maxime est donc \emph{contradictoire} : la tricherie est moralement condamnable, indépendamment de ses conséquences. De plus, Kant interdit de traiter autrui seulement comme un moyen : tricher, c'est utiliser l'institution et les autres candidats comme de simples instruments de son succès.

\textbf{Application de l'utilitarisme.} L'utilitarisme (Bentham, Stuart Mill) juge une action selon ses \cle{conséquences} : est bonne l'action qui produit le plus grand bonheur pour le plus grand nombre. La tricherie procure un plaisir immédiat à un seul individu, mais elle produit du malheur pour beaucoup : injustice envers les candidats honnêtes, dévalorisation du diplôme pour tous, risque de sanction collective. La balance des conséquences est négative : l'utilitarisme condamne aussi la tricherie.

\textbf{Étape 5 — Conclusion.} Les deux théories convergent ici : consulter le sujet volé est une action mauvaise. Mais elles ne la condamnent pas pour la même raison : Kant invoque la contradiction de la maxime et le respect du devoir, l'utilitarisme invoque le calcul des conséquences. Cela montre que des morales différentes peuvent aboutir au même jugement par des chemins différents.
\end{exoresolu}

\courssub{Savoir-faire 2 : Distinguer les trois morales de l'intérêt}

\begin{methode}[Identifier la morale de l'intérêt mobilisée dans un texte ou un cas]
Les morales de l'intérêt fondent la valeur de l'action sur un avantage pour l'agent. Pour les distinguer sans confusion :
\begin{enumerate}
\item \textbf{Repérer le mot-clé du texte} : « plaisir » renvoie à l'\cle{épicurisme} (Épicure) ; « bonheur » renvoie à l'\cle{eudémonisme} (Aristote) ; « utilité », « plus grand bonheur du plus grand nombre » renvoient à l'\cle{utilitarisme} (Bentham, Stuart Mill).
\item \textbf{Identifier le critère} : le plaisir est un critère \emph{immédiat et individuel} ; le bonheur est un critère de \emph{vie entière accomplie} ; l'utilité est un critère \emph{collectif} (le bonheur de tous).
\item \textbf{Vérifier la nuance épicurienne} : Épicure ne prône pas la débauche mais le plaisir \emph{stable} (l'absence de douleur du corps et de trouble de l'âme, l'\emph{ataraxie}) et le choix raisonnable des plaisirs.
\item \textbf{Conclure} : nommer la théorie, citer son auteur et justifier en une phrase pourquoi le critère correspond.
\end{enumerate}
\textbf{Formule à retenir} : « Toutes les morales de l'intérêt visent le bonheur, mais elles diffèrent sur ce qui le constitue : le plaisir, l'accomplissement de soi ou l'utilité générale. »
\end{methode}

\begin{exoresolu}[Trois attitudes, trois morales]
\textbf{Énoncé.} Trois élèves justifient leur conduite : (a) « Je sors avec mes amis parce que cela me fait plaisir, mais je refuse les excès qui me rendraient malade » ; (b) « Je travaille dur à l'atelier parce qu'une vie réussie est une vie où l'on accomplit ses talents » ; (c) « Je fais du bénévolat au quartier parce que cela rend le plus de gens heureux ». Identifiez la théorie morale qui correspond à chaque attitude, en justifiant.
\tcblower
\textbf{Solution détaillée.}

\textbf{Attitude (a) — L'épicurisme.} Le critère invoqué est le \cle{plaisir}. Mais l'élève choisit ses plaisirs et refuse les excès : c'est exactement la position d'\textbf{Épicure}, qui distingue les désirs naturels et nécessaires des désirs vains, et qui définit le souverain bien comme l'absence de douleur et la tranquillité de l'âme. Le plaisir épicurien est un plaisir mesuré, choisi par la raison.

\textbf{Attitude (b) — L'eudémonisme.} Le critère est le \cle{bonheur} compris comme accomplissement de soi sur une vie entière. C'est la position d'\textbf{Aristote} : le bonheur (\emph{eudaimonia}) est la fin suprême, et il s'obtient par l'exercice des vertus et le développement de ses capacités propres. Travailler à l'atelier pour accomplir ses talents relève de cette morale de la réalisation de soi.

\textbf{Attitude (c) — L'utilitarisme.} Le critère est l'\cle{utilité}, c'est-à-dire le plus grand bonheur du plus grand nombre. C'est la formule de \textbf{Bentham}, reprise par \textbf{Stuart Mill} : une action est bonne si elle maximise le bonheur collectif. Le bénévolat est justifié non par le plaisir personnel ni par l'accomplissement de soi, mais par le bien-être produit pour la communauté.

\textbf{Conclusion.} Les trois attitudes relèvent toutes des morales de l'intérêt, car chacune fonde la valeur de l'action sur un avantage (plaisir, bonheur, utilité). Mais elles se distinguent par la nature du bien visé : plaisir individuel mesuré, accomplissement d'une vie, ou bonheur collectif.
\end{exoresolu}

\courssub{Savoir-faire 3 : Expliquer et appliquer la morale du devoir de Kant}

\begin{methode}[Appliquer le test kantien de l'universalité]
Pour juger une action selon Kant, il faut suivre une procédure rigoureuse :
\begin{enumerate}
\item \textbf{Formuler la maxime} de l'action : la règle subjective que suit l'agent (« Je mens pour me sortir d'embarras »).
\item \textbf{Universaliser la maxime} : imaginer que cette règle devient une loi pour tous (« Tout le monde ment pour se sortir d'embarras »).
\item \textbf{Chercher la contradiction} : si la maxime universalisée se détruit elle-même (si tous mentent, la parole ne veut plus rien dire, donc mentir devient impossible), l'action est \emph{immorale}.
\item \textbf{Vérifier le respect de la personne} : l'action traite-t-elle autrui comme une \cle{fin} en soi, ou seulement comme un moyen ? Utiliser autrui comme simple instrument est immoral.
\item \textbf{Conclure} : l'action n'est morale que si elle est faite par \cle{devoir}, c'est-à-dire par respect de la loi morale, et non par intérêt ou inclination.
\end{enumerate}
\textbf{Réflexe} : chez Kant, ce n'est pas le résultat qui compte, mais l'\emph{intention}. Une action utile faite par intérêt n'a pas de valeur morale proprement dite.
\end{methode}

\begin{exoresolu}[Le mensonge au commerçant]
\textbf{Énoncé.} Un commerçant de Douala, pour vendre davantage, affirme à ses clients que sa marchandise est de qualité supérieure alors qu'il sait qu'elle est défectueuse. Jugez cette conduite à la lumière de la morale kantienne.
\tcblower
\textbf{Solution détaillée.}

\textbf{Étape 1 — Formulation de la maxime.} La maxime du commerçant est : « Je mens aux clients chaque fois que le mensonge augmente mon profit. »

\textbf{Étape 2 — Universalisation.} Imaginons que cette maxime devienne une loi universelle : « Tout commerçant ment chaque fois que le mensonge est profitable. »

\textbf{Étape 3 — Recherche de la contradiction.} Si tous les commerçants mentent systématiquement, alors aucun client ne croit plus aucun commerçant : la confiance, condition de tout échange commercial, disparaît. Le mensonge universalisé détruit donc la possibilité même du mensonge (on ne peut tromper que quelqu'un qui nous fait confiance). La maxime est \emph{contradictoire} : elle ne peut pas être érigée en loi universelle.

\textbf{Étape 4 — Respect de la personne.} En mentant, le commerçant traite ses clients comme de simples \emph{moyens} de son enrichissement, et non comme des personnes raisonnables capables de décider librement en connaissance de cause. Or la deuxième formule de l'impératif catégorique exige d'agir de telle sorte que l'humanité, en soi et en autrui, soit toujours traitée en même temps comme une \cle{fin}.

\textbf{Étape 5 — Conclusion.} La conduite du commerçant est immorale au sens de Kant, et ce pour deux raisons : sa maxime ne supporte pas l'universalisation, et elle instrumentalise autrui. On remarquera que ce jugement ne dépend pas des conséquences : même si le mensonge enrichissait le commerçant sans jamais être découvert, il resterait moralement condamnable, car la valeur morale réside dans l'intention et non dans le résultat.
\end{exoresolu}

\courssub{Savoir-faire 4 : Rédiger une introduction de dissertation sur la morale}

\begin{methode}[Construire une introduction en quatre temps]
Une introduction de dissertation au Probatoire ou au Baccalauréat technique comporte quatre moments obligatoires :
\begin{enumerate}
\item \textbf{L'accroche (amorce)} : partir d'un constat large et concret (une situation de la vie courante, un fait d'expérience) qui mène au sujet.
\item \textbf{La présentation du sujet} : définir brièvement les termes clés du sujet (ici : la \cle{morale}, ensemble des règles de conduite qui distinguent le bien du mal).
\item \textbf{La problématique} : formuler le problème sous forme d'une question (ou de deux questions liées) qui montre une tension. Exemple : la morale doit-elle viser notre intérêt ou nous commander des devoirs ?
\item \textbf{L'annonce du plan} : présenter en une phrase les deux ou trois parties du développement, sans les numéroter.
\end{enumerate}
\textbf{Réflexes} : l'introduction tient en 8 à 12 lignes ; on n'y donne jamais son opinion personnelle ; on évite les définitions de dictionnaire (« Selon le dictionnaire... »).
\end{methode}

\begin{exoresolu}[Sujet : « Faut-il agir par intérêt ou par devoir ? »]
\textbf{Énoncé.} Rédigez l'introduction complète d'une dissertation sur le sujet : « Faut-il agir par intérêt ou par devoir ? »
\tcblower
\textbf{Solution détaillée (modèle rédigé).}

\textbf{Accroche.} Dans la vie quotidienne, chacun de nos actes peut être expliqué de deux manières : le commerçant honnête l'est-il parce que l'honnêteté attire la clientèle, ou parce qu'il estime devoir être honnête ? L'élève qui travaille le fait-il pour réussir, ou parce qu'il juge que travailler est son devoir ?

\textbf{Présentation du sujet et définitions.} Cette alternative renvoie à la question du \cle{fondement de la morale}, c'est-à-dire de ce qui donne sa valeur à une action. Agir par \cle{intérêt}, c'est rechercher un avantage : le plaisir, le bonheur ou l'utilité. Agir par \cle{devoir}, c'est obéir à une obligation morale indépendamment de tout avantage attendu.

\textbf{Problématique.} Le problème est alors le suivant : une action n'a-t-elle de valeur morale que si elle sert notre intérêt et contribue à notre bonheur, ou bien la morale exige-t-elle au contraire que nous accomplissions notre devoir pour lui-même, même contre notre intérêt ? Autrement dit, l'intérêt et le devoir sont-ils conciliables ?

\textbf{Annonce du plan.} Nous verrons d'abord que les morales de l'intérêt fondent la valeur de l'action sur le plaisir, le bonheur ou l'utilité ; nous montrerons ensuite que la morale du devoir, avec Kant, exige d'agir par pur respect de la loi morale ; nous pourrons enfin nous demander si ces deux exigences ne se rejoignent pas dans la vie concrète.

\textbf{Commentaire méthodologique.} Cette introduction respecte les quatre étapes : l'accroche part d'exemples camerounais concrets ; les termes « intérêt » et « devoir » sont définis ; la problématique est formulée en questions ; le plan est annoncé en trois parties sans numérotation.
\end{exoresolu}

\courssub{Savoir-faire 5 : Rédiger un paragraphe argumenté et une conclusion}

\begin{methode}[La règle du paragraphe unique]
Chaque paragraphe de développement suit la structure \textbf{I-D-E-A} :
\begin{enumerate}
\item \textbf{Idée} : une phrase qui énonce clairement la thèse du paragraphe.
\item \textbf{Définition} : préciser les notions mobilisées (ex. : l'\cle{impératif catégorique} est un commandement inconditionné : « Agis de telle sorte que la maxime de ta volonté puisse valoir en même temps comme principe d'une législation universelle »).
\item \textbf{Exemple} : illustrer par un cas concret, de préférence tiré de la vie courante camerounaise.
\item \textbf{Argumentation} : expliquer pourquoi l'exemple prouve l'idée, en citant brièvement l'auteur (Épicure, Aristote, Bentham, Stuart Mill, Kant).
\end{enumerate}
Pour la \textbf{conclusion} : (1) répondre directement à la problématique ; (2) résumer le chemin parcouru en deux phrases ; (3) éventuellement ouvrir sur une question nouvelle, sans introduire de notion hors programme.
\end{methode}

\begin{exoresolu}[Rédiger un paragraphe sur l'utilitarisme]
\textbf{Énoncé.} Dans le cadre du sujet « Le bonheur des autres fait-il mon bonheur ? », rédigez un paragraphe de développement présentant l'utilitarisme.
\tcblower
\textbf{Solution détaillée (modèle rédigé).}

\textbf{(Idée)} L'utilitarisme soutient que la morale consiste à rechercher le bonheur du plus grand nombre, et non seulement le sien propre. \textbf{(Définition)} Selon Bentham, puis Stuart Mill, une action est bonne si elle est \cle{utile}, c'est-à-dire si elle produit le maximum de bonheur et le minimum de souffrance pour l'ensemble des personnes concernées : c'est le principe du « plus grand bonheur du plus grand nombre ». \textbf{(Exemple)} Ainsi, lorsqu'un chef de quartier organise la vaccination des enfants ou le nettoyage d'un marché, son action est moralement bonne non parce qu'elle lui rapporte un profit personnel, mais parce qu'elle augmente le bien-être de toute la communauté. \textbf{(Argumentation)} L'intérêt individuel n'est donc pas rejeté, mais il est élargi : mon propre bonheur est compté, mais il ne pèse pas plus lourd que celui d'autrui. Comme le souligne Stuart Mill, l'utilitarisme exige une impartialité stricte entre mon bonheur et celui des autres. Le bonheur des autres devient ainsi une partie constitutive de mon propre bonheur, puisque je ne puis être durablement heureux dans une communauté malheureuse.

\textbf{Commentaire méthodologique.} Le paragraphe respecte la structure I-D-E-A : une idée unique (l'utilitarisme fonde la morale sur le bonheur collectif), une définition précise avec les auteurs, un exemple camerounais concret, et une argumentation qui relie l'exemple à la thèse et répond au sujet posé.
\end{exoresolu}

\courssub{Savoir-faire 6 : Commenter une citation d'auteur}

\begin{methode}[Expliquer une citation en trois étapes]
Au commentaire de texte ou lorsqu'une citation est donnée à l'examen :
\begin{enumerate}
\item \textbf{Situer} : nommer l'auteur et la théorie à laquelle la citation appartient (ex. : Kant, morale du devoir).
\item \textbf{Expliquer} : reformuler la citation avec ses propres mots, en définissant chaque terme technique (maxime, impératif catégorique, ataraxie, utilité...). C'est l'étape la plus longue.
\item \textbf{Discuter} : montrer la portée de la thèse (ce qu'elle apporte) puis sa limite éventuelle (ce qu'une théorie rivale lui objecterait), en restant dans les notions du programme.
\end{enumerate}
\textbf{Réflexe} : ne jamais paraphraser mot à mot ; une explication doit éclairer ce que la citation suppose sans le dire.
\end{methode}

\begin{exoresolu}[Commenter une formule de Kant]
\textbf{Énoncé.} Expliquez la formule suivante de Kant : « Agis de telle sorte que la maxime de ta volonté puisse valoir en même temps comme principe d'une législation universelle. »
\tcblower
\textbf{Solution détaillée.}

\textbf{Étape 1 — Situation.} Cette formule est la première formulation de l'\cle{impératif catégorique}, c'est-à-dire du commandement suprême de la morale du devoir développée par Emmanuel Kant.

\textbf{Étape 2 — Explication.} Une \emph{maxime} est la règle subjective selon laquelle j'agis concrètement (« je rends l'argent trouvé », « je mens pour réussir »). Kant demande de soumettre cette maxime à un test : puis-je vouloir, sans contradiction, que ma règle personnelle devienne une loi valable pour tous les hommes ? Si oui, l'action est morale ; si la maxime universalisée se contredit elle-même, l'action est immorale. Par exemple, la maxime « je fais une promesse que je n'ai pas l'intention de tenir » ne peut être universalisée : si tout le monde faisait de fausses promesses, personne ne croirait plus aux promesses, et la promesse elle-même deviendrait impossible. L'impératif est dit \emph{catégorique} parce qu'il commande sans condition : il ne dit pas « si tu veux réussir, sois honnête » (ce serait un impératif hypothétique, lié à l'intérêt), mais « sois honnête », quoi qu'il en coûte.

\textbf{Étape 3 — Discussion.} La force de cette formule est de donner un critère universel et rationnel du bien moral, indépendant des conséquences et des intérêts particuliers : elle garantit le respect égal de toutes les personnes. Sa limite, objectée par les morales de l'intérêt, est sa rigueur : elle ignore les conséquences concrètes et peut sembler exigeante au point d'opposer le devoir au bonheur. Mais c'est précisément cette exigence qui, selon Kant, distingue la moralité authentique du simple calcul d'intérêt.

\textbf{Conclusion.} La formule kantienne signifie donc que la moralité d'une action se mesure à l'universalité possible de sa règle : n'est bon que ce que tous pourraient faire sans contradiction.
\end{exoresolu}

\begin{attention}[Les 5 erreurs qui coûtent des points au Bac]
\begin{enumerate}
\item \textbf{Confondre épicurisme et recherche des plaisirs excessifs.} Épicure prône un plaisir \emph{mesuré} : l'absence de douleur et la tranquillité de l'âme. Écrire « Épicure invite à jouir sans limites » est une contresens éliminatoire.
\item \textbf{Confondre les trois morales de l'intérêt.} Plaisir (épicurisme), bonheur-accomplissement (eudémonisme d'Aristote) et utilité collective (utilitarisme) ne sont pas synonymes : attribuer la formule « le plus grand bonheur du plus grand nombre » à Épicure est une erreur grave.
\item \textbf{Présenter Kant comme un moraliste du résultat.} Chez Kant, la valeur morale tient à l'\emph{intention} et au devoir, jamais aux conséquences. Dire « selon Kant, mentir est mal parce que cela a de mauvaises conséquences » renverse sa pensée.
\item \textbf{Oublier le test de l'universalisation.} Beaucoup de copies citent l'impératif catégorique sans montrer la contradiction de la maxime : c'est pourtant ce raisonnement (formuler la maxime, l'universaliser, dégager la contradiction) qui est attendu et noté.
\item \textbf{Négliger la méthode de la dissertation.} Introduction sans problématique claire, plan annoncé avec des numéros (« Premièrement... Deuxièmement... »), paragraphes sans exemple concret, conclusion qui ne répond pas à la question posée : autant de fautes de forme sanctionnées, même quand le contenu est bon.
\end{enumerate}
\end{attention}

\newpage

\coursec{Exercices}

\courssub{Je vérifie mes connaissances}

\begin{exercice}[QCM : repérer les doctrines]
Pour chaque question, une seule réponse est correcte. Entoure la lettre correspondante.
\begin{enumerate}
\item Selon Épicure, le souverain bien est :
\begin{enumerate}
\item[a)] le devoir \item[b)] le plaisir \item[c)] l'utilité sociale \item[d)] la richesse
\end{enumerate}
\item Pour Kant, une action a une valeur morale lorsqu'elle est accomplie :
\begin{enumerate}
\item[a)] par intérêt personnel \item[b)] par inclination \item[c)] par devoir \item[d)] par habitude
\end{enumerate}
\item L'utilitarisme évalue une action selon :
\begin{enumerate}
\item[a)] son intention \item[b)] ses conséquences pour le plus grand nombre \item[c)] la tradition \item[d)] la loi divine
\end{enumerate}
\item L'eudémonisme identifie la fin de la morale à :
\begin{enumerate}
\item[a)] le bonheur \item[b)] le plaisir immédiat \item[c)] l'obéissance \item[d)] la gloire
\end{enumerate}
\end{enumerate}
\end{exercice}

\begin{exercice}[Vrai ou faux, en justifiant]
Réponds par vrai ou faux, puis justifie ta réponse en deux ou trois lignes.
\begin{enumerate}
\item La morale se confond avec le droit, car toutes deux reposent sur des sanctions extérieures.
\item Selon Épicure, tous les plaisirs doivent être recherchés sans distinction.
\item Chez Kant, l'impératif catégorique commande inconditionnellement, sans condition ni fin extérieure.
\item L'utilitarisme peut justifier le sacrifice d'un individu si cela procure le bonheur du plus grand nombre.
\end{enumerate}
\end{exercice}

\begin{exercice}[Définitions à compléter]
Recopie et complète les définitions suivantes avec les mots de la liste : \emph{devoir, plaisir, bonheur, utilité, obligation, intérêt}.
\begin{enumerate}
\item La morale est l'ensemble des règles de conduite qui s'imposent à la conscience comme une \dots\dots
\item L'épicurisme fait du \dots\dots le souverain bien et le critère des actions.
\item L'eudémonisme fait du \dots\dots la fin dernière de l'action humaine.
\item L'utilitarisme fait de l'\dots\dots le critère de la valeur morale d'une action.
\item La morale kantienne est une morale du \dots\dots, fondée sur la raison et non sur l'\dots\dots
\end{enumerate}
\end{exercice}

\begin{exercice}[Association auteur-doctrine]
Relie chaque auteur ou courant à la formule qui le caractérise, puis cite un exemple de la vie courante illustrant chaque formule.
\begin{center}
\begin{tabularx}{\linewidth}{|l|X|}
\hline
\rowcolor{popblueL} \textbf{Doctrine} & \textbf{Formule} \\
\hline
1. Épicurisme & a) « Agis de telle sorte que ta maxime puisse être érigée en loi universelle » \\
\hline
2. Eudémonisme & b) « Le plus grand bonheur du plus grand nombre » \\
\hline
3. Utilitarisme & c) « Il faut choisir les plaisirs durables et renoncer aux désirs vains » \\
\hline
4. Morale du devoir & d) « Tout homme agit naturellement en vue du bonheur » \\
\hline
\end{tabularx}
\end{center}
\end{exercice}

\courssub{Je m'entraîne}

\begin{exercice}[Situations concrètes camerounaises]
Voici cinq situations de la vie courante. Pour chacune, indique la doctrine morale (épicurisme, eudémonisme, utilitarisme ou morale du devoir) qui la \emph{justifierait} et celle qui la \emph{condamnerait}, puis rédige la justification en trois à cinq lignes.
\begin{enumerate}
\item Un fonctionnaire de Yaoundé accepte un pot-de-vin pour accélérer la délivrance d'un acte de naissance, arguant que « tout le monde le fait et que sa famille en profite ».
\item Dans un village du Nord-Ouest, les habitants organisent une collecte (\emph{njangi}) pour payer les frais d'hôpital d'un voisin malade.
\item Un élève de Première recopie sur son voisin à l'examen du Probatoire « pour ne pas redoubler et faire honte à ses parents ».
\item Un homme d'affaires de Douala finance la construction d'un forage d'eau potable dans son village natal.
\item Un conducteur de taxi-moto respecte le feu rouge à minuit, alors qu'aucun policier ni aucun véhicule n'est en vue.
\end{enumerate}
\end{exercice}

\begin{exercice}[Analyse de données : la vaccination obligatoire]
Une politique publique fictive rend un vaccin obligatoire. Les données recueillies après un an sont les suivantes :
\begin{center}
\begin{tabularx}{\linewidth}{|X|c|c|}
\hline
\rowcolor{popblueL} \textbf{Indicateur} & \textbf{Avant} & \textbf{Après} \\
\hline
Nombre de cas de la maladie & 50\,000 & 5\,000 \\
\hline
Nombre de décès & 2\,000 & 150 \\
\hline
Effets secondaires graves déclarés & --- & 80 \\
\hline
Personnes sanctionnées pour refus de vaccination & --- & 1\,200 \\
\hline
\end{tabularx}
\end{center}
\begin{enumerate}
\item Calcule le taux de diminution du nombre de cas et du nombre de décès (en \%).
\item Rédige en une vingtaine de lignes le jugement qu'un \emph{utilitariste} porterait sur cette politique : pèse-t-elle positivement pour le plus grand nombre ?
\item Rédige en une vingtaine de lignes le jugement qu'un \emph{kantien} porterait sur cette même politique : l'obligation respecte-t-elle la dignité et la liberté des personnes ?
\end{enumerate}
\end{exercice}

\begin{exercice}[Classification des désirs selon Épicure]
Épicure distingue les désirs naturels et nécessaires, les désirs naturels et non nécessaires, et les désirs vains.
\begin{enumerate}
\item Rappelle la définition de chacune de ces trois catégories.
\item Classe les désirs suivants dans la bonne catégorie en justifiant : manger du ndolé quand on a faim ; posséder le dernier téléphone à la mode ; boire de l'eau ; devenir célèbre sur les réseaux sociaux ; dormir ; manger des mets raffinés chaque jour.
\item Déduis-en la conduite que l'épicurien adopte face à chaque catégorie de désirs.
\end{enumerate}
\end{exercice}

\begin{exercice}[Intérêt ou devoir ?]
Deux amis, Aminatou et Basile, rendent visite à un camarade hospitalisé à l'hôpital de district. Aminatou déclare : « Je viens parce que cela me fait plaisir et parce qu'il fera de même pour moi un jour. » Basile déclare : « Je viens parce que c'est mon devoir d'ami, même si cela me coûte. »
\begin{enumerate}
\item Identifie la conception morale qui inspire chacun des deux amis.
\item Selon Kant, laquelle des deux conduites possède une véritable valeur morale ? Justifie.
\item Peut-on dire pour autant que la conduite d'Aminatou est immorale ? Discute en une dizaine de lignes.
\end{enumerate}
\end{exercice}

\courssub{Je me prépare au Bac}

\begin{exercice}[Dissertation : devoir ou intérêt ?]
Sujet : \emph{« Faut-il agir par devoir ou par intérêt ? »}
\begin{enumerate}
\item Dégage le problème posé par ce sujet et formule une problématique en une ou deux questions.
\item Rédige l'introduction complète de la dissertation (accroche, définition des termes, problématique, annonce du plan).
\item Construis le plan détaillé d'un développement en trois parties : première thèse (agir par intérêt : épicurisme, eudémonisme, utilitarisme), deuxième thèse (agir par devoir : Kant), troisième partie (confrontation et position personnelle argumentée).
\item Rédige la conclusion de la dissertation.
\end{enumerate}
\end{exercice}

\begin{exercice}[Commentaire de texte : Épicure, \emph{Lettre à Ménécée}]
Texte : « Il faut méditer sur ce qui procure le bonheur, puisque, le bonheur présent, nous possédons tout, et, le bonheur absent, nous faisons tout pour le posséder. [\dots] Nous considérons le plaisir comme le commencement et la fin de la vie heureuse. [\dots] Mais quand nous disons que le plaisir est la fin, nous ne visons pas les plaisirs des débauchés [\dots] : c'est au contraire ne pas souffrir du corps et ne pas être troublé dans l'âme. [\dots] Il faut donc rapporter tout désir à cette fin : certains désirs sont naturels, d'autres sont vains ; et parmi les désirs naturels, les uns sont nécessaires, les autres simplement naturels. »
\begin{enumerate}
\item Situe le texte : auteur, œuvre, doctrine défendue.
\item Explique la thèse centrale du texte : en quel sens le plaisir est-il « le commencement et la fin de la vie heureuse » ?
\item Pourquoi Épicure refuse-t-il d'identifier le plaisir aux « plaisirs des débauchés » ? Quelle conception du plaisir défend-il ?
\item Reconstitue la classification des désirs présentée à la fin du texte et explique son rôle dans la recherche du bonheur.
\item Discussion : la recherche du plaisir ainsi comprise est-elle une morale égoïste ? Peut-elle fonder la vie en société au Cameroun aujourd'hui ? Justifie ta réponse en une quinzaine de lignes.
\end{enumerate}
\end{exercice}

\begin{exercice}[Question de synthèse type Probatoire]
\begin{enumerate}
\item Construis un tableau comparatif des quatre doctrines morales étudiées (épicurisme, eudémonisme, utilitarisme, morale du devoir) selon les rubriques suivantes : fin poursuivie, fondement de la morale, critère de l'action bonne, auteur(s) de référence, critique possible.
\item À partir de ce tableau, rédige un paragraphe de synthèse (une quinzaine de lignes) répondant à la question : qu'est-ce qui oppose fondamentalement les morales de l'intérêt à la morale du devoir ?
\item Application : un jeune entrepreneur de Bafoussam hésite entre délocaliser son atelier pour augmenter ses profits et le maintenir pour préserver l'emploi de ses dix ouvriers. Analyse ce dilemme successivement du point de vue utilitariste et du point de vue kantien, puis donne ton avis personnel argumenté.
\end{enumerate}
\end{exercice}

\newpage

\coursec{Corrigés des exercices}

\begin{corrige}[Exercice 1 — QCM : repérer les doctrines]
\begin{enumerate}
\item \textbf{Réponse b) : le plaisir.} Pour Épicure, le plaisir est « le commencement et la fin de la vie heureuse » : c'est le souverain bien, c'est-à-dire la fin dernière à laquelle toutes nos actions doivent être rapportées.
\item \textbf{Réponse c) : par devoir.} Chez Kant, une action n'a de valeur morale authentique que si elle est accomplie \emph{par devoir}, et non simplement \emph{conformément au devoir} par intérêt ou par inclination.
\item \textbf{Réponse b) : ses conséquences pour le plus grand nombre.} L'utilitarisme (Bentham, Mill) juge une action selon son \emph{utilité}, c'est-à-dire sa capacité à produire « le plus grand bonheur du plus grand nombre ».
\item \textbf{Réponse a) : le bonheur.} L'eudémonisme (du grec \emph{eudaimonia}, bonheur) fait du bonheur la fin dernière de l'action humaine : tout homme agit naturellement en vue d'être heureux.
\end{enumerate}
\textbf{Points de vigilance :} ne pas confondre \emph{plaisir} (Épicure) et \emph{bonheur} (eudémonisme) ; ne pas attribuer à Kant un critère fondé sur les conséquences, ce qui est la marque de l'utilitarisme.
\end{corrige}

\begin{corrige}[Exercice 2 — Vrai ou faux, en justifiant]
\begin{enumerate}
\item \textbf{Faux.} La morale s'impose à la conscience de l'intérieur et sa sanction est intérieure (remords, approbation de soi) ; le droit repose au contraire sur des règles extérieures assorties de sanctions extérieures (tribunaux, amendes, prison). On peut d'ailleurs obéir au droit sans être moral, et accomplir un devoir moral que la loi n'impose pas.
\item \textbf{Faux.} Épicure distingue les désirs : il faut satisfaire les désirs naturels et nécessaires, modérer les désirs naturels non nécessaires et renoncer aux désirs vains. Certains plaisirs doivent être écartés parce qu'ils engendrent des douleurs plus grandes (excès, troubles de l'âme).
\item \textbf{Vrai.} L'impératif catégorique commande \emph{inconditionnellement} : « agis de telle sorte que la maxime de ton action puisse être érigée en loi universelle ». Il ne dépend d'aucune fin extérieure, contrairement à l'impératif hypothétique (« si tu veux ceci, fais cela »).
\item \textbf{Vrai, mais avec une réserve.} Le calcul utilitariste peut en effet conduire à sacrifier un individu si le bonheur global du plus grand nombre s'en trouve augmenté : c'est la logique des conséquences. C'est précisément la critique que Kant adresse à l'utilitarisme : nul ne doit être traité seulement comme un moyen.
\end{enumerate}
\textbf{Points de vigilance :} la justification doit citer la distinction fondamentale (sanction intérieure/extérieure ; impératif catégorique/hypothétique ; classification des désirs) et non seulement affirmer.
\end{corrige}

\begin{corrige}[Exercice 3 — Définitions à compléter]
\begin{enumerate}
\item La morale est l'ensemble des règles de conduite qui s'imposent à la conscience comme une \textbf{obligation}.
\item L'épicurisme fait du \textbf{plaisir} le souverain bien et le critère des actions.
\item L'eudémonisme fait du \textbf{bonheur} la fin dernière de l'action humaine.
\item L'utilitarisme fait de l'\textbf{utilité} le critère de la valeur morale d'une action.
\item La morale kantienne est une morale du \textbf{devoir}, fondée sur la raison et non sur l'\textbf{intérêt}.
\end{enumerate}
\textbf{Points de vigilance :} le mot \emph{devoir} ne peut convenir à la phrase 1 (on demande ce que la morale fait éprouver à la conscience : une obligation) ; dans la phrase 5, le second terme attendu est \emph{intérêt}, car Kant oppose l'action par devoir à l'action par intérêt ou par inclination.
\end{corrige}

\begin{corrige}[Exercice 4 — Association auteur-doctrine]
\textbf{Associations :} 1-c ; 2-d ; 3-b ; 4-a.
\begin{itemize}
\item \textbf{1-c, Épicurisme :} « Il faut choisir les plaisirs durables et renoncer aux désirs vains ». \emph{Exemple :} un élève renonce à une soirée la veille d'une composition pour préserver sa tranquillité et sa réussite.
\item \textbf{2-d, Eudémonisme :} « Tout homme agit naturellement en vue du bonheur ». \emph{Exemple :} un jeune de Bafoussam choisit une filière de formation parce qu'elle lui promet une vie épanouie et stable.
\item \textbf{3-b, Utilitarisme :} « Le plus grand bonheur du plus grand nombre ». \emph{Exemple :} une mairie de Douala installe l'éclairage public dans les quartiers populaires pour la sécurité du plus grand nombre.
\item \textbf{4-a, Morale du devoir :} « Agis de telle sorte que ta maxime puisse être érigée en loi universelle ». \emph{Exemple :} un commerçant rend la monnaie exacte même quand le client ne compte pas, parce que l'honnêteté doit valoir pour tous.
\end{itemize}
\textbf{Points de vigilance :} l'exemple doit illustrer la \emph{formule} et non seulement la doctrine en général ; pour Kant, l'exemple doit montrer une action faite par devoir, sans recherche d'avantage.
\end{corrige}

\begin{corrige}[Exercice 5 — Situations concrètes camerounaises]
\begin{enumerate}
\item \textbf{Le pot-de-vin.} \emph{Justifierait :} un intérêt égoïste déguisé (mais aucune des quatre doctrines ne le justifie vraiment ; à la rigueur un épicurisme mal compris). \emph{Condamnerait :} la morale du devoir (la corruption ne peut être érigée en loi universelle : si tous corrompaient, le service public s'effondrerait) et l'utilitarisme (la corruption nuit au plus grand nombre en bloquant les services et en creusant les inégalités). L'argument « tout le monde le fait » n'est pas une justification morale : une maxime ne devient pas bonne parce qu'elle est répandue.
\item \textbf{La collecte (\emph{njangi}).} \emph{Justifierait :} l'utilitarisme (le geste soulage un malade et renforce le bien-être du groupe : le plus grand bonheur du plus grand nombre) et la morale du devoir (aider son prochain est un devoir universalisable). \emph{Condamnerait :} un épicurisme étroit et égoïste, qui renoncerait à la dépense ; mais l'épicurisme authentique valorise l'amitié, donc ne condamne pas réellement.
\item \textbf{La fraude au Probatoire.} \emph{Justifierait :} un calcul d'intérêt personnel (éviter le redoublement), mais aucune doctrine étudiée ne l'approuve. \emph{Condamnerait :} la morale du devoir (la fraude universalisée détruit la valeur du diplôme et la confiance dans l'examen) et l'utilitarisme (la fraude nuit à la qualité de la formation de tous). L'épicurien ajouterait que la fraude entretient l'inquiétude d'être découvert, donc le trouble de l'âme.
\item \textbf{Le forage d'eau potable.} \emph{Justifierait :} l'utilitarisme (l'accès à l'eau profite durablement au plus grand nombre) ; la morale du devoir si le donateur agit par pur devoir de solidarité. \emph{Condamnerait :} aucune doctrine ne condamne ce geste ; seul un égoïsme radical s'y refuserait.
\item \textbf{Le feu rouge respecté à minuit.} \emph{Justifierait :} la morale du devoir (le respect de la règle vaut même sans témoin ni sanction, car la maxime « respecter le code de la route » doit valoir universellement). \emph{Condamnerait :} aucune doctrine ne condamne ; un utilitariste calculant à court terme pourrait juger l'arrêt inutile, mais l'utilitarisme réglé approuve le respect général des feux, source de sécurité pour tous.
\end{enumerate}
\textbf{Points de vigilance :} bien distinguer l'action \emph{conforme} au devoir (par intérêt, par crainte du gendarme) de l'action \emph{par devoir} ; dans la situation 5, l'absence de témoin est précisément ce qui révèle la valeur morale de la conduite.
\end{corrige}

\begin{corrige}[Exercice 6 — Analyse de données : la vaccination obligatoire]
\textbf{1. Calculs des taux de diminution.}
\[ \text{Cas} : \frac{50\,000 - 5\,000}{50\,000}\times 100 = \frac{45\,000}{50\,000}\times 100 = 90\ \% \]
\[ \text{Décès} : \frac{2\,000 - 150}{2\,000}\times 100 = \frac{1\,850}{2\,000}\times 100 = 92{,}5\ \% \]
Le nombre de cas a diminué de \textbf{90 \%} et le nombre de décès de \textbf{92,5 \%}.

\textbf{2. Jugement utilitariste.} L'utilitariste évalue la politique à ses conséquences pour le plus grand nombre. Or le bilan est massivement positif : 45\,000 malades et 1\,850 décès évités, contre 80 effets secondaires graves. Le rapport est écrasant : pour un effet secondaire grave, on compte environ 23 vies sauvées ($1\,850/80 \approx 23$). Le « plus grand bonheur du plus grand nombre » est donc largement servi : la santé publique s'améliore, les familles sont épargnées, les hôpitaux désengorgés, l'économie préservée. Certes, les 1\,200 personnes sanctionnées subissent un coût, et les 80 victimes d'effets secondaires paient un prix réel ; l'utilitariste regrette ces dommages mais les juge inférieurs au bien collectif obtenu. Il recommandera toutefois d'améliorer la surveillance des effets secondaires et de préférer la persuasion à la sanction, car une politique acceptée produit encore plus d'utilité qu'une politique subie. Conclusion utilitariste : la politique est moralement justifiée, car elle maximise le bonheur global.

\textbf{3. Jugement kantien.} Le kantien ne juge pas d'abord aux conséquences mais au principe de l'action et au respect des personnes. Il demande : l'obligation vaccinale traite-t-elle les individus comme des fins, ou seulement comme des moyens au service de la statistique sanitaire ? Or la seconde formule de l'impératif catégorique interdit de traiter l'humanité, en soi ou en autrui, uniquement comme un moyen. Contraindre un corps sans le consentement de la personne, sous peine de sanction, semble violer son autonomie, c'est-à-dire sa capacité de se donner à soi-même sa loi. Les 1\,200 sanctions témoignent de cette contrainte. Cependant, le kantien peut aussi défendre l'obligation : refuser le vaccin met autrui en danger, et la maxime « je me soustrais à la protection collective » ne peut être universalisée sans détruire la santé de tous. Le kantien exigera donc que la loi soit universelle, égale pour tous, débattue rationnellement, et accompagnée d'information plutôt que de pure répression. Conclusion kantienne : la fin sanitaire est légitime, mais le moyen coercitif n'est admissible que s'il respecte la dignité et la liberté de chacun.

\textbf{Barème indicatif :} calculs exacts (4 pts) ; argument utilitariste chiffré et nuancé (8 pts) ; argument kantien citant l'impératif catégorique et la dignité (8 pts).
\textbf{Points de vigilance :} le taux se calcule par rapport à la valeur \emph{initiale} ; le kantien ne doit pas être caricaturé en adversaire de la santé publique : il juge le \emph{principe}, pas seulement le résultat.
\end{corrige}

\begin{corrige}[Exercice 7 — Classification des désirs selon Épicure]
\textbf{1. Définitions.}
\begin{itemize}
\item \textbf{Désirs naturels et nécessaires :} désirs dont la satisfaction est indispensable à la vie et à la tranquillité (manger, boire, dormir, s'abriter).
\item \textbf{Désirs naturels et non nécessaires :} désirs qui prolongent la nature sans être indispensables (mets raffinés, confort recherché) ; ils ne suppriment pas la faim, ils la varient.
\item \textbf{Désirs vains :} désirs sans fondement naturel et sans limite, nés des opinions creuses (richesse, gloire, célébrité, pouvoir).
\end{itemize}
\textbf{2. Classement.}
\begin{itemize}
\item \emph{Manger du ndolé quand on a faim :} naturel et nécessaire (il apaise la faim, besoin vital).
\item \emph{Boire de l'eau :} naturel et nécessaire (la soif doit être étanchée pour vivre).
\item \emph{Dormir :} naturel et nécessaire (le repos est vital).
\item \emph{Manger des mets raffinés chaque jour :} naturel mais non nécessaire (le raffinement varie le plaisir sans être indispensable).
\item \emph{Posséder le dernier téléphone à la mode :} vain (désir né de l'opinion et de la comparaison sociale, sans limite).
\item \emph{Devenir célèbre sur les réseaux sociaux :} vain (la gloire est insatiable et dépend d'autrui).
\end{itemize}
\textbf{3. Conduite de l'épicurien.} Il satisfait simplement les désirs naturels et nécessaires ; il goûte avec modération, sans en dépendre, les désirs naturels non nécessaires ; il renonce aux désirs vains, sources de trouble et d'insatisfaction. Ainsi atteint-il l'\emph{ataraxie}, la tranquillité de l'âme, et l'absence de douleur du corps : le vrai plaisir.
\textbf{Points de vigilance :} le critère de classement porte sur le \emph{désir} et sa limite naturelle, non sur l'objet lui-même ; le téléphone n'est pas condamné en soi, mais le désir illimité de la nouveauté.
\end{corrige}

\begin{corrige}[Exercice 8 — Intérêt ou devoir ?]
\textbf{1. Doctrines inspirantes.} Aminatou agit par inclination et par calcul d'intérêt réciproque : sa conduite relève des morales de l'intérêt (eudémonisme, voire épicurisme de l'amitié). Basile agit par devoir, indépendamment de son plaisir : sa conduite relève de la morale kantienne.

\textbf{2. Valeur morale selon Kant.} Selon Kant, seule la conduite de Basile possède une véritable valeur morale, car elle est accomplie \emph{par devoir} : la maxime « venir en aide à un ami malade » peut être érigée en loi universelle, et elle est suivie même quand elle coûte. Aminatou agit certes \emph{conformément} au devoir, mais par inclination (plaisir) et par intérêt (réciprocité) : son mobile est pathologique au sens kantien, c'est-à-dire sensible, et non rationnel.

\textbf{3. Discussion.} On ne peut pourtant pas dire que la conduite d'Aminatou est \emph{immorale} : elle n'est pas contraire au devoir, elle est seulement dépourvue de valeur morale au sens strict. Kant lui-même reconnaît que les actions conformes au devoir sont utiles et louables socialement : un ami visité par plaisir est tout de même réconforté. De plus, l'amitié vécue dans la joie a une valeur humaine évidente, que l'épicurien et l'eudémoniste reconnaissent. Mais la valeur \emph{morale} proprement dite, c'est-à-dire la pureté de l'intention, manque : si le plaisir disparaît ou si la réciprocité fait défaut, la visite cessera. La morale du devoir garantit la constance et l'universalité de la conduite. On dira donc que la conduite d'Aminatou est moralement neutre ou simplement légale au plan moral, tandis que celle de Basile est morale au sens plein.
\textbf{Points de vigilance :} bien employer la distinction kantienne \emph{agir par devoir} / \emph{agir conformément au devoir} ; ne pas conclure trop vite que l'intérêt est « immoral » : il est \emph{amoral} ou non moral, ce qui est différent.
\end{corrige}

\begin{corrige}[Exercice 9 — Dissertation : « Faut-il agir par devoir ou par intérêt ? »]
\textbf{1. Problème et problématique.} Le sujet oppose deux mobiles de l'action : l'intérêt (recherche d'un avantage, d'un plaisir, d'un bonheur) et le devoir (obéissance à une obligation rationnelle universelle). Le problème : la moralité réside-t-elle dans les conséquences utiles de l'acte ou dans la pureté de son intention ? Problématique : \emph{agir moralement, est-ce rechercher son intérêt et le bonheur du plus grand nombre, ou accomplir son devoir pour lui-même ? Les deux mobiles peuvent-ils se concilier ?}

\textbf{2. Introduction (modèle rédigé).} Au marché de Mokolo, un commerçant honnête rend la monnaie : le fait-il parce que l'honnêteté fidélise sa clientèle, ou parce que l'honnêteté est un devoir ? Cette question quotidienne engage toute la réflexion morale. Le \emph{devoir} désigne l'obligation qui s'impose à la conscience indépendamment de tout avantage ; l'\emph{intérêt} désigne ce qui est utile ou agréable à l'agent ou au plus grand nombre. La morale semble exiger le désintéressement, pourtant toute action humaine vise une fin qui satisfait. Dès lors, faut-il agir par devoir ou par intérêt ? L'intérêt peut-il fonder la moralité, ou la moralité exige-t-elle le pur devoir ? Nous verrons d'abord que les morales de l'intérêt font du plaisir, du bonheur ou de l'utilité le critère de l'action ; puis que Kant fonde la moralité sur le seul devoir ; enfin nous confronterons ces positions pour dégager une réponse personnelle.

\textbf{3. Plan détaillé.}
\begin{itemize}
\item \textbf{I. Agir par intérêt : les morales eudémonistes et utilitaires.} a) Tout homme recherche naturellement le bonheur (eudémonisme) : l'intérêt est un fait de nature. b) L'épicurisme : le plaisir bien compris (ataraxie, classification des désirs) est le souverain bien. c) L'utilitarisme : l'intérêt peut être universel — le plus grand bonheur du plus grand nombre ; exemple de la \emph{njangi} ou d'une politique de santé publique.
\item \textbf{II. Agir par devoir : la morale kantienne.} a) Seule la bonne volonté est bonne sans restriction ; agir par devoir, non par inclination. b) L'impératif catégorique : universaliser sa maxime ; traiter la personne comme fin. c) Critique des morales de l'intérêt : elles rendent la morale instable, relative aux circonstances, et peuvent sacrifier l'individu.
\item \textbf{III. Confrontation et position personnelle.} a) Limites du pur intérêt (la corruption « rentable » reste immorale) et limites du devoir rigide (un devoir sans humanité peut devenir inhumain). b) Position argumentée : le devoir est le fondement de la moralité, mais le bonheur reste la fin légitime de l'homme ; l'idéal est une conduite où le devoir accompli concourt au bonheur de tous — l'honnêteté du commerçant sert à la fois sa conscience et sa clientèle.
\end{itemize}

\textbf{4. Conclusion (modèle rédigé).} Agir par intérêt est naturel, et l'utilitarisme montre que l'intérêt peut s'élargir au bonheur du plus grand nombre ; mais l'intérêt, même généralisé, ne garantit ni la justice envers chacun ni la constance de la volonté. Agir par devoir, au contraire, donne à l'action sa valeur morale propre : une intention universalisable et respectueuse des personnes. Nous répondrons donc qu'il faut agir par devoir, sans renoncer au bonheur : la moralité n'interdit pas d'être heureux, elle exige seulement que le bonheur ne soit jamais acheté au prix du devoir. L'homme accompli est celui qui, comme le conducteur respectueux du feu rouge à minuit, fait son devoir même lorsque rien ne l'y oblige extérieurement.

\textbf{Barème indicatif :} problématique (2 pts) ; introduction complète (4 pts) ; plan cohérent en trois parties avec auteurs et exemples (8 pts) ; conclusion répondant à la problématique (4 pts) ; langue et présentation (2 pts).
\textbf{Points de vigilance :} définir les deux termes du sujet dès l'introduction ; citer au moins un auteur par partie ; la troisième partie doit trancher avec des arguments, non par un simple compromis.
\end{corrige}

\begin{corrige}[Exercice 10 — Commentaire de texte : Épicure, \emph{Lettre à Ménécée}]
\textbf{1. Situation du texte.} Le texte est extrait de la \emph{Lettre à Ménécée} d'Épicure (philosophe grec, III\textsuperscript{e} siècle av. J.-C.). Il expose l'épicurisme, morale du plaisir : le plaisir y est présenté comme le souverain bien, à condition d'être correctement compris.

\textbf{2. Thèse centrale.} Le plaisir est « le commencement et la fin de la vie heureuse » : c'est de lui que part toute action (nous fuyons la douleur et cherchons le plaisir dès la naissance) et c'est à lui qu'elle aboutit (le bonheur est l'état de plaisir stable). Le plaisir est donc à la fois le point de départ naturel et le but ultime de l'existence : le bonheur présent, nous possédons tout.

\textbf{3. Refus des « plaisirs des débauchés ».} Épicure refuse l'assimilation de sa doctrine à la débauche : les plaisirs violents et désordonnés engendrent des douleurs plus grandes qu'eux (maladies, remords, dépendances). Le vrai plaisir épicurien est négatif et stable : c'est l'absence de douleur du corps (\emph{aponie}) et l'absence de trouble de l'âme (\emph{ataraxie}). Le plaisir n'est pas l'ivresse mais la santé, non l'agitation mais la tranquillité.

\textbf{4. Classification des désirs.} Épicure distingue : les désirs \emph{naturels et nécessaires} (manger, boire), qu'il faut satisfaire ; les désirs \emph{naturels et non nécessaires} (mets raffinés), qu'il faut modérer ; les désirs \emph{vains} (richesse, gloire), qu'il faut éliminer. Cette classification est l'instrument pratique du bonheur : en rapportant tout désir à la fin du plaisir stable, on évite les troubles inutiles et l'on se contente de peu, ce qui rend le bonheur facile à atteindre.

\textbf{5. Discussion.} Cette morale est-elle égoïste ? Elle l'est en apparence, car chacun y cherche son propre plaisir. Mais l'égoïsme épicurien est éclairé : Épicure célèbre l'amitié comme le plus sûr des biens, car sans amis la vie est insécurisée et triste ; la justice elle-même est définie comme un accord utile pour ne pas se nuire mutuellement. La recherche du plaisir bien compris conduit donc à la modération, à la justice et à l'amitié, vertus indispensables à la vie sociale. Au Cameroun aujourd'hui, une telle sagesse aurait sa place : renoncer aux désirs vains (course aux biens de consommation, prestige des réseaux sociaux) libérerait de l'endettement et de la corruption ; valoriser les plaisirs simples (famille, amitié, travail accompli) renforcerait le vivre-ensemble. Cependant, cette morale reste centrée sur le bonheur individuel : elle peine à fonder le sacrifice de soi pour autrui, que la morale du devoir, elle, peut exiger. Elle est donc une sagesse de vie précieuse, mais insuffisante comme fondement complet de la morale sociale.

\textbf{Barème indicatif :} situation du texte (2 pts) ; explication de la thèse (4 pts) ; conception du plaisir (4 pts) ; classification des désirs (4 pts) ; discussion argumentée et personnelle (6 pts).
\textbf{Points de vigilance :} ne pas caricaturer Épicure en partisan de la débauche — c'est l'erreur la plus fréquente ; employer les notions d'\emph{ataraxie} et de \emph{souverain bien} ; la discussion doit mobiliser un exemple camerounais.
\end{corrige}

\begin{corrige}[Exercice 11 — Question de synthèse type Probatoire]
\textbf{1. Tableau comparatif.}
\begin{center}
\begin{tabularx}{\linewidth}{|l|X|X|X|X|}
\hline
\rowcolor{popblueL} \textbf{Rubrique} & \textbf{Épicurisme} & \textbf{Eudémonisme} & \textbf{Utilitarisme} & \textbf{Morale du devoir} \\
\hline
Fin poursuivie & Le plaisir (absence de douleur, ataraxie) & Le bonheur & Le plus grand bonheur du plus grand nombre & L'accomplissement du devoir pour lui-même \\
\hline
Fondement & La nature (le plaisir est le premier bien) & La tendance naturelle au bonheur & Le calcul des conséquences & La raison pure pratique \\
\hline
Critère de l'action bonne & Ce qui procure un plaisir durable sans douleur & Ce qui conduit au bonheur & L'utilité, le bilan des conséquences & La conformité de la maxime à la loi universelle \\
\hline
Auteur(s) & Épicure & Aristote & Bentham, Mill & Kant \\
\hline
Critique possible & Risque d'égoïsme et de retrait du monde & Le bonheur est vague et relatif à chacun & Peut sacrifier l'individu au plus grand nombre & Rigueur sans égard aux conséquences ni aux sentiments \\
\hline
\end{tabularx}
\end{center}

\textbf{2. Paragraphe de synthèse.} Ce qui oppose fondamentalement les morales de l'intérêt à la morale du devoir, c'est la place accordée aux \emph{conséquences} et au \emph{mobile} de l'action. Les morales de l'intérêt — épicurisme, eudémonisme, utilitarisme — jugent l'action d'après la fin qu'elle atteint : plaisir, bonheur ou utilité collective. Elles sont des morales \emph{conséquentialistes} et \emph{naturelles} : elles partent du désir humain d'être heureux et font de la satisfaction le critère du bien. La morale du devoir, au contraire, est une morale de l'\emph{intention} : pour Kant, une action n'est bonne que si sa maxime est universalisable et si elle est accomplie par pur respect du devoir, indépendamment de tout résultat agréable ou utile. Là où l'intérêt demande « que vais-je gagner, que gagnera le plus grand nombre ? », le devoir demande « puis-je vouloir que tous agissent ainsi ? ». La première perspective rend la morale relative aux circonstances et aux calculs ; la seconde lui donne un fondement absolu et inconditionné dans la raison. L'opposition est donc celle de la \emph{fin} et du \emph{principe}, du \emph{bonheur recherché} et de l'\emph{obligation respectée}.

\textbf{3. Application : le dilemme de l'entrepreneur de Bafoussam.}
\emph{Point de vue utilitariste :} il faut comparer les conséquences globales. Délocaliser augmente les profits, donc le bonheur de l'entrepreneur, de sa famille, peut-être de nouveaux employés ailleurs et des consommateurs (prix plus bas) ; mais elle jette dix ouvriers et leurs familles dans la précarité, appauvrit le quartier. Si la délocalisation crée plus d'emplois et de richesse qu'elle n'en détruit, l'utilitariste la justifie ; sinon il la condamne. Le jugement dépend du calcul, donc des circonstances.
\emph{Point de vue kantien :} la question n'est pas le bilan mais le principe. Licencier dix ouvriers pour le seul profit, c'est les traiter uniquement comme des moyens (des coûts à supprimer) et non comme des fins en soi : l'impératif catégorique l'interdit. La maxime « je sacrifie autrui à mon profit » n'est pas universalisable sans détruire la confiance qui fonde toute entreprise. Le kantien exige donc le maintien de l'atelier, ou au minimum que les ouvriers soient respectés ( reclassement, concertation).
\emph{Avis personnel :} la position kantienne me paraît devoir primer : dix familles ne sauraient être sacrifiées à un calcul de profit, car la dignité des personnes n'a pas de prix mais une valeur absolue. Toutefois, l'utilitarisme rappelle utilement qu'une entreprise qui fait faillite par excès de fidélité ne protège plus personne : l'idéal est une solution où le devoir envers les ouvriers inspire une reconversion ou une modernisation de l'atelier, conciliant dignité et viabilité économique.

\textbf{Barème indicatif :} tableau complet et exact (6 pts) ; synthèse dégageant l'opposition conséquences/intention (6 pts) ; application aux deux doctrines et avis argumenté (8 pts).
\textbf{Points de vigilance :} le tableau doit rester fidèle aux auteurs du programme ; dans l'application, formuler explicitement l'impératif catégorique pour le point de vue kantien et le calcul du plus grand bonheur pour l'utilitarisme ; l'avis personnel doit être \emph{argumenté}, pas seulement affirmé.
\end{corrige}

\newpage

\coursec{Fiche bilan}

\begin{popfigure}
\begin{tikzpicture}[mindmap, grow cyclic, every node/.style={concept}, text width=2.6cm, align=center,
root concept/.append style={concept color=popblue, font=\large\bfseries, text=white, minimum size=3.2cm},
level 1/.append style={level distance=4.4cm, sibling angle=51.5, font=\small\bfseries},
level 2/.append style={level distance=2.9cm, font=\footnotesize, text width=2.2cm}]
\node[root concept]{LA MORALE}
 child[concept color=poporange]{ node {Définition et fondement}
   child[concept color=poporange!60]{ node {Ensemble de règles de conduite} }
   child[concept color=poporange!60]{ node {Fondement : raison, nature, Dieu, société} }
 }
 child[concept color=popgreen]{ node {Morales de l'intérêt}
   child[concept color=popgreen!60]{ node {But : l'intérêt, le bonheur de l'agent} }
 }
 child[concept color=poppurple]{ node {Épicurisme}
   child[concept color=poppurple!60]{ node {Plaisir stable : ataraxie} }
   child[concept color=poppurple!60]{ node {Désirs naturels et nécessaires} }
 }
 child[concept color=poppink]{ node {Eudémonisme}
   child[concept color=poppink!60]{ node {Bonheur = fin suprême} }
   child[concept color=poppink!60]{ node {Aristote : vie vertueuse} }
 }
 child[concept color=popgold]{ node {Utilitarisme}
   child[concept color=popgold!60]{ node {Bentham, Mill} }
   child[concept color=popgold!60]{ node {Plus grand bonheur du plus grand nombre} }
 }
 child[concept color=popblue]{ node {Morale du devoir}
   child[concept color=popblue!60]{ node {Kant : impératif catégorique} }
   child[concept color=popblue!60]{ node {Agir par devoir, non par inclination} }
 };
\end{tikzpicture}
\legende{Carte mentale de la leçon : la morale, ses fondements, les trois morales de l'intérêt et la morale kantienne du devoir.}
\end{popfigure}

\begin{bilan}
\textbf{Définitions clés}
\begin{itemize}
\item \cle{Morale} : ensemble des règles et des valeurs qui distinguent le \cle{bien} du \cle{mal} et orientent la conduite humaine vers le bien.
\item \cle{Fondement de la morale} : ce sur quoi repose l'obligation morale (la raison, la nature humaine, Dieu, la société, l'intérêt).
\item \cle{Morale de l'intérêt} : morale qui fonde l'action sur la recherche de l'intérêt, du plaisir ou du bonheur de l'agent.
\item \cle{Morale du devoir} : morale qui commande d'agir par respect de la loi morale, indépendamment de tout intérêt.
\end{itemize}

\textbf{Les doctrines à connaître par c\oe ur}
\begin{center}
\begin{tabularx}{\linewidth}{|l|X|X|}
\hline
\rowcolor{popblueL} \textbf{Doctrine} & \textbf{Auteur(s)} & \textbf{Idée maîtresse} \\
\hline
Épicurisme & Épicure & Le plaisir est le souverain bien ; viser l'\emph{ataraxie} (absence de trouble) en satisfaisant les seuls désirs naturels et nécessaires. \\
\hline
Eudémonisme & Aristote & Le bonheur (\emph{eudaimonia}) est la fin suprême ; il s'atteint par une vie vertueuse et rationnelle. \\
\hline
Utilitarisme & Bentham, Mill & Est bonne l'action utile qui procure le plus grand bonheur pour le plus grand nombre. \\
\hline
Morale du devoir & Kant & Agir par \cle{devoir}, selon l'\cle{impératif catégorique} : « Agis de telle sorte que la maxime de ton action puisse être érigée en loi universelle. » \\
\hline
\end{tabularx}
\end{center}

\textbf{Oppositions fondamentales}
\begin{itemize}
\item Impératif \cle{hypothétique} (« si tu veux X, fais Y ») $\neq$ impératif \cle{catégorique} (« fais ton devoir, sans condition »).
\item Agir \emph{par intérêt} (conséquences, plaisir) $\neq$ agir \emph{par devoir} (respect de la loi morale).
\item Plaisir \emph{en mouvement} $\neq$ plaisir \emph{stable} (ataraxie) chez Épicure.
\end{itemize}

\textbf{Méthodes express}
\begin{itemize}
\item \emph{Dissertation} : définir la morale dans l'introduction, opposer morales de l'intérêt et morale du devoir, conclure par une position argumentée.
\item \emph{Analyse de situation} : identifier la fin poursuivie (plaisir ? bonheur ? utilité ? devoir ?) pour reconnaître la doctrine en jeu.
\item Toujours illustrer par un exemple concret de la vie courante (honnêteté au marché, corruption, entraide au village).
\end{itemize}
\end{bilan}

\begin{aretenir}[Checklist avant le Bac]
\begin{itemize}
\item[$\square$] Je sais définir la morale et distinguer bien et mal.
\item[$\square$] Je connais au moins trois fondements possibles de la morale (raison, nature, Dieu, société).
\item[$\square$] Je sais ce qui caractérise une morale de l'intérêt.
\item[$\square$] Je connais la doctrine d'Épicure : plaisir, ataraxie, classification des désirs.
\item[$\square$] Je sais que l'eudémonisme fait du bonheur la fin suprême (Aristote).
\item[$\square$] Je connais le principe de l'utilitarisme : le plus grand bonheur du plus grand nombre (Bentham, Mill).
\item[$\square$] Je sais citer et expliquer l'impératif catégorique de Kant.
\item[$\square$] Je distingue impératif hypothétique et impératif catégorique.
\item[$\square$] Je sais opposer « agir par devoir » et « agir par intérêt » avec des exemples.
\item[$\square$] Je sais appliquer ces doctrines à une situation concrète de la vie courante camerounaise.
\item[$\square$] Je sais construire une dissertation : problématique, thèses opposées, exemples, conclusion personnelle.
\item[$\square$] Je cite les auteurs brièvement et correctement (Épicure, Aristote, Bentham, Mill, Kant).
\end{itemize}
\end{aretenir}

\findecours

\end{document}
